\documentclass[11pt,oneside]{article}
\usepackage[T1]{fontenc}
\usepackage[utf8]{inputenc}
\usepackage{lmodern}
\usepackage[a4paper,margin=27mm,headheight=15pt]{geometry}
\usepackage{microtype}
\usepackage{amsmath,amssymb,amsthm,mathtools,aliascnt}
\usepackage{booktabs,longtable,array,enumitem}
\usepackage[dvipsnames]{xcolor}
\usepackage{fancyhdr}
\usepackage{listings}
\usepackage{xurl}
\usepackage{hyperref}
\hypersetup{colorlinks=true,linkcolor=MidnightBlue,citecolor=MidnightBlue,
 urlcolor=MidnightBlue,pdftitle={Notations for Omnific Integers},
 pdfauthor={Merged research report for the Surreal project},
 pdfsubject={Holonomic towers, shielded and rational-omega calculi, and the complexity of equality for omnific notations}}
\usepackage[nameinlink,noabbrev]{cleveref}
\pagestyle{fancy}
\fancyhf{}
\fancyhead[L]{\small Notations for omnific integers}
\fancyhead[R]{\small\thepage}
\renewcommand{\headrulewidth}{0.3pt}
\setlength{\parskip}{4pt}
\setlength{\parindent}{1.2em}
\setlength{\emergencystretch}{2.5em}
\setcounter{tocdepth}{2}
\numberwithin{equation}{section}
\newtheorem{theorem}{Theorem}[section]
\crefname{theorem}{theorem}{theorems}
\newaliascnt{proposition}{theorem}
\newtheorem{proposition}[proposition]{Proposition}
\aliascntresetthe{proposition}
\crefname{proposition}{proposition}{propositions}
\newaliascnt{lemma}{theorem}
\newtheorem{lemma}[lemma]{Lemma}
\aliascntresetthe{lemma}
\crefname{lemma}{lemma}{lemmas}
\newaliascnt{corollary}{theorem}
\newtheorem{corollary}[corollary]{Corollary}
\aliascntresetthe{corollary}
\crefname{corollary}{corollary}{corollaries}
\theoremstyle{definition}
\newaliascnt{definition}{theorem}
\newtheorem{definition}[definition]{Definition}
\aliascntresetthe{definition}
\crefname{definition}{definition}{definitions}
\newaliascnt{example}{theorem}
\newtheorem{example}[example]{Example}
\aliascntresetthe{example}
\crefname{example}{example}{examples}
\newaliascnt{question}{theorem}
\newtheorem{question}[question]{Question}
\aliascntresetthe{question}
\crefname{question}{question}{questions}
\theoremstyle{remark}
\newaliascnt{remark}{theorem}
\newtheorem{remark}[remark]{Remark}
\aliascntresetthe{remark}
\crefname{remark}{remark}{remarks}
\crefname{part}{Part}{Parts}
\crefname{section}{section}{sections}
\crefname{subsection}{subsection}{subsections}
\crefname{appendix}{appendix}{appendices}
\crefname{table}{table}{tables}
\crefname{equation}{equation}{equations}
\crefname{enumi}{item}{items}
\crefname{lstlisting}{listing}{listings}

\newcommand{\No}{\mathbf{No}}
\newcommand{\Oz}{\mathbf{Oz}}
\newcommand{\On}{\mathbf{On}}
\newcommand{\N}{\mathbb N}
\newcommand{\Z}{\mathbb Z}
\newcommand{\Q}{\mathbb Q}
\newcommand{\R}{\mathbb R}
\newcommand{\Ralg}{\R_{\mathrm{alg}}}
\newcommand{\Hh}{\mathcal H}
\newcommand{\Dd}{\mathcal D}
\newcommand{\FH}{\operatorname{Frac}(\mathcal H)}
\newcommand{\SE}{\mathcal S_{\varepsilon_0}}
\newcommand{\QO}{\mathcal Q}
\newcommand{\IO}{\mathcal I}
\newcommand{\KG}{\mathcal K}
\newcommand{\sB}{\mathsf B}
\newcommand{\DF}{\operatorname{DF}}
\newcommand{\Frac}{\operatorname{Frac}}
\newcommand{\supp}{\operatorname{supp}}
\newcommand{\ct}{\operatorname{ct}}
\newcommand{\lc}{\operatorname{lc}}
\newcommand{\lt}{\operatorname{lt}}
\newcommand{\ddeg}{\operatorname{deg}_{\omega}}
\newcommand{\ord}{\operatorname{ord}}
\newcommand{\sgn}{\operatorname{sgn}}
\newcommand{\adj}{\operatorname{adj}}
\newcommand{\ev}{\operatorname{ev}}
\newcommand{\can}{\operatorname{can}}
\newcommand{\Ex}{\operatorname{Ex}}
\newcommand{\Fac}{\operatorname{Fac}}
\newcommand{\Geo}{\operatorname{Geo}}
\newcommand{\WF}{\mathrm{WF}}
\newcommand{\KB}{\mathrm{KB}}
\newcommand{\kb}{<_{\mathrm{KB}}}
\newcommand{\Pos}{\mathrm{Pos}}
\newcommand{\Eq}{\mathrm{Eq}}
\newcommand{\CK}{\omega_1^{\mathrm{CK}}}
\newcommand{\Om}{\Omega}
\newcommand{\tr}[1]{\tau_{\ge #1}}
\newcommand{\trs}[1]{\tau_{> #1}}
\newcommand{\fl}[1]{\left\lfloor #1\right\rfloor_{\Oz}}
\newcommand{\floor}[1]{\left\lfloor #1\right\rfloor}
\newcommand{\hb}{\bar h}
\newcommand{\src}[1]{\textup{\textsf{\small[#1]}}}
\newcommand{\mergetag}{\textup{\textsf{\small[merge]}}}
\newcommand{\replabel}[1]{\nolinkurl{#1}}
\newcommand{\fn}[1]{\nolinkurl{#1}}
\newcommand{\status}[1]{\par\smallskip\noindent\textit{Provenance.} #1\par\smallskip}
\newcolumntype{P}[1]{>{\raggedright\arraybackslash}p{#1}}
\lstset{basicstyle=\small\ttfamily,columns=fullflexible,breaklines=true,
 frame=single,rulecolor=\color{Gray},showstringspaces=false}

\begin{document}
\hypersetup{pageanchor=false}
\begin{titlepage}
\centering
\vspace*{0mm}
{\Large Merged research report for the Surreal project\par}
\vspace{6mm}
{\huge\bfseries Notations for Omnific Integers\par}
\vspace{5mm}
{\Large Holonomic towers, shielded and rational--omega calculi,\\
and the complexity of equality\par}
\vspace{5mm}
{\large Built from three manuscripts dated September 23, 2026\par}
\vspace{4mm}
\begin{minipage}{\textwidth}
\begin{abstract}
An omnific integer may have a transfinite Conway normal form, arbitrary real
coefficients at its positive exponents and an ordinary integer constant term.
Finite notation is therefore not the same as finite normal form, and decidable
syntax is not the same as decidable equality. This report studies which finite
representations of omnific integers have decidable validity, equality, order,
membership in $\Oz$ and omnific floor, and measures the failure of these
properties for the representations that lack them.

On the positive side three effective systems are constructed. Over any
effective ordered subfield of $\R$, certified Euler operators with finite
compatible jets name D-finite series, a leading-polynomial argument computes the
exceptional integer indices even over non-Archimedean coefficient fields, and
iterated fraction fields give effective finite-rank surreal fields whose
omnific parts have a unique triangular form, exact truncation, a computable
floor and a decidable ordinal slice; a union over hereditary Cantor labels has
ordinal trace exactly the ordinals below $\varepsilon_0$. A general shielding
theorem embeds $D+XL[X]$ into $\Oz$ whenever the monomial interpreting $X$
dominates the exponents of a Hahn field $L$, and transfers the equality problem
of $L$ many-one; with hereditary finite forms and $X=\varepsilon_0$ it gives a
decidable omnific ring with infinite normal forms, and finite matrix data
certify leading terms at arbitrary positive scales, with an exact
collision-aware form. Finite expressions in rational constants, field operations
and Conway's map $x\mapsto\omega^x$ compile into separated rational grids; this
decides membership in $\Oz$ and computes the floor, every truncation at a named
exponent and every coefficient. Its ordinals are exactly those below
$\varepsilon_0$, depth by depth, and its supports have order type below
$\omega^\omega$, sharply.

On the negative side, primitive-recursive coefficient streams on $d$ fixed
positive blocks always denote omnific integers, but their equality is
$\Pi^0_1$-complete and their positivity is many-one complete for the $d$-c.e.\
sets, already on values with at most $d$ monomials; no algorithm translates
these names into any system with decidable equality, although every value has a
rational name; for full-support rational switch streams the minimal realization
dimension is exactly $s+1$ and has no computable bound. With a decidable dyadic
support in $[1,2)$ and all coefficients one, support validity is
$\Pi^1_1$-complete.

This report merges three independently written manuscripts; it states which
source proved each result, prints shared results once, and keeps every
source's limitations. Several results were already in the collection and are
cited rather than reprinted. It is an AI-assisted draft: it has not been
refereed, and nothing in it has been formalized in Lean.
\end{abstract}
\end{minipage}
\vfill
\begin{minipage}{\textwidth}
\footnotesize
Sources: manuscripts 08 (the base text), 05 and 09 of the batch placed in
commit \texttt{a4dcb91}, all pinned to repository commit \texttt{a45d722};
their titles are in \cref{onot:sub:provenance}. The provenance of every
section is recorded in Appendix~\ref{onot:app:provenance}; non-claims are
collected in Appendix~\ref{onot:app:nonclaims}. The accompanying programs check
finite identities and examples only.
\end{minipage}
\end{titlepage}
\hypersetup{pageanchor=true}
\tableofcontents
\newpage

\section{Scope, results, and conventions}\label{onot:sec:intro}

An ordinal notation is a finite object with an ordinal interpretation.
Hereditary Cantor normal forms are the favourable case: validity and comparison
are decidable, and canonical expressions name exactly the ordinals below
$\varepsilon_0$. More general recursive ordinal systems make well-foundedness
part of the semantics, so recognizing valid expressions can itself be highly
undecidable. For omnific integers the question has more dimensions. A name must
describe real coefficients, surreal exponents and possibly transfinite supports,
and must then account for cancellation; an ordinal notation for the order type
of a support supplies none of these data. A system can be very rich in its
values while being algorithmically tame, or very restricted in its values while
having an undecidable equality relation. This report makes both possibilities
explicit. The ordinal-notation overview \cite{ordinalwiki} and Kleene's original
paper \cite{Kleene} motivate the distinction; the arguments below use the
underlying mathematical facts, not that overview, as their technical source
(08, 05, 09).

The positive constructions are three effective notation systems:
\emph{certified holonomic towers} (\cref{onot:hol:part}), \emph{hereditary finite
forms and an $\varepsilon_0$-shielded calculus} with finite-state series
(\cref{onot:sh:part}), and \emph{rational--omega expressions}
(\cref{onot:grid:part}). They share a general shielding theorem
(\cref{onot:sec:shield}). The negative constructions (\cref{onot:neg:part})
remove, one at a time, the certificates that make the positive systems work:
first the certificate that a coefficient stream vanishes identically, then the
certificate that a support is well ordered.

\begin{center}\small
\renewcommand{\arraystretch}{1.2}
\begin{tabular}{@{}P{.27\textwidth}P{.2\textwidth}P{.21\textwidth}P{.22\textwidth}@{}}
\toprule
Representation & Semantic validity & Equality & Order / positivity\\
\midrule
Certified holonomic towers (08) & Decidable & Decidable & Decidable\\
Hereditary forms; $\varepsilon_0$-shielded fractions; matrix series (05) &
Decidable & Decidable (finite certificates) & Decidable\\
Rational--omega expressions (09) & Decidable & Decidable & Decidable\\
$d$ primitive-recursive positive blocks (08; $d=1$: 05, 09) & Every
well-formed name is valid & $\Pi^0_1$-complete & $d$-c.e.-complete\\
Rational switch streams, coefficients $1,2$ (05) & Every name is valid &
$\Pi^0_1$-complete; no computable realization bound & Strict comparison c.e.\
(one block)\\
Tree-controlled dyadic supports (05, 08, 09) & $\Pi^1_1$-complete & Co-c.e.\
under the validity promise; checked equality $\Pi^1_1$-complete & Not needed for
the validity lower bound\\
\bottomrule
\end{tabular}
\end{center}
The last row distinguishes an equality question under a promise from a
question that must also check the promise (08). Here ``c.e.'' means
computably enumerable, equivalently semidecidable; a $\Pi^0_1$ predicate has the
form $\forall n\,R(n)$ with $R$ computable; a $\Pi^1_1$ predicate allows a
universal quantifier over functions $\N\to\N$ followed by an arithmetical
condition. Completeness is under computable many-one reductions, with the valid
input domains specified explicitly (05).

\subsection{Principal results}\label{onot:sub:results}

Bracketed numbers name the source manuscripts (\cref{onot:sub:provenance}).
Throughout, $k\subseteq\R$ is an effective ordered coefficient field, typically
$\Q$ or the real algebraic numbers $\Ralg$ (\cref{onot:sub:conventions}).

\begin{enumerate}[label=\textup{(\arabic*)},leftmargin=*,itemsep=3pt]
\item \textbf{Shielding and equality transfer} \src{05}. If $L\subseteq
k((t^G))$ is a subfield of a set-sized Hahn field, $\Z\subseteq D\subseteq
L\cap\Oz$, and $|g|<\Lambda$ for all $g\in G$, evaluation at $X=\omega^\Lambda$
embeds $D+XL[X]$ into $\Oz$ with unique coefficient lists and the sign of the top
block; with effective arithmetic, equality in the shielded ring is many-one
equivalent to equality in $L$ (\cref{onot:thm:shield,onot:thm:eqtransfer}). The
triangular theorem of the towers and the grid-ring theorem of the rational--omega
expressions are the characterizations of the whole omnific part in their fields
(\cref{onot:rem:characterizations}).
\item \textbf{Certified holonomic towers} \src{08}. Over an effective
leading-data field the nonnegative integer roots of a polynomial are computable
(\cref{onot:hol:thm:integerroots}); a certified Euler operator with a jet through
the last exceptional index names a unique D-finite series, with effective
validity, coefficients, zero test and ring operations
(\cref{onot:hol:thm:certificate,onot:hol:thm:dfclosure}); iterating fraction
fields gives countable subfields $F_r\subseteq\No$ with decidable equality and
order (\cref{onot:hol:thm:tower}), effective truncation
(\cref{onot:hol:thm:truncation}), triangular omnific parts
(\cref{onot:hol:thm:triangular}) and a computable floor
(\cref{onot:hol:thm:floor}). Supports have order type at most $\omega^r$, and
every $\eta<\omega^r$ is realized in $F_r\cap\Oz$ (\cref{onot:hol:thm:support});
the ordinals of $F_r$ are $\N[U_1,\ldots,U_r]$ (\cref{onot:hol:thm:ordinalslice});
the union over finite sets of hereditary Cantor labels has decidable equality,
order, recognition and floor, a computable canonicalization, and ordinal trace
$\On_{<\varepsilon_0}$ (\cref{onot:hol:thm:epsilon}).
\item \textbf{Hereditary forms and the $\varepsilon_0$-shielded calculus}
\src{05}. Finite hereditary forms over $k=\Q$ or $\Ralg$ have canonical trees,
decidable order and omnific membership, and a computable floor; their ordinals
are exactly those below $\varepsilon_0$, with natural arithmetic
(\cref{onot:sh:thm:hereditary,onot:sh:thm:floor,onot:sh:thm:ordinalfragment}).
The shielded ring $\SE=\Dd+\varepsilon_0\FH[\varepsilon_0]$ has decidable
validity, equality and order and contains infinite normal forms, including
supports of type $\omega^r$ for every $r$
(\cref{onot:sh:thm:epsilon,onot:sh:cor:omegar}); $\FH$ is not real closed
(\cref{onot:sh:prop:notrealclosed}).
\item \textbf{Finite fronts of matrix series} \src{05}. For commuting matrix
data of dimension $m$ at independent positive scales, the leading term occurs
at total degree below $m$; at arbitrary positive scales, including dependent and
non-Archimedean ones, the collected adjugate numerator certifies zero and the
leading term, with the sharp bound $(m-1)\sum\delta_i$
(\cref{onot:sh:thm:front,onot:sh:thm:collision}).
\item \textbf{Rational--omega expressions} \src{09}. For finite terms in
rational constants, $+,-,\times,/$ and $\Om(x)=\omega^x$, validity, equality and
order are decidable (\cref{onot:grid:thm:equality}); every term compiles into a
rational function of separated monomials (\cref{onot:grid:thm:grid}); membership
in $\Oz$, the floor, truncation at every named exponent and every coefficient
are computable (\cref{onot:grid:thm:effectivefloor}). Supports have order type
below $\omega^\omega$, sharply at depth two; an explicit computable strong sum
of such values escapes the language, and so does $\sum_n\omega^{1/n}$
(\cref{onot:grid:thm:supportbound,onot:grid:thm:sharp,onot:grid:thm:escape,onot:grid:prop:denominators}).
\item \textbf{Ordinal traces} \src{09; 08; 05}. At omega-nesting depth $d$ the
ordinals named are exactly those below $\beta_d$, where $\beta_0=\omega$ and
$\beta_{d+1}=\omega^{\beta_d}$ (\cref{onot:grid:thm:ordinaltrace}). The labelled
tower union, the hereditary ring and the rational--omega field all have ordinal
trace exactly $\On_{<\varepsilon_0}$, yet the omnific rings of the three systems
are pairwise incomparable, with one exception: for towers over $k=\Q$ it
is shown only that the rational--omega ring is not contained in the tower ring
(\cref{onot:grid:prop:incomparable}, \mergetag).
\item \textbf{Stream equality and no translation} \src{08; 05, 09}. For every
$d\ge1$, equality of raw primitive-recursive $d$-block names is
$\Pi^0_1$-complete, already on values with at most one monomial
(\cref{onot:neg:thm:rawequality}); no algorithm translates these names, or the
full-support rational switch names, into any system with decidable equality,
although every value has a finite rational name
(\cref{onot:neg:thm:notranslator,onot:neg:thm:rationalhard}). The two reductions
are those of \replabel{cas:thm-zerotest} in the computer-algebra report; what
is new is their omnific placement, the explicit translator corollaries, and the
exact minimal realization dimension $s+1$ with no computable bound
(\cref{onot:neg:thm:statebound}).
\item \textbf{Positivity} \src{08}. For each $d\ge1$ positivity of raw
$d$-block names is many-one complete for the $d$-c.e.\ sets, with at most $d$
monomials in the hard instances (\cref{onot:neg:thm:ershov}); this strengthens the
two-scale sign obstruction \replabel{cas:prop-leading}.
\item \textbf{Support validity} \src{05; 08, 09}. With an explicit dyadic
coding, validity of $\sum_{s\in T}\omega^{\kappa(s)}$ is $\Pi^1_1$-complete while
every support is a decidable subset of the dyadics in $[1,2)$
(\cref{onot:neg:thm:pionesone}); checked equality is $\Pi^1_1$-complete and
promised equality is co-c.e.\ (\cref{onot:neg:prop:treeequality}). This
strengthens the undecidability of \replabel{cas:prop-validation} to an exact
classification.
\end{enumerate}

These statements are proved here, but proving a statement here is not a claim
that nobody has previously proved it. Each source names its own proposed
contributions and credits its classical inputs
(Appendix~\ref{onot:app:nonclaims}); no exhaustive priority claim is made.

\subsection{Provenance of this report}\label{onot:sub:provenance}

This report is built from three manuscripts written independently and dated
September 23, 2026. They are items 08, 05 and 09 of the batch placed in commit
\texttt{a4dcb91}; their shipped file prefixes are \fn{08-holonomic-towers-},
\fn{05-hereditary-shielding-} and \fn{09-rational-omega-}. All three pin
repository commit \texttt{a45d722} (30 commits before that placement). At the
pin the collection already contained the computer-algebra report, the
computable-surreals report and the omnific Diophantine report with the labels
cited below; the manuscripts read only parts of them
(Appendix~\ref{onot:app:stale}).

\begin{center}
\renewcommand{\arraystretch}{1.2}
\begin{tabular}{>{\raggedright\arraybackslash}p{0.07\textwidth}>{\raggedright\arraybackslash}p{0.34\textwidth}>{\raggedright\arraybackslash}p{0.49\textwidth}}
\toprule
Source & Manuscript & Principal contribution to this report\\
\midrule
08 & \emph{Notations for Omnific Integers: Exact Holonomic Towers and the
Complexity of Equality} (31 pp.) & Base text and structure; notation systems and
checked equality; the integer-root algorithm; certified Euler series; holonomic
towers, truncation, triangular integer parts, floor, support spectrum, ordinal
slice and the $\varepsilon_0$-labelled union (\cref{onot:hol:part}); raw
$d$-block names, $\Pi^0_1$ equality, the translator theorem, the $d$-c.e.\
positivity classification, checked equality of tree supports.\\
05 & \emph{Equality for Omnific Notations} (26 pp.) & The shielding and
equality-transfer theorems (\cref{onot:sec:shield}); hereditary finite forms, the
$\varepsilon_0$-shielded calculus, finite and collision-aware fronts of matrix
series (\cref{onot:sh:part}); strict comparison; the rational-realization gap
with exact minimal dimension; the explicit dyadic coding with decidable
supports; effective boundedness.\\
09 & \emph{Effective Notations for Omnific Integers} (28 pp.) &
Rational--omega expressions: decidable equality, separated grids, grid-ring
theorem, recognition, floor, truncation at named exponents, per-depth ordinal
trace, support barrier, escape and lattice obstructions, recurrence and jet
certificates (\cref{onot:grid:part}); the $\CK$ barrier with its caveat.\\
\bottomrule
\end{tabular}
\end{center}

\paragraph{Where the merge had to choose.}
\begin{itemize}[leftmargin=*]
\item \emph{Base.} Among the proofs of the shared theorems, 08's have the
weakest hypotheses: any effective ordered coefficient field, an abstract
leading-data field, $d$ blocks where 05 and 09 treat one, an exact order
classification, and checked equality besides validity. Its text is the base,
with three exceptions printed from the members as main statements: 05's
shielding and transfer theorems as the general embedding, 05's dyadic coding
because only it gives decidable supports, and 09's per-depth ordinal trace.
\item \emph{Embedding versus characterization.} 05 proves only that $D+XL[X]$
embeds in $\Oz$; it does not prove that the image is the whole omnific part of
any field. That converse is 08's triangular theorem and 09's grid-ring theorem,
each in its own field. The three are printed as one general embedding and two
characterizations (\cref{onot:rem:characterizations}).
\item \emph{Printed once.} The omnific floor formula is
\replabel{odg:thm:floor} of the omnific Diophantine report and
\replabel{cas:eq-floor} of the computer-algebra report; it is cited
(\cref{onot:sub:floor}), and only the effectivity statements of the three
systems are proved here. 09's division with remainder and its Diophantine
transfer are \replabel{odg:prop:division} and \replabel{odg:cor:H10}, stronger
there; the computable part is kept (\cref{onot:grid:cor:division,onot:grid:prop:diophantine}).
The stream reductions are those of \replabel{cas:thm-zerotest}; 08's two-block
sign example is \replabel{cas:prop-leading}; 08's integer-root algorithm is the
effective form of \replabel{hol:lem:integer} of the holonomic-rigidity report;
09's grid fields reduce to \replabel{cas:thm-core} after compilation. The
equality theorems for one stream block, the no-translator theorems, the
canonicalization observation, the size obstruction and the $\CK$ barrier, each
proved by two or three sources, are printed once.
\item \emph{Second routes.} 08's and 09's midpoint codings of the
Kleene--Brouwer order are kept as second routes to the validity theorem; they
do not give decidable supports (\cref{onot:neg:rem:secondroute}).
\item \emph{Corrections.} 08's statement of its promised-equality proposition
says ``a co-c.e.\ test for inequality''; its proof and its own summary table
show that inequality is c.e., so promised equality is co-c.e.; the statement is
corrected (\cref{onot:neg:prop:treeequality}).
\end{itemize}
Results added by the merge carry the tag \mergetag{} and have complete proofs.
Appendix~\ref{onot:app:provenance} records, section by section, which source
proved what and every symbol renamed from a source.

\subsection{What ``decidable'' means here}\label{onot:sub:decidable}

All algorithms are ordinary terminating algorithms on finite strings, with no
oracle unless explicitly stated. A procedure need not have a practical
complexity bound to be decidable. The coefficient field at the bottom is part of
the input convention, not an oracle for real numbers: it is an effectively
presented ordered subfield $k\subseteq\R$ whose valid codes, equality, order and
field operations are decidable or computable. Typical choices are $\Q$ and the
real algebraic numbers $\Ralg$ with exact algebraic codes, for instance a
normalized irreducible primitive polynomial over $\Z$ with the index of the root
in increasing order; isolating intervals may accompany this pair as certificates
but are not part of its canonical identity, and standard polynomial and
real-root algorithms provide the field interface (05). Arbitrary Cauchy names for
all computable reals do not satisfy this equality contract
(\cref{onot:neg:prop:realcoeff}). Some results need more: the floor needs an
effective integer test and ordinary floor in $k$, and the separated-basis lemma
of \cref{onot:grid:part} needs $k=\Q$.

Every set and sum in the proofs is set-sized. The proper class $\No$ is used
only as an ambient semantic domain; in a universe-relative formalization one may
choose a universe containing the countable fields and supports under
consideration, and no truth predicate for the entire universe is required (08).
All index sets and supports are sets; no proper-class sum is formed (09).

\subsection{Notation and conventions}\label{onot:sub:conventions}

\paragraph{Sign convention.} Exponents are those of Conway normal forms:
supports are reverse well ordered and larger exponents correspond to larger
monomials. In the valuation convention $t^\gamma=\omega^{-\gamma}$, Hahn series
have increasing well-ordered supports; the realization
$\sum_ga_gt^g\mapsto\sum_ga_g\omega^{-g}$ carries a minus sign, which is
essential (05, 09). All three sources use this growth convention, as do the
omnific reports of the collection.

\paragraph{One meaning per symbol.} The sources use several letters for
different notions. The following choices are fixed for the whole report; the
complete list of renamings is \cref{onot:tab:renames}.
\begin{itemize}[leftmargin=*]
\item $\ddeg x=\max\supp x$ is the leading (greatest) exponent of $x\neq0$ and
$\lc(x)$ its coefficient, as in the omnific Diophantine report. 08 wrote
$\ell(x)$ and 09 $\lambda(x)$ for this. The order type of $\supp x$, read in
decreasing order, is $\ord(\supp x)$; 09 wrote $\ell(x)$ for that, so the letter
$\ell$ is not used for either notion. In the $t$-convention the least exponent of
a Hahn series $\Phi$ is its valuation $v(\Phi)$ (05's $\lambda$).
\item $\varepsilon_0$ is always written out. 05 abbreviated it $E$; 08's
exponential series $E(t)$ is $\Ex(t)$ here.
\item $k$ is the effective ordered coefficient field inside $\R$ (08's $C$, 05's
$K$); $K$ is 08's effective leading-data field, a subfield of $\No$. $\Ralg$ is the
field of real algebraic numbers (05's $\mathbb A$).
\item The code set of a notation system is $N$ (08 and 09 wrote $D$). In the
shielding theorem $D$ is the constant ring and $L$ the coefficient field (05
wrote $F$); $\Dd=\Hh\cap\Oz$ is the hereditary omnific core.
\item The three fields: $F_r$, $\mathcal F_A$, $\mathcal F_\varepsilon$ (08, with
omnific parts $R_r$, $\mathcal R_A$, $\mathcal R_\varepsilon$); $\Hh$, $\FH$ and
the ring $\SE$ (05; 05 wrote $\mathcal F$ for $\FH$ and $\mathcal C$ for $\SE$);
the rational--omega field $\QO$ with depth levels $\QO_d$ and omnific part
$\IO=\QO\cap\Oz$ (09's $\mathcal F$, $\mathcal F_d$, $\mathcal A$).
\item $\tr g$ and $\trs g$ are the truncations keeping the exponents $\ge g$,
respectively $>g$ (08's $\operatorname{Tr}$, 09's $\tau$).
\item $h_e(n)=1$ exactly when machine $e$ halts \emph{for the first time} at
step $n$ (08, 09); $\hb_e(n)=1$ exactly when $e$ has halted \emph{within} $n$
steps (05, who wrote $h_e$). These are Reductions B and A of
\replabel{cas:thm-zerotest}.
\item $\sB=\omega^\omega/(1-\omega^{-1})=\sum_{n\ge0}\omega^{\omega-n}$ is the
same element in all three sources (05's $B$, 09's $B$, 08's $U\Geo(t)$).
\end{itemize}

\paragraph{Promises.} $\Eq_\nu$ is equality on valid names (a promise
problem), $\Eq^{\mathrm{chk}}_\nu$ is \emph{checked} equality, which must also
verify validity (\cref{onot:def:notation}). A promise algorithm may diverge or
answer wrongly on inputs violating the promise (08).

\subsection{Relationship with the repository and the literature}\label{onot:sub:repo}

Conway normal forms and the characterization of omnific integers are
classical \cite{Conway,Gonshor}; L'Innocente and Mantova develop a factorization
theory for generalized power series and omnific integers and settle a specific
primality conjecture \cite{LM}. The subject here is not factorization but finite
representations and their decision problems; no claim about the decidability of
factorization, of elementary theories or of Diophantine solvability is inferred
from the equality algorithms (05). Lurie studied the effective content of surreal
algebra \cite{Lurie}; Hamkins's account of joint work with Turetsky describes a
much larger notion of computable surreal and records its rediscovery of Lurie's
work \cite{Hamkins}. The fields here are deliberately smaller: the real trace of
the rational--omega field is only $\Q$, and ``computable'' always refers to a
stated representation (09). D-finite power-series algebra is classical
\cite{Stanley}, and Chen, Fang and van der Hoeven give a zero test for
D-algebraic transseries \cite{CFH}; the towers of \cref{onot:hol:part} are not
presented as the first zero test beyond rational functions (08).

Within the collection, the computable-surreals report
(\fn{docs/foundations-and-computation/computable-surreals}) separates structural
cut names from numerical series names and stresses that equality and
normalization must not be inferred from abstract algebraic closure; its no-go
theorem says a successful alternative representation must, among other options,
use a more informative real representation, which is the route all three systems
here take. The computer-algebra report
(\fn{docs/foundations-and-computation/computer-algebra}, prefix \texttt{cas:})
already contains the effective rational monomial core, both halting reductions
for coefficient streams, a two-scale sign obstruction and the $\Pi\oplus\Z$
floor. The omnific Diophantine report (\fn{docs/surreal/omnific-diophantine-geometry},
prefix \texttt{odg:}) contains the floor, division with remainder and the
Hilbert's-tenth-problem transfer. The precise relations are in
\cref{onot:sec:neighbours}.

\section{Normal forms, notation systems, and the floor}\label{onot:sec:normal}

\subsection{The semantic domain}

We use Conway's monomial map $x\mapsto\omega^x$, an order-preserving
homomorphism from $(\No,+)$ to the positive monomials with
$\omega^{x+y}=\omega^x\omega^y$ and $\omega^0=1$; it is not an unspecified
extension of ordinary exponentiation and is not defined as $\exp(x\log\omega)$
(05, 08, 09). Every surreal has a unique normal form
\begin{equation}\label{onot:eq:normal}
 x=\sum_{\gamma\in S}c_\gamma\omega^\gamma,
 \qquad c_\gamma\in\R\setminus\{0\},
\end{equation}
where $S=\supp x\subseteq\No$ is a set and is reverse well ordered: each
nonempty subset has a greatest element. For $x\ne0$ put $\ddeg x=\max S$ and
$\lc(x)=c_{\ddeg x}$; the sign of $x$ is the sign of $\lc(x)$, and
\begin{equation}\label{onot:eq:leading}
 \ddeg(xy)=\ddeg x+\ddeg y,\qquad \lc(xy)=\lc(x)\lc(y).
\end{equation}
No leading exponent is assigned to zero. Addition collects equal exponents;
multiplication uses $\omega^\gamma\omega^\eta=\omega^{\gamma+\eta}$ and a finite
convolution at each exponent. These operations are justified by the standard
support lemmas for generalized power series \cite{Conway,Gonshor,LM}. For a
set-sized ordered additive subgroup $G\subseteq\No$ we write $k((t^G))$ for the
Hahn field of series with well-ordered support in the $t$-convention; we use the
Hahn support and product theorem: sums of well-ordered supports are well
ordered, and the coefficient fibres of a product are finite. For a finite set of
strictly positive $g_i$ the nonnegative integer monoid they generate is well
ordered and each of its elements has finitely many representations
$\sum n_ig_i$; this is the positive-support form of Neumann's lemma (05; see
also \cite[Proposition~3.7]{Camacho}).

\paragraph{Strong sums are not limits.} An indexed family of normal forms is
strongly summable if the union of the supports is reverse well ordered and only
finitely many summands contribute at each exponent; its sum is then defined
coefficientwise. Every infinite sum in this report is meant in this sense; its
existence is a support assertion, not convergence of partial sums in the order
topology of $\No$, and no such topological convergence is needed anywhere
(05, 08, 09).

\begin{definition}[Omnific integers]\label{onot:def:oz}
The class of omnific integers is
\begin{equation}\label{onot:eq:oz}
 \Oz=\{x\in\No:\ \supp x\subseteq\No_{\ge0}\ \text{and}\ \ct(x)\in\Z\},
\end{equation}
where $\ct(x)$ is the coefficient of $\omega^0$, read as zero when $0\notin\supp
x$. Equivalently $z=m+\sum_{\gamma\in S}c_\gamma\omega^\gamma$ with $m\in\Z$ and
$S\subseteq\No_{>0}$ reverse well ordered. The positive-exponent coefficients
need not be ordinary integers.
\end{definition}

This is the classical criterion \cite{Conway,Gonshor,LM}, and it agrees with
the collection's $\Oz=\Pi\oplus\Z$ (\replabel{found:eq:omnific}). For
example $\sqrt2\,\omega-\tfrac13\omega^{1/2}+7$, $\omega-1$ and $\omega^{1/2}$
are omnific integers, whereas $\tfrac12$, $\omega+\tfrac12$ and
$\omega\pm\omega^{-1}$ are not (08, 05). The support criterion shows that $\Oz$
is a subring of $\No$: a sum of nonnegative exponents is zero only if both are
zero, so $\ct:\Oz\to\Z$ is a ring homomorphism which is the identity on $\Z$
(09). $\Oz$ is discretely ordered: a positive element with nonempty positive
support exceeds every real number, and otherwise it is a positive ordinary
integer, so $1$ is its least positive element. These are elementary classical
consequences, not new theorems (09).

It is useful to write
\begin{equation}\label{onot:eq:threeparts}
 x=x_{>0}+\ct(x)+x_{<0},
\end{equation}
where $x_{>0}$ collects the terms with positive exponents and $x_{<0}$ those with
negative exponents. The constant term is an ordinary real number, $x_{<0}$ is
zero or infinitesimal ($|x_{<0}|<1/n$ for every positive integer $n$), and
$x_{>0}$ is zero or purely infinite. This is a normal-form decomposition;
$x_{>0}$ is not generally the positive part in the order-theoretic sense (08).

\subsection{The omnific floor}\label{onot:sub:floor}

For every surreal $x$ there is exactly one $z\in\Oz$ with $z\le x<z+1$, namely
\begin{equation}\label{onot:eq:floor}
 \fl{x}=x_{>0}+\floor{\ct(x)}-
 \begin{cases}
 1,&\ct(x)\in\Z\text{ and }x_{<0}<0,\\
 0,&\text{otherwise.}
 \end{cases}
\end{equation}
This is the integer-part formula \replabel{odg:thm:floor} of the omnific
Diophantine report, credited there to Ehrlich and Kaplan, and the
$\Pi\oplus\Z$ floor \replabel{cas:eq-floor} of the computer-algebra report
(whose ring $\Pi\oplus\Z$ is $\Oz$); both reports prove it, and all three sources
prove it again (05's Theorem~3.2, 08's Theorem~7.4, 09's equation (8)). It is
printed once, here. The omnific continued-fraction report proves it again as
\replabel{ocf:lem:floor} (batch 33), where it is the digit map of the
continued-fraction algorithm. If $\ct(x)$ is not an integer, its distances from the
adjacent integers are positive reals which an infinitesimal cannot cross; if it
is an integer, the sign of $x_{<0}$ selects the side. The correction is not
optional:
\begin{gather*}
 \fl{-\omega+3-\omega^{-1}}=-\omega+2,\qquad
 \fl{\omega-\omega^{-1}}=\omega-1,\\
 \fl{\omega+\tfrac32-\omega^{-1}}=\omega+1,\qquad
 \fl{\omega+\tfrac12-\omega^{-1}}=\omega,
\end{gather*}
and deleting the negative-exponent terms and flooring the constant gives the
wrong answer in the first two cases (08, 05, 09). What each system of this report
adds is \emph{effectivity}: its names allow $x_{>0}$, $\ct(x)$ and the sign of
$x_{<0}$ to be computed as names of the same system
(\cref{onot:hol:thm:floor,onot:sh:thm:floor,onot:grid:thm:effectivefloor}).
The computer-algebra report records that the operation is effective only when
the backend can split the normal form, decide the ordinary integer predicate
and determine the infinitesimal sign; the three systems supply exactly these.

\subsection{Notation systems, validity, and canonical codes}

\begin{definition}[Notation system]\label{onot:def:notation}
A notation system consists of a decidable set $N$ of finite syntactic codes, a
subset $V\subseteq N$ of semantically valid codes, and a denotation map
$\nu:V\to\Oz$ (or into $\No$ for an ambient field). Its equality problem is
the promise problem
\[
 \Eq_\nu=\{(a,b)\in V^2:\nu(a)=\nu(b)\}
\]
on valid inputs. Its \emph{checked equality} relation on all syntactic inputs is
\[
 \Eq^{\mathrm{chk}}_\nu=\{(a,b)\in N^2:\ a,b\in V\text{ and }\nu(a)=\nu(b)\}.
\]
The system is total if $V=N$. A canonicalization is a computable map
$\can:V\to V$ with $\nu(\can(a))=\nu(a)$, and $\can(a)=\can(b)$ exactly when
$\nu(a)=\nu(b)$.
\end{definition}

Four problems are logically distinct: recognizing valid codes, comparing
values, producing a preferred code for a value, and deciding whether a value
belongs to a specified subring; a promise such as ``this program is total'' is
not a validity algorithm (09). A parser may decide membership in a syntactic
language without deciding semantic validity; a comparison procedure may be
required to work only under a validity promise; a code for an exact coefficient
is different from a code for approximations to it. These alternatives change the
decision problem and are never left implicit (05). A coefficient-access
algorithm computes a specified coefficient or a specified finite collection of
coefficients; it does not by itself certify that all remaining coefficients
vanish, and a support generator is not automatically a support-validity
certificate (08).

\begin{proposition}[Canonicalization and decidable quotients \src{08; 05, 09}]\label{onot:prop:canonical}
For a total notation system with effectively coded finite strings, equality is
decidable if and only if it has a computable canonicalization.
\end{proposition}
\begin{proof}
A canonicalizer decides equality by comparing its outputs. Conversely, number
all finite strings by natural numbers. For an input code $a$, test the finitely
many codes $b\le a$ for syntactic validity and for semantic equality to $a$, and
return the least successful one. The tested set contains $a$, so it is nonempty.
The result is the least code in the whole equivalence class, since no code larger
than $a$ can be the least; equal inputs therefore give the same result.
\end{proof}
This supplies a canonical form in the computability sense, not a useful
algebraic expression or an efficient normalizer (08). A canonical denotation in
the mathematical sense does not supply an algorithm to find its canonical code;
the least-code search may be useless computationally, but it separates existence
from efficient normal-form design (05). A reduced rational function is canonical
only relative to a fixed exponent coordinate system (09).

\begin{proposition}[Size obstruction \src{08, 05, 09}]\label{onot:prop:size}
No notation system over a finite or countable alphabet names every omnific
integer, or even every $c\omega$ with $c\in\R$. More generally, the image of any
set of names is set-sized.
\end{proposition}
\begin{proof}
Finite strings over a countable alphabet form a countable set, whereas
$c\mapsto c\omega$ is injective on $\R$. Moreover $\Oz$ contains every ordinal
and is a proper class; replacement gives the general assertion.
\end{proof}
This obstruction is not an undecidability theorem, and it does not prevent
large useful countable subrings. A language containing a symbol for a very large
specified ordinal may have trivial equality on that symbol; restrictions
involving $\CK$ concern computably presented well-orders, not the largest ordinal
one may mention by a finite semantic convention (05).

\subsection{A sparse baseline}\label{onot:sub:sparse}

Let $G\subseteq(\No,+,<)$ be an effectively presented ordered abelian group
with fixed embedding and group operations. Over $k$, use finite records
\begin{equation}\label{onot:eq:sparse}
 (m;(\gamma_1,c_1),\ldots,(\gamma_s,c_s)),\qquad m\in\Z,\ \gamma_i>0,
\end{equation}
allowing zero coefficients and repeated exponents in raw records; sort the
exponents decreasingly, collect equal exponents and delete zero coefficients.

\begin{proposition}[Sparse normalization \src{08}]\label{onot:prop:sparse}
The records \eqref{onot:eq:sparse} have decidable validity, equality and order,
and computable ring operations. After collection and sorting, their list of
exponent--coefficient values and their integer constant are unique.
\end{proposition}
\begin{proof}
Only finitely many comparisons and coefficient operations occur. Normal-form
uniqueness gives equality and uniqueness. The sign of the first nonzero term
decides order; if no positive term remains, the integer constant decides it.
Products have exponents in the finite sumset of the input supports, and
positive exponents stay positive.
\end{proof}
If coefficient and exponent codes are not themselves canonical, this is
uniqueness of values rather than of literal strings. The sparse core is
familiar generalized-polynomial algebra and already occurs in the
computer-algebra report (\replabel{cas:thm-core}) (08). Adding infinite series is
not adding longer lists:
\begin{equation}\label{onot:eq:guardbasic}
 \sum_{n\ge0}a_n\omega^{\omega-n}=\omega^\omega\sum_{n\ge0}a_n(\omega^{-1})^n
\end{equation}
always has positive exponents and a valid support, whatever the real $a_n$, but
an algorithm for each $a_n$ need not decide whether every $a_n$ vanishes. A
finite certificate can supply the missing information (\cref{onot:hol:part}),
including the singular initial indices, which cannot simply be ignored (08).
The holonomic-rigidity report reads the series \eqref{onot:eq:guardbasic} as a
recurrence observable (batch 33). Starting from $y_0$ equal to that series,
the iterates $y_n=\omega^n y_0$ under multiplication by $\omega$ stay purely
infinite, and the coefficient of
$\omega^\omega$ in $y_n$ is $a_n$ (\replabel{hol:rc:thm:coding}). So the
fixed-coefficient zero sets of first-order omnific recurrences realize every
subset of the natural numbers, and no coefficientwise
Skolem--Mahler--Lech theorem holds for them (\replabel{hol:rc:cor:nosml}).
That report credits the series to this equation.

\section{Shielding and the transfer of equality}\label{onot:sec:shield}

This section separates the source of equality difficulty from the fact that the
output is an omnific integer. The ring form $D+XL[X]$ is a familiar polynomial
pullback construction; the concern here is its surreal realization and its
effective equality (05).

Let $k\subseteq\R$ be an ordered field, $G\subseteq\No$ a set-sized ordered
additive subgroup, and $L\subseteq k((t^G))$ a subfield, realized in $\No$ by
$\sum a_gt^g\mapsto\sum a_g\omega^{-g}$. Let $D\subseteq L\cap\Oz$ be a subring
containing $\Z$. Choose a positive surreal $\Lambda$ with
\begin{equation}\label{onot:eq:shieldbound}
 |g|<\Lambda\qquad(g\in G).
\end{equation}
Such a bound exists because $G$ is a set; its existence is not a uniform
algorithm for extracting a bound from arbitrary programs, and in an effective
application $\Lambda$ and a justification of the bound are part of the chosen
presentation. If $G\subseteq\R$, the fixed choice $\Lambda=\omega$ works. Put
$X=\omega^\Lambda$ and
\begin{equation}\label{onot:eq:shieldring}
 \mathcal S_\Lambda(D,L)=\Bigl\{d+\sum_{j=1}^mX^jf_j:\ d\in D,\ f_j\in L,\ m\in\N\Bigr\}.
\end{equation}
Here $X$ is a surreal; in the abstract polynomial ring below, $X$ is an
indeterminate.

\begin{theorem}[Shielding and unique degree separation \src{05}]\label{onot:thm:shield}
Evaluation at $X=\omega^\Lambda$ induces an injective ring homomorphism
\[
 D+XL[X]\hookrightarrow\Oz
\]
with image \eqref{onot:eq:shieldring}. Every element has a unique finite
coefficient list, up to trailing zeros. If its largest nonzero positive degree
is $m$, its sign is the sign of $f_m$; if no positive degree survives, its sign
is that of $d$.
\end{theorem}
\begin{proof}
Expand $f_j=\sum_ga_{j,g}\omega^{-g}$. Then
\begin{equation}\label{onot:eq:shieldexpand}
 X^jf_j=\sum_ga_{j,g}\omega^{j\Lambda-g}.
\end{equation}
For $j\ge1$ every exponent is strictly positive by \eqref{onot:eq:shieldbound},
and the support is reverse well ordered because $g\mapsto j\Lambda-g$ reverses
the well-ordered Hahn support; a finite union of such supports is reverse well
ordered. The degree-zero element $d$ is omnific by assumption, so the evaluated
expression is omnific.

If $j>i$ and $g,h\in G$, then
\[
 (j\Lambda-g)-(i\Lambda-h)=(j-i)\Lambda-(g-h)>0,
\]
since $g-h\in G$. Thus every exponent in a higher degree exceeds every exponent in
a lower degree, including degree zero, whose exponents lie in $-G$. Different
degree blocks cannot cancel, and normal-form uniqueness gives injectivity and
uniqueness of the coefficient list. The leading coefficient of the highest
nonzero block is that of $f_m$, which gives the sign. The finite polynomial
identities together with Hahn multiplication show that evaluation is a ring
homomorphism.
\end{proof}

For $D=\Z$ the omnific constant term is simply $d$; in general it is $\ct(d)$.
The construction does \emph{not} embed $L$ as a subfield of $\Oz$, since a
discretely ordered ring cannot contain $\tfrac12$; rather $f\mapsto Xf$ is an
injective additive encoding with $(Xf)(Xg)=X^2fg$ (05).

\begin{remark}[Embedding and characterization \mergetag]\label{onot:rem:characterizations}
\Cref{onot:thm:shield} proves one inclusion: the image of $D+XL[X]$ lies in
$\Oz$. It does not say that this image is the whole omnific part of a field
containing $L$ and $X$. Two such characterizations are proved in this report,
each in its own field and each by an argument not contained in
\cref{onot:thm:shield}.
\begin{itemize}[leftmargin=*]
\item In the holonomic tower (\cref{onot:hol:thm:triangular}) take
$L=F_{r-1}$, $D=R_{r-1}$, $G=G_{r-1}$ and $\Lambda=\delta_r$: every element of
$G_{r-1}$ is dominated by $\delta_r$, so \eqref{onot:eq:shieldbound} holds and
$\mathcal S_{\delta_r}(R_{r-1},F_{r-1})=R_{r-1}+U_rF_{r-1}[U_r]$. The theorem
there adds $F_r\cap\Oz=R_r$ equal to this ring.
\item In a separated rational grid (\cref{onot:grid:thm:gridring}) take $L$ the
tail field $\Q(X_2,\ldots,X_r)$, $D$ its omnific part and $\Lambda=\gamma_1$;
since $\gamma_1\gg\gamma_2\gg\cdots$, every element of the tail lattice is
dominated by $\gamma_1$. The theorem there adds that the omnific part of
$\Q(X_1,\ldots,X_r)$ equals this ring.
\end{itemize}
The converse inclusion is where both proofs use the field: an element of the
field whose top Laurent block has positive index, or whose grid expansion has a
negative first coordinate, has a negative exponent and so is not omnific.
\end{remark}

\subsection{Exact transfer of equality}

Assume $L$ has a system of valid finite names on which $+$, $-$ and
multiplication are uniformly computable, and that $D$-names translate computably
to $L$-names. No equality algorithm for $L$ is assumed. Polynomial-list names for
the shielded ring allow zero coefficients; they need not be normalized.

\begin{theorem}[Many-one equality transfer \src{05}]\label{onot:thm:eqtransfer}
Under these representation contracts the valid-name promise problems satisfy
\[
 \Eq_{\mathcal S_\Lambda(D,L)}\equiv_m\Eq_L,
\]
by computable maps on promised valid inputs. If both naming domains are
decidable, these are ordinary total many-one reductions, invalid strings being
rejected, and decidability and every many-one completeness classification of
$\Eq_L$ transfer to the shielded omnific notation system. Without decidable
validity the assertion concerns promised equality, not the complexity of
recognizing valid names.
\end{theorem}
\begin{proof}
For the forward reduction map $(f,g)$ to $(Xf,Xg)$; injectivity in
\cref{onot:thm:shield} proves equivalence. Conversely, pad two coefficient lists
with zeros to the same length, put $h_0=d-e$ and $h_j=f_j-g_j$ for $j\ge1$, and
compute the single coefficient-field expression
\begin{equation}\label{onot:eq:squares}
 H=h_0^2+h_1^2+\cdots+h_m^2.
\end{equation}
In an ordered field $H=0$ exactly when every $h_j=0$, and degree separation
identifies this with equality of the two omnific values. The output pair $(H,0)$
is a single instance of $\Eq_L$, not a list of oracle queries; only the
stipulated arithmetic and translation operations were used. If the naming
domains are decidable, first reject invalid inputs and map them to a fixed
unequal pair of valid names; this extends the promised reductions to total
reductions of the ordinary equality sets.
\end{proof}

The field $L$ is ordered, being a subfield of the ordered Hahn field
$k((t^G))$ with $k\subseteq\R$, so the sum-of-squares step needs no hypothesis
beyond those stated. Over an arbitrary coefficient field, individual coefficient
tests still give a finite truth-table reduction, but this many-one argument may
fail (05). If equality and order in $L$ are decidable, comparison in the shielded
ring follows by finding the highest nonzero degree and applying the sign rule;
the abstract existence of $L$ does not by itself make its raw-name equality
decidable.

\part{Certified holonomic towers}\label{onot:hol:part}

This part is source 08's. It builds effective fields inside $\No$ by iterating
fraction fields of certified D-finite series, and describes their omnific parts,
floors, supports and ordinals exactly. The base field is an effective ordered
subfield $k\subseteq\R$ (\cref{onot:sub:decidable}).

\section{Computing integer singularities over non-Archimedean fields}\label{onot:hol:sec:integerroots}

The recursive construction needs to solve one modest problem: given
$P(X)\in K[X]\setminus\{0\}$, compute all nonnegative ordinary integers $n$ with
$P(n)=0$, although $K$ is non-Archimedean. Bounding roots by the absolute values
of the coefficients does not solve this: the bound may be an infinite surreal
with no ordinary integer above it.

\begin{definition}[Effective leading-data field]\label{onot:hol:def:leadingdata}
An effective leading-data field over $k$ is an effectively presented ordered
subfield $K\subseteq\No$ with support in an effective ordered exponent group $G$,
together with algorithms that compute $\ddeg a\in G$ and $\lc(a)\in k$ for
$a\ne0$.
\end{definition}
At rank zero, $K=k$: every nonzero element has exponent zero and is its own
leading coefficient.

\begin{lemma}[Archimedean integer search \src{08}]\label{onot:hol:lem:archsearch}
Given $c\in k$, an ordinary integer greater than $|c|$ can be found effectively,
and $\floor c\in\Z$ can be computed.
\end{lemma}
\begin{proof}
Test $2^j>|c|$ for $j=0,1,\ldots$; Archimedeanness guarantees termination.
Within the resulting finite integer interval, exact comparisons determine the
greatest integer at most $c$; binary search is optional.
\end{proof}

\begin{theorem}[Leading-polynomial integer-root algorithm \src{08}]\label{onot:hol:thm:integerroots}
Let $K$ be an effective leading-data field over $k$. There is an algorithm that,
given $0\ne P\in K[X]$, returns the finite set
\[
 I(P)=\{n\in\N:P(n)=0\},
\]
and in particular computes $B(P)=\max(I(P)\cup\{-1\})$.
\end{theorem}
\begin{proof}
Write $P(X)=\sum_{i=0}^ma_iX^i$, omitting zero coefficients when taking leading
exponents, and put
\[
 g=\max_{a_i\ne0}\ddeg a_i,\qquad p(X)=\sum_{\ddeg a_i=g}\lc(a_i)X^i\in k[X].
\]
The polynomial $p$ is nonzero, since its contributions sit at different powers
of $X$. For each ordinary $n\ge0$ the coefficient of $\omega^g$ in $P(n)$ is
exactly $p(n)$, so
\begin{equation}\label{onot:hol:eq:rootcontain}
 P(n)=0\ \Longrightarrow\ p(n)=0,
\end{equation}
including $n=0$. If $p$ is a nonzero constant, return the empty set. Otherwise
write $p(X)=b_dX^d+\cdots+b_0$ with $b_d\ne0$ and choose an integer
$M>1+\max_{i<d}|b_i/b_d|$ by \cref{onot:hol:lem:archsearch}. By the elementary
Cauchy bound $p$ has no real root of absolute value at least $M$: for
$|x|>1+\max_{i<d}|b_i/b_d|$ the lower-degree terms have total absolute value
below $|b_dx^d|$ by the geometric-series estimate. Thus
\eqref{onot:hol:eq:rootcontain} reduces the problem to $n=0,\ldots,M$; evaluate
$P(n)$ in $K$ and keep exactly its zeros.
\end{proof}

The theorem does not compute all roots in $K$ and does not require $K$ to be
real closed. The leading polynomial is only a filter: cancellation at one of its
roots can leave a nonzero lower-scale term, which is why the final exact tests
are needed.

\begin{remark}[Relation to the holonomic-rigidity report \mergetag]\label{onot:hol:rem:holinteger}
The same leading-polynomial reduction is \replabel{hol:lem:integer} (``integer
evaluation'') of the holonomic-rigidity report
(\fn{docs/surcomplex/holonomic-rigidity-for-entire-hahn-functions}), which
shows that $P(n)\ne0$, with valuation equal to the least coefficient valuation,
for all but finitely many $n$. That lemma is not effective: it does not bound the
exceptional $n$. \Cref{onot:hol:thm:integerroots} adds the Cauchy bound on the
reduced polynomial and the exact tests, which make it an algorithm over any
effective leading-data field. 08 did not cite the holonomic-rigidity report; the
credit is the merge's.
\end{remark}

\begin{example}\label{onot:hol:ex:integerroot}
In a field containing $t=\omega^{-1}$ let $P(X)=(X-3)\omega+(X-2)$. Its leading
polynomial is $X-3$; the only candidate is $3$, but $P(3)=1$, so
$I(P)=\varnothing$. By contrast $(X-3)(\omega+1)$ has the same leading polynomial
and the integer root $3$.
\end{example}

\section{Certified differential-series names}\label{onot:hol:sec:certificates}

Fix an effective leading-data field $K$ and an independent formal variable $t$.
Differentiation with respect to $t$ kills $K$; put $\theta=t\,d/dt$.

\subsection{Euler-normalized operators}

An Euler-normalized operator is
\begin{equation}\label{onot:hol:eq:euleroperator}
 \mathsf L=\sum_{j=0}^Jt^jP_j(\theta),\qquad P_j\in K[X],\ P_0\ne0.
\end{equation}
Every nonzero linear differential operator over $K(t)$ can be put into this form
without changing its power-series solutions: clear rational-function
denominators, use
\begin{equation}\label{onot:hol:eq:falling}
 t^a\Bigl(\frac d{dt}\Bigr)^b=t^{a-b}\theta(\theta-1)\cdots(\theta-b+1),
\end{equation}
collect equal powers of $t$, and multiply by a power of $t$ to make the least
occurring power zero. Multiplying the equation by a nonzero rational function
does not change its solution set in $K((t))$, and the triangular change of basis
between powers and falling factorials keeps a nonzero operator nonzero. For
$f=\sum_{n\ge0}a_nt^n$ the equation $\mathsf Lf=0$ is equivalent to
\begin{equation}\label{onot:hol:eq:recurrence}
 P_0(n)a_n+\sum_{j=1}^JP_j(n-j)a_{n-j}=0\quad(n\ge0),\qquad a_m=0\ (m<0).
\end{equation}
The argument of $P_j$ is $n-j$, not $n$.

\begin{definition}[Certified holonomic name]\label{onot:hol:def:certificate}
A certified holonomic name over $K$ is an operator
\eqref{onot:hol:eq:euleroperator} together with the finite jet
$(a_0,\ldots,a_B)$, $B=B(P_0)$, empty if $B=-1$, such that every equation
\eqref{onot:hol:eq:recurrence} with $0\le n\le B$ holds. A well-formed but
incompatible jet is invalid and is rejected.
\end{definition}

\begin{theorem}[Finite certification, zero testing, and leading data \src{08}]\label{onot:hol:thm:certificate}
A certified holonomic name has a unique denotation in $K[[t]]$. Validity,
coefficient extraction and the zero test are effective; explicitly,
\begin{equation}\label{onot:hol:eq:jetzero}
 f=0\iff a_0=\cdots=a_B=0.
\end{equation}
For a nonzero denotation the least nonzero index is the least nonzero index in
the jet, hence at most $B$. Conversely, every D-finite power series over $K$ has
such a name.
\end{theorem}
\begin{proof}
Compute $B$ by \cref{onot:hol:thm:integerroots}; there are finitely many
compatibility equations to check. For $n>B$, $P_0(n)\ne0$, so
\begin{equation}\label{onot:hol:eq:extendjet}
 a_n=-P_0(n)^{-1}\sum_{j=1}^JP_j(n-j)a_{n-j}
\end{equation}
is a uniquely determined finite computation from earlier coefficients. It gives
a formal series satisfying every coefficient equation, hence $\mathsf Lf=0$, and
any solution with that jet satisfies the same recursion, which proves
uniqueness. If the jet is zero, induction in \eqref{onot:hol:eq:extendjet} gives
the zero series; the converse of \eqref{onot:hol:eq:jetzero} is immediate. If
$B=-1$ the induction starts at $n=0$ and the only solution is zero. A D-finite
series satisfies a nonzero linear differential equation over $K(t)$; normalize
it as above and supply the jet of that series through the computed bound.
\end{proof}

``Every D-finite series'' refers to series with coefficients in $K$, whose
finitely many initial coefficients have $K$-names. No algorithm is asserted for
recognizing that an arbitrary coefficient program is D-finite, or for finding
its differential equation (08).

\begin{example}[The order is not a jet bound \src{08}]\label{onot:hol:ex:resonance}
For an ordinary integer $M>0$, $\mathsf L=\theta(\theta-M)$ has the solutions $1$
and $t^M$. Its order is two, but $t^M$ agrees with zero through index $M-1$; the
bound from $P_0(X)=X(X-M)$ is $B=M$, and a name for $t^M$ must include its
exceptional coefficient at $M$. No valid argument replaces the computed
singular-index bound by the operator order.
\end{example}

\subsection{Closure algorithms, not just closure assertions}

Let $\DF_K[[t]]$ be the ring of D-finite series over $K$. Its familiar closure
theorem \cite{Stanley} has the following effective form.

\begin{theorem}[Effective addition and multiplication \src{08}]\label{onot:hol:thm:dfclosure}
From certified names for $f,g\in\DF_K[[t]]$ one can compute certified names for
$f+g$, $-f$ and $fg$. No oracle for algebraic independence or for minimal
differential equations is needed.
\end{theorem}
\begin{proof}
A nonzero series satisfying an operator of order $r\ge1$ has a derivative vector
$v_f=(f,f',\ldots,f^{(r-1)})^{\mathsf T}$ with $v_f'=A_fv_f$, where $A_f$ is the
companion matrix over $K(t)$; the zero case is decidable and treated separately.
For a sum use the direct sum $A_f\oplus A_g$ with an output row selecting $f+g$;
for a product use $A=A_f\otimes I+I\otimes A_g$ on $v_f\otimes v_g$. These
systems have dimensions $r+s$ and $rs$. If $h=c_0v$ in a system $v'=Av$, define
$c_{i+1}=c_i'+c_iA$, so $h^{(i)}=c_iv$. In dimension $q$ the rows
$c_0,\ldots,c_q$ are linearly dependent over $K(t)$, and Gaussian elimination
finds a nontrivial relation $\sum_ib_i(t)c_i=0$, giving the nonzero annihilator
$\sum_ib_i(t)(d/dt)^i$ of $h$. All rational-function operations and zero tests
reduce to operations in $K$. Normalize the operator, compute its bound $B$, and
compute the jet by finite coefficient addition or Cauchy convolution; the genuine
sum or product is a solution, so the jet is compatible, and
\cref{onot:hol:thm:certificate} verifies the name. Negation keeps the operator
and negates the jet.
\end{proof}

Repeated operator construction can be expensive and can produce nonminimal
orders; that affects performance, not termination. Two different equations for
the same series are compared by computing an annihilator and the finite jet of
their difference, not by comparing the operators literally (08).

\subsection{Passing to a field}

D-finite series do not in general form a field, so we take
\begin{equation}\label{onot:hol:eq:H}
 H(K,t)=\Frac(\DF_K[[t]])\subseteq K((t)).
\end{equation}
A name is a pair $(p,q)$ of certified power-series names with $q\ne0$, a
decidable condition by \cref{onot:hol:thm:certificate}.

\begin{proposition}[Effective fraction field \src{08}]\label{onot:hol:prop:fractions}
$H(K,t)$ has computable arithmetic and equality; its Laurent valuation, leading
coefficient in $K$ and individual Laurent coefficients are computable.
\end{proposition}
\begin{proof}
Use the fraction formulas and \cref{onot:hol:thm:dfclosure}; equality is
$p/q=r/s\iff ps-rq=0$. For a nonzero fraction write $p=t^a(p_a+p_{a+1}t+\cdots)$
and $q=t^b(q_b+q_{b+1}t+\cdots)$ with $p_aq_b\ne0$, reading $a,b$ from the
certified jets; the valuation is $a-b$ and the leading coefficient $p_a/q_b$. A
coefficient at any specified Laurent index follows by triangular series division
after removing $t^b$, using finitely many earlier coefficients and division by
$q_b\ne0$.
\end{proof}
A reciprocal in $H(K,t)$ is represented as a fraction; it is not asserted to be
D-finite. Rational functions are included, and every D-finite Laurent series
belongs to the field after multiplication by a power of $t$.

\section{Finite-rank holonomic towers inside the surreal field}\label{onot:hol:sec:towers}

\subsection{Scales and their order}

Choose a finite increasing list of ordinal labels $\alpha_1<\cdots<\alpha_r$ and
set
\begin{equation}\label{onot:hol:eq:scales}
 \delta_j=\omega^{\alpha_j},\qquad t_j=\omega^{-\delta_j},\qquad U_j=t_j^{-1}=\omega^{\delta_j},\qquad G_j=\Z\delta_1+\cdots+\Z\delta_j.
\end{equation}
The coefficients in $G_j$ are ordinary integers and addition is surreal
addition. The $\delta_j$ are linearly independent over $\Q$ by normal-form
uniqueness, and their integer span is ordered lexicographically from the
\emph{largest} index: the sign of a nonzero vector is the sign of its last
nonzero coordinate, so every positive multiple of $\delta_j$ dominates $G_{j-1}$.
The labels are not exponents of a formal analytic exponential; they specify
actual Conway monomials and so a definite embedding of the exponent lattice.
Different label choices can have the same finite-rank algorithmic structure but
very different magnitudes. Define
\begin{equation}\label{onot:hol:eq:tower}
 F_0=k,\qquad F_j=H(F_{j-1},t_j)\quad(1\le j\le r).
\end{equation}
At level $j$ the field $F_{j-1}$ is constant for differentiation in $t_j$. The
recursion on levels is essential: an operator may not use its own unknown
denotation among its coefficients.

\begin{lemma}[Iterated Laurent embedding \src{08}]\label{onot:hol:lem:embedding}
There is an injective ordered-field embedding
$k((t_1))\cdots((t_r))\to\No$ with the substitutions \eqref{onot:hol:eq:scales}.
For $f=\sum_{n\ge N}a_nt_r^n\ne0$ with $a_N\ne0$,
\begin{equation}\label{onot:hol:eq:leadingtower}
 \ddeg f=\ddeg a_N-N\delta_r,\qquad \lc(f)=\lc(a_N).
\end{equation}
\end{lemma}
\begin{proof}
Induct on $r$. Every exponent of $a_nt_r^n$ has the form $g-n\delta_r$ with
$g\in G_{r-1}$. If $n<m$, every exponent of the $n$th block exceeds every
exponent of the $m$th, because $(m-n)\delta_r$ dominates $G_{r-1}$. Each block is
reverse well ordered by induction and the blocks are indexed by a set of integers
bounded below, so the ordered union is reverse well ordered; distinct integer
vectors give distinct exponents, so blocks do not overlap. In a product, a
prescribed top coordinate receives finitely many pairs of top indices, and within
each pair the lower-rank convolution is legitimate by induction. Nonzero series
have the leading term \eqref{onot:hol:eq:leadingtower}, which gives injectivity
and order preservation, and a field homomorphism respects inverses. Rank zero is
$k\subseteq\R\subseteq\No$.
\end{proof}

\begin{theorem}[Effective holonomic tower \src{08}]\label{onot:hol:thm:tower}
For every finite list of labels, \eqref{onot:hol:eq:tower} gives a countable
subfield $F_r\subseteq\No$ with a decidable finite-name representation, and
computable field operations, equality, order, leading exponent, real leading
coefficient and coefficient extraction at each specified exponent of $G_r$.
\end{theorem}
\begin{proof}
Rank zero is the contract for $k$. Suppose the assertions hold at rank $j-1$.
\Cref{onot:hol:thm:integerroots} computes the exceptional indices of every Euler
operator over $F_{j-1}$, so certificates over that field have decidable validity
and zero tests, and \cref{onot:hol:thm:dfclosure,onot:hol:prop:fractions} give
the field algorithms at rank $j$. \Cref{onot:hol:lem:embedding} supplies the
denotation and computes the leading exponent and real leading coefficient by
descending one rank; the sign is that of the latter. To extract the coefficient
at $g=m\delta_j+h$, take the Laurent coefficient at index $-m$ and then its
coefficient at $h\in G_{j-1}$. Each algorithm calls lower-rank algorithms and
finitely many same-rank operations, never an unconstructed next rank. The codes
are finite trees, so there are countably many.
\end{proof}

There is no claim of analytic convergence: a formal factorial series is as
legitimate as an exponential series. Nor is real closedness asserted: the tower
is a specific effective subfield, not an algebraic closure or the field of all
surreals with a given support group (08).

\subsection{A modular generalization}

The role of differential equations can be isolated. Suppose a class
$\mathfrak C(K,t)\subseteq K[[t]]$ has decidable finite names, computable ring
operations and coefficients, and a computable finite zero-test bound for each
name, and contains $K[t]$. Then the same fraction-field and finite-rank
construction applies, provided the resulting fields carry the leading data needed
at the next stage; certified D-finite series are one completely specified
realization. This permits later replacement by stronger certified classes without
changing the omnific arguments below. It does \emph{not} justify inserting
arbitrary differential-equation solvers: a zero test, a chosen-solution
convention and terminating coefficient algorithms remain hypotheses (08). This is
stated as an observation, not as a theorem with proof.

\section{Truncation, triangular integer forms, and floor}\label{onot:hol:sec:integerpart}

\subsection{Finite descriptions of infinite truncations}

For $g\in G_r$, $\tr g(f)$ is the sum of the terms of $f$ with exponent at least
$g$, and $\trs g(f)$ the sum of those above $g$; a truncation need not have
finite support.

\begin{theorem}[Effective truncation closure \src{08}]\label{onot:hol:thm:truncation}
$F_r$ is effectively closed under $\tr g$ and $\trs g$ for every $g\in G_r$. In
particular $f_{>0}$, $\ct(f)$ and $f_{<0}$ are computable as finite names.
\end{theorem}
\begin{proof}
Rank zero is immediate. At rank $r$ take $f=\sum_{n\ge N}a_nt_r^n$ and
$g=m\delta_r+h$ with $h\in G_{r-1}$; for $f=0$ return zero. Block dominance gives
\begin{equation}\label{onot:hol:eq:truncation}
 \tr g(f)=\sum_{N\le n<-m}a_nt_r^n+t_r^{-m}\tr h(a_{-m}),
\end{equation}
with an empty sum zero and coefficients below $N$ zero; the same formula with
$\trs h$ in the last term gives $\trs g(f)$. The first sum is finite, each
coefficient has an effective $F_{r-1}$-name, and the last term is available by
induction. Constant extraction is the iterated coefficient at zero, and
subtraction gives the negative part.
\end{proof}

An exponent cutoff and a number-of-terms cutoff are different interfaces:
\eqref{onot:hol:eq:truncation} returns a finite \emph{expression} for the exact
truncation, whose support can have order type $\omega^{r-1}$ or more; it does not
print an infinite list as a finite one (08).

\subsection{The exact omnific part}

Put $R_r=F_r\cap\Oz$, with $R_0=\Z\subseteq k$.

\begin{theorem}[Triangular omnific normal form \src{08}]\label{onot:hol:thm:triangular}
For each $r\ge1$,
\begin{equation}\label{onot:hol:eq:Rrecursion}
 R_r=R_{r-1}+U_rF_{r-1}[U_r].
\end{equation}
Every element of $R_r$ has a unique finite decomposition by coefficient values
\begin{equation}\label{onot:hol:eq:triangular}
 z=m+\sum_{j=1}^r\sum_{i=1}^{d_j}U_j^ia_{j,i},\qquad m\in\Z,\ a_{j,i}\in F_{j-1},
\end{equation}
after deleting trailing zero coefficients at each level. Recognition of $R_r$
inside $F_r$ and computation of this decomposition are decidable, and
arithmetic, equality and order in this notation are computable.
\end{theorem}
\begin{proof}
Write $f=\sum_{n\ge N}a_nt_r^n$. All exponents of a nonzero block with $n<0$ are
positive, since $-n\delta_r$ dominates every lower-rank exponent; all exponents of
a nonzero block with $n>0$ are negative; the block $n=0$ has exactly its
lower-rank normal form. So $f$ is omnific if and only if the positive-$t_r$ tail
is zero and $a_0\in R_{r-1}$. There are finitely many negative Laurent indices,
with arbitrary coefficients in $F_{r-1}$. This proves \eqref{onot:hol:eq:Rrecursion}
inductively.

For the algorithm, extract the finite nonpositive Laurent prefix and compare it
with $f$ by the field equality test; if they differ, the positive tail is nonzero
and $f\notin R_r$. Otherwise apply the lower-rank recognition to $a_0$. At rank
zero, \cref{onot:hol:lem:archsearch} computes $\floor c$ and decides whether it
equals $c$, which is membership in $\Z$. Different top indices have disjoint
support blocks, so equality of two such expressions forces equality of every top
coefficient and of the constant block; descending the ranks proves uniqueness.
The field algorithms decide arithmetic and order, and results are normalized
again. Closure also follows from \eqref{onot:hol:eq:Rrecursion}: products of
positive-degree polynomials keep positive top degree, and $R_{r-1}$ is a subring
of the coefficient field.
\end{proof}

The triangular form is a finite outer normal form with exact field coefficients,
not a claim that every coefficient has a finite Conway normal form. Applying
\cref{onot:prop:canonical} to the coefficient names would give a literal
canonical string, which is not needed for the equality test (08). In the language
of \cref{onot:sec:shield}, \eqref{onot:hol:eq:Rrecursion} says that the omnific
part of $F_r$ is exactly the shielded ring
$\mathcal S_{\delta_r}(R_{r-1},F_{r-1})$ (\cref{onot:rem:characterizations}).

\begin{corollary}[Normal-form membership test \src{08}]\label{onot:hol:cor:oztest}
For $x\in F_r$ let $p=x_{>0}$, $c=\ct(x)$ and $\epsilon=x_{<0}$. Then
$x\in R_r$ iff $\epsilon=0$ and $c\in\Z$. Both tests are effective, and
$p\in R_r$ always.
\end{corollary}

\begin{theorem}[Computable omnific floor \src{08}]\label{onot:hol:thm:floor}
For every $x\in F_r$ the omnific floor $\fl x$ lies in $R_r$ and is computable;
with the decomposition of \cref{onot:hol:cor:oztest} it is given by
\eqref{onot:eq:floor}.
\end{theorem}
\begin{proof}
By \cref{onot:hol:thm:truncation} the names of $p$, $c$ and $\epsilon$ are
computable; $\floor c$ and the test $c\in\Z$ are computable at rank zero by
\cref{onot:hol:lem:archsearch}, and the sign of $\epsilon$ by
\cref{onot:hol:thm:tower}. The right side of \eqref{onot:eq:floor} lies in $R_r$
by \cref{onot:hol:cor:oztest}, and it is the floor by \cref{onot:sub:floor}.
\end{proof}

\section{Support order types and the ordinal part}\label{onot:hol:sec:ordinals}

\subsection{Sharp finite-rank support bounds}

$\ord(\supp f)$ is the order type of the support read in decreasing exponent
order, zero for $f=0$. It is a property of the normal form, not the birthday.

\begin{theorem}[Support spectrum \src{08}]\label{onot:hol:thm:support}
For $r\ge1$,
\[
 f\in F_r\Rightarrow\ord(\supp f)\le\omega^r,\qquad z\in R_r\Rightarrow\ord(\supp z)<\omega^r.
\]
The field bound is attained by a rational expression, and for each
$\eta<\omega^r$ a rational expression in $R_r$ has support order type exactly
$\eta$.
\end{theorem}
\begin{proof}
At rank zero a nonzero element has one term. At rank $r$ an iterated Laurent
series has at most countably many top blocks, each of type at most
$\omega^{r-1}$, so the ordinal sum has type at most $\omega^r$. An omnific
integer has finitely many negative-$t_r$ blocks and then one $R_{r-1}$ block; for
$r=1$ its support is finite, and for $r>1$ its type is at most
$\omega^{r-1}d+\beta<\omega^r$ with $d<\omega$, $\beta<\omega^{r-1}$.

For sharpness, $\Geo_r=\prod_{j=1}^r(1-t_j)^{-1}$ has one nonzero coefficient at
every vector of $\N^r$, read in the iterated lexicographic order, so its support
has type $\omega^r$. For $\eta=\omega^{r-1}d_r+\cdots+d_1$ with $d_j\in\N$, the
element
\begin{equation}\label{onot:hol:eq:realizesupport}
 z_\eta=\sum_{j=1}^r\sum_{i=1}^{d_j}U_j^i\prod_{l<j}(1-t_l)^{-1}
\end{equation}
is omnific by \cref{onot:hol:thm:triangular}. Each $(j,i)$ block has support type
$\omega^{j-1}$, all coefficients are positive and the blocks are disjoint, so
nothing cancels; higher $j$ blocks precede lower ones, giving exactly the Cantor
normal form of $\eta$. For $\eta=0$ take zero.
\end{proof}

These bounds contrast with \cref{onot:neg:sec:raw}: arbitrarily large finite-rank
support types are compatible with decidable equality, while a raw name with at
most one monomial can have an undecidable zero test (08).

\subsection{Recognizing actual ordinals}

Three operations must not be conflated: surreal addition, ordinary ordinal
addition and natural (Hessenberg) addition. On ordinals the first agrees with the
third, not generally with the second: $1+\omega=\omega+1$ in $\No$, while
$1+_{\mathrm{ord}}\omega=\omega$. The constructors here are field operations;
ordinary ordinal operations can be performed on recognized Cantor normal forms
(08; also 05, \cref{onot:sh:thm:ordinalfragment}).

\begin{lemma}[Ordinal points of the exponent group \src{08}]\label{onot:hol:lem:ordinalgroup}
For the scales \eqref{onot:hol:eq:scales},
$G_r\cap\On=\N\delta_1+\cdots+\N\delta_r$.
\end{lemma}
\begin{proof}
A nonnegative integer combination is the Cantor normal form of an ordinal.
Conversely an element of $G_r$ has Conway normal form $\sum_jn_j\omega^{\alpha_j}$
with distinct ordinal exponents; an ordinal's Conway normal form is its Cantor
normal form, whose coefficients are positive integers, so every nonzero $n_j$ is
positive.
\end{proof}

\begin{theorem}[Ordinal slice and its recognition \src{08}]\label{onot:hol:thm:ordinalslice}
The ordinals in $F_r$, equivalently in $R_r$, are exactly
\begin{equation}\label{onot:hol:eq:ordinaltrace}
 F_r\cap\On=\N[U_1,\ldots,U_r].
\end{equation}
Given an $F_r$-name one can decide whether its value is an ordinal and, if so,
compute its Cantor normal form. For the labels $\alpha_j=j-1$ the slice is exactly
$\On_{<\omega^{\omega^r}}$.
\end{theorem}
\begin{proof}
An ordinal has a finite Cantor normal form with ordinal exponents and positive
integer coefficients; every exponent of an element of $F_r$ lies in $G_r$, so by
\cref{onot:hol:lem:ordinalgroup} it is a natural-coefficient polynomial in the
$U_j$, and every such polynomial is an ordinal. For recognition, extract the
finite Laurent prefix of indices at most zero and test that the positive-index
tail vanishes; if not, some exponent has negative top coordinate and is not an
ordinal exponent. Otherwise recursively test each coefficient of the finite
polynomial in $U_r$ for membership in $\N[U_1,\ldots,U_{r-1}]$; at rank zero test
membership in $\N$. The recursion terminates and yields the Cantor normal form,
since $\prod_jU_j^{n_j}$ has exponent $\sum_jn_j\omega^{\alpha_j}$ and these
exponents can be sorted exactly. For consecutive labels the exponent-group
ordinals are those below $\omega^r$, and finite Cantor normal forms with such
exponents give exactly the ordinals below $\omega^{\omega^r}$.
\end{proof}

The recognition algorithm does not decide whether all terms of an arbitrary
recurrence are nonnegative; it uses a zero test for the positive-index tail and
finitely many lower-rank tests, so no unproved positivity procedure is hidden
(08).

\section{One notation system for all ordinals below epsilon-zero}\label{onot:hol:sec:epsilon}

\subsection{Canonical labels}

Use hereditary Cantor-normal-form trees: zero is a tree, and a nonzero tree is a
finite list $(\beta_1,n_1),\ldots,(\beta_s,n_s)$ with trees
$\beta_1>\cdots>\beta_s$ and $n_i\in\N_{>0}$, interpreted as
$\omega^{\beta_1}n_1+\cdots+\omega^{\beta_s}n_s$. Comparison is lexicographic,
recursing into the exponents; finite-tree induction proves termination and
decidable validity. The classical Cantor-normal-form theorem identifies these
trees with $\On_{<\varepsilon_0}$ (see also the formal treatment \cite{Bancerek}),
$\varepsilon_0$ being the least fixed point of $\beta\mapsto\omega^\beta$. For a
finite set $A=\{\alpha_1<\cdots<\alpha_r\}$ of labels write $\mathcal F_A$,
$\mathcal R_A$ for the fields and rings above; a global name is a finite label set
with a valid local name, and empty label sets give $k$ and $\Z$.

\begin{lemma}[Effective insertion of scales \src{08}]\label{onot:hol:lem:insert}
If $A\subseteq A'$ there are computable denotation-preserving embeddings of local
names for $\mathcal F_A$ into those for $\mathcal F_{A'}$, preserving the omnific
parts; the semantic embeddings commute under successive insertions.
\end{lemma}
\begin{proof}
Insert missing labels in increasing order. At an existing variable, map all
lower-level operator coefficients and jet entries through the embedding already
constructed; a nonzero polynomial stays nonzero and its ordinary integer roots
are unchanged under an injective field embedding, so the exceptional bound and
compatibility equations are preserved, and the induced series has the mapped
coefficients by the uniqueness recursion. Fractions map likewise. At an inserted
variable, regard every old lower-field element $c$ as the constant series with
certificate $\theta f=0$ and jet $f(0)=c$. \Cref{onot:hol:lem:embedding}
identifies both expressions with the same surreal under the fixed scales, so
different insertion orders agree and omnific membership is preserved.
\end{proof}

\begin{theorem}[Epsilon-labelled decidable omnific notation \src{08}]\label{onot:hol:thm:epsilon}
The directed unions
\[
 \mathcal F_\varepsilon=\bigcup_{A\subseteq_{\mathrm{fin}}\varepsilon_0}\mathcal F_A,\qquad
 \mathcal R_\varepsilon=\bigcup_{A\subseteq_{\mathrm{fin}}\varepsilon_0}\mathcal R_A
\]
have finite-name representations with decidable validity, equality and order,
computable arithmetic, computable omnific recognition and floor, and computable
truncation at every finitely named exponent of their exponent groups. Moreover
\begin{equation}\label{onot:hol:eq:epsilontrace}
 \mathcal F_\varepsilon\cap\On=\mathcal R_\varepsilon\cap\On=\On_{<\varepsilon_0}.
\end{equation}
Recognition and conversion of this slice are effective, and the whole omnific
system has a computable canonicalization.
\end{theorem}
\begin{proof}
Given two names, embed both in the field of the sorted union of their label sets
(\cref{onot:hol:lem:insert}) and apply the local algorithms there; a cutoff
exponent is handled by the same union. Validity is finite label comparison plus
the finite certificate checks. Local floor and recognition commute with
insertion, so the answers do not depend on the common field.

If an element of the union is an ordinal, \cref{onot:hol:thm:ordinalslice}
writes it as a natural-coefficient polynomial in finitely many
$\omega^{\omega^\alpha}$ with $\alpha<\varepsilon_0$; every exponent of its Cantor
form is a finite natural sum of such $\omega^\alpha$ and is below $\varepsilon_0$,
so the ordinal is below $\varepsilon_0$. Conversely, for $\xi<\varepsilon_0$
collect the exponents of the Cantor normal forms of the exponents of $\xi$ as
labels $A$; each $\omega^\beta$ is then a monomial in the
$U_\alpha=\omega^{\omega^\alpha}$, $\alpha\in A$, so $\xi\in\mathcal R_A$, and the
local recognition converts between the forms. Finally apply
\cref{onot:prop:canonical} to the total system of all locally valid names.
\end{proof}

``Total notation system'' uses the decidable set of locally valid codes as its
syntactic domain; malformed strings and incompatible certificates are rejected
first, and no noncomputable choice of representative is hidden (08).

\begin{corollary}[Cofinality and an ambient ordinal bound \src{08}]\label{onot:hol:cor:ambientbound}
For a nonempty finite label set $A$ with greatest label $\alpha$, the positive
powers of $U_\alpha=\omega^{\omega^\alpha}$ are cofinal in $\mathcal F_A$, and its
least ordinal upper bound is $\omega^{\omega^{\alpha+1}}$. For the union,
$\mathcal F_\varepsilon\subset(-\varepsilon_0,\varepsilon_0)$, and $\varepsilon_0$
is its least ordinal upper bound. These are ordinal upper bounds, not suprema in
$\No$.
\end{corollary}
\begin{proof}
For nonzero $x\in\mathcal F_A$ the leading exponent lies in the span of the
scales; choose $N$ with $\ddeg x<N\omega^\alpha$, so $|x|<U_\alpha^N$, an ordinal
of $\mathcal F_A$. In the ordinal order,
$\sup_N\omega^{\omega^\alpha\cdot N}=\omega^{\omega^\alpha\cdot\omega}=\omega^{\omega^{\alpha+1}}$
(ordinary ordinal products). For the empty label set the naturals are cofinal in
$k$ and the bound is $\omega$. For $\alpha<\varepsilon_0$ the local bound is below
$\varepsilon_0$, so every element of the union has absolute value below
$\varepsilon_0$, and \eqref{onot:hol:eq:epsilontrace} shows no smaller ordinal
bounds the union. Indeed $\varepsilon_0/2$ is already a smaller surreal upper
bound: every element is below some ordinal $\beta<\varepsilon_0$, and $\beta+\beta$
is still an ordinal below $\varepsilon_0$, so $\beta<\varepsilon_0/2$. This
prevents reading ``least ordinal upper bound'' as ``surreal supremum''.
\end{proof}

\begin{remark}[Epsilon-zero is a design choice \src{08}]\label{onot:hol:rem:design}
The theorem does not make $\varepsilon_0$ an absolute ceiling for effective
omnific notation. The labels are canonical hereditary Cantor trees; a different,
independently justified decidable ordinal-label system can be substituted in the
finite-rank construction, but its validity and comparison properties must be
proved separately, and an unrestricted well-foundedness promise cannot be
smuggled in as decidable syntax (08). The union names $\varepsilon_0$ in no form,
which is compatible with 05's use of $\varepsilon_0$ as a primitive symbol
(\cref{onot:sh:thm:epsilon}), because the languages differ; compare the remark
after \cref{onot:prop:size} (\mergetag).
\end{remark}

\begin{example}[Large magnitudes at small rank \src{08}]\label{onot:hol:ex:magnitude}
The labels $A=\{0,\omega\}$ give the ordinal $\omega^{\omega^\omega}+\omega^3+2$ at
rank two, far beyond the consecutive-label cutoff $\omega^{\omega^2}$, with two
nested coefficient fields. Label magnitude, finite rank, support order type and
birthday are four different measurements.
\end{example}

\section{Worked infinite-support notations}\label{onot:hol:sec:examples}

Take $k=\Q$, consecutive labels $0,1$, and abbreviate $t=t_1=\omega^{-1}$ and
$U=U_2=\omega^\omega$. For every $f\in F_1$, $Uf$ is omnific, because the exponents
of $f$ are ordinary integers bounded above and $\omega$ dominates them; so
$\Z+UF_1[U]\subseteq R_2$. This restricted ring is what the delivered prototype
implements (\cref{onot:sec:implementation}).

\paragraph{A factorial tail.} $\Fac(t)=\sum_{n\ge0}n!\,t^n$ has the certificate
$\mathsf L=\theta-t(\theta+1)^2$ with $a_0=1$: $P_0(X)=X$ has the sole integer
root $0$ and \eqref{onot:hol:eq:recurrence} reads $na_n-n^2a_{n-1}=0$ for
$n\ge1$, so $a_n=n!$. Thus
\[
 U\Fac(t)=\omega^\omega+\omega^{\omega-1}+2\omega^{\omega-2}+6\omega^{\omega-3}+\cdots
\]
is a finitely named omnific integer of infinite support. That $\Fac$ has no
positive radius of convergence is irrelevant to its formal or surreal
denotation.

\paragraph{Exponential notation without a global exponential.}
$\Ex(t)=\sum_nt^n/n!$ and $\Ex(-t)$ have the certificates $\theta-t$ and
$\theta+t$ with constant coefficient one. The formal product has derivative zero
and constant term one, and in characteristic zero a series with zero derivative
is constant, so
\begin{equation}\label{onot:hol:eq:expidentity}
 (U\Ex(t))(U\Ex(-t))=U^2.
\end{equation}
Both factors have infinite normal forms and are omnific. The identity is proved
in formal series and transported to $\No$; it needs no choice of a global surreal
exponential. The closure algorithm independently produces a certificate for the
difference and checks its finite jet, without recognizing the identity from the
printed expressions.

\paragraph{Cancellations and leading terms.} $\Geo(t)=(1-t)^{-1}$ has the
certificate $\theta-t(\theta+1)$ with initial coefficient one. The first two
coefficients of $\Geo-\Ex$ vanish and the next is $\tfrac12$, so
\begin{equation}\label{onot:hol:eq:leadingexample}
 \lt\bigl(U(\Geo(t)-\Ex(t))\bigr)=\tfrac12\omega^{\omega-2},
\end{equation}
found after a finite certified zero test, not a fixed truncation length. At higher
rank a whole lower-rank infinite series can cancel in a single top coefficient;
the recursive equality test lets the leading-term algorithm move past it, which
enumerating the expanded normal form would not guarantee. Note that
$U\Geo(t)=\sB$.

\paragraph{Ordinal versus nonordinal outputs.}
$U(t^{-2}+3)+5t^{-2}+7=\omega^{\omega+2}+3\omega^\omega+5\omega^2+7$ is an
ordinal, and recognition outputs this Cantor normal form. $U\Fac(t)$ is not an
ordinal: it has infinitely many terms, and $\omega-n$ for $n>0$ is not an ordinal
although it is a positive exponent; integer coefficients do not remedy either
issue. Similarly $\sqrt2\,U$ is omnific when $\sqrt2\in k$ but is not an ordinal;
confusing arbitrary real coefficients with Cantor coefficients would exclude much
of the intended domain.

\part{Hereditary forms, the epsilon-zero-shielded calculus, and finite-state series}\label{onot:sh:part}

This part is source 05's. Its coefficient field is an exact effective ordered
subfield $k\subseteq\R$ with canonical codes, effective field operations and
decidable comparison; the implementation uses $k=\Q$ and the mathematics also
applies to $k=\Ralg$ (\cref{onot:sub:decidable}). No arbitrary computable-real
coefficient oracle is assumed.

\section{A hereditary finite normal-form calculus}\label{onot:sh:sec:hereditary}

\subsection{Syntax and denotation}

A hereditary form is a finite tree of the shape
\begin{equation}\label{onot:sh:eq:hereditary}
 a_0+\sum_{i=1}^ma_i\omega^{e_i},
\end{equation}
where $a_0\in k$, $a_i\in k\setminus\{0\}$, and the $e_i$ are nonzero hereditary
forms of strictly smaller tree height in strictly decreasing numerical order; at
height zero there are only constants. Let $\Hh=\Hh_k$ be the class of values of
such trees. The exponent zero is absorbed into $a_0$. Negative exponents are
allowed in this auxiliary calculus; omnific membership is imposed afterwards
rather than making exponent arithmetic artificially one-sided. A convenient raw
expression language has constants from $k$ and constructors $+$, $-$, $\cdot$
and $\Om(e)$ with denotation $\omega^{\nu(e)}$; its normal forms are exactly the
hereditary trees. It contains, for example, $(\omega+1)^2=\omega^2+2\omega+1$ and
$\omega^{\omega-1}+3\omega^{1/2}-7$, and over algebraic coefficients
$\omega^{\sqrt2\,\omega+1/3}$. It has no division and no limits as raw
constructors.

\begin{theorem}[Hereditary normalization \src{05}]\label{onot:sh:thm:hereditary}
There are terminating exact algorithms for normalization, equality, comparison,
addition, negation, multiplication and the omega map on hereditary expressions.
Two canonical trees denote the same surreal exactly when their coefficient and
exponent trees are identical. In particular $\Hh$ is an ordered subring of
$\No$.
\end{theorem}
\begin{proof}
Proceed by tree height. At height zero use the field algorithms. For normalized
exponents of lower height, comparison is available by induction: merge equal
exponents, add their coefficients, remove zero coefficients, absorb exponent zero
into the constant and sort the rest. For addition merge the two finite lists;
for multiplication form all pairs of terms, using
$(a\omega^e)(b\omega^f)=ab\,\omega^{e+f}$, where the exponent addition takes place
at strictly smaller height, and normalize. Negation changes signs. To normalize
$\Om(e)$, normalize $e$ and output $1$ if $e=0$, and otherwise the single monomial
with coefficient one and exponent $e$. A raw expression is normalized along its
finite syntax tree. To compare normal forms, inspect the greatest exponent at
which the lists differ; the sign is that of the coefficient difference, and
comparing candidate exponents uses only smaller heights. The procedures terminate
by induction on height. Normal-form uniqueness says equal values have the same
exponents and coefficients, and the induction hypothesis turns equality of
exponent values into identity of their canonical trees. Hahn arithmetic gives the
ring identities.
\end{proof}

The proof is a decidability argument, not a polynomial bit-complexity bound: a
short product can have many monomials, and exponent trees may contain large
expressions. Hash-consing and directed acyclic graphs are useful implementation
choices but do not change the semantics (05).

\subsection{Omnific membership and the floor}

Put $\Dd=\Hh\cap\Oz$, and assume also that the integer test and the ordinary
floor are effective in $k$, as they are for $\Q$ and $\Ralg$.

\begin{theorem}[Decidable omnific core and floor \src{05}]\label{onot:sh:thm:floor}
Membership in $\Dd$ is decidable; $\Dd$ has decidable equality and order; and
the omnific floor restricts to an effective map $\fl{\cdot}:\Hh\to\Dd$ with
$\fl x\le x<\fl x+1$.
\end{theorem}
\begin{proof}
In the normal form of $x$, test whether every nonconstant exponent is positive
and whether the constant is in $\Z$; criterion \eqref{onot:eq:oz} gives
membership, and closure under ring operations follows from positive supports or
from the ring property of $\Oz$. The decomposition
\eqref{onot:eq:threeparts} of a hereditary form is read off its finite normal
form: $x_{>0}$ is the finite sum of positive-exponent terms, $\ct(x)\in k$, and
$x_{<0}$ is the finite sum of negative-exponent terms, whose sign is the sign of
its leading coefficient. Formula \eqref{onot:eq:floor} therefore produces a
hereditary form with positive exponents and integer constant, which lies in $\Dd$
and is the floor by \cref{onot:sub:floor}.
\end{proof}

\subsection{The ordinal fragment}

\begin{theorem}[Ordinal fragment \src{05}]\label{onot:sh:thm:ordinalfragment}
For $k=\Q$ or $k=\Ralg$, the ordinals in $\Hh$ are exactly the ordinals below
$\varepsilon_0$. They are recognized recursively: the normal form must have
natural-number coefficients and hereditarily ordinal exponents. Arithmetic
inherited from $\No$ restricts to natural (Hessenberg) addition and
multiplication, not ordinary ordinal arithmetic.
\end{theorem}
\begin{proof}
The normal form of an ordinal is its Cantor normal form: positive natural
coefficients, ordinal exponents, natural constant; conversely such a finite form
is an ordinal. Apply the criterion recursively to the exponent subtrees. Finite
hereditary Cantor normal forms are exactly the ordinals below $\varepsilon_0$:
downward recursion on leading exponents constructs each such ordinal, and
induction on tree height bounds each constructed ordinal below a finite tower of
omega powers, whose supremum in the ordinal order is $\varepsilon_0$. The
coefficientwise sum and distributive product of Cantor forms are the natural sum
and product, which are the operations the normalizer uses.
\end{proof}

A program supporting both surreal and ordinary ordinal operations must give them
different names or types (05, 08).

\begin{lemma}[Hereditary values have smaller ordinal bounds \src{05}]\label{onot:sh:lem:epsilonbound}
For every $h\in\Hh_k$, $k\subseteq\R$, there is an ordinal $\alpha<\varepsilon_0$
with $|h|<\alpha$.
\end{lemma}
\begin{proof}
For a constant take a large finite ordinal. For a normal form, induction gives a
common positive ordinal $\beta<\varepsilon_0$ exceeding the absolute values of all
exponent subtrees. Each monomial is smaller than $\omega^\beta$ up to a finite
real factor, so a finite sum is bounded by $\omega^{\beta+1}$; increasing once
more gives a strict bound, still below $\varepsilon_0$.
\end{proof}

Although $\varepsilon_0$ is the least \emph{ordinal} bounding all hereditary
ordinals, it is not their least surreal upper bound: $\varepsilon_0-1$ is also an
upper bound of every ordinal below $\varepsilon_0$. No surreal supremum assertion
is made (05; compare \cref{onot:hol:cor:ambientbound}, where 08 exhibits
$\varepsilon_0/2$).

\section{The epsilon-zero-shielded calculus}\label{onot:sh:sec:epsilon}

Let $k=\Q$ or $\Ralg$, and consider $\FH$ and $\Dd=\Hh\cap\Oz$. A valid
$\FH$-name is a pair $(p,q)$ of hereditary forms with $q\ne0$, denoting $p/q$;
equality is decided by
\begin{equation}\label{onot:sh:eq:crossmultiply}
 \frac pq=\frac rs\iff ps-rq=0
\end{equation}
using \cref{onot:sh:thm:hereditary}, and signs come from the signs of numerator
and denominator. Division introduces infinite normal forms without destroying
exact equality.

\begin{lemma}[Support of hereditary fractions \src{05}]\label{onot:sh:lem:fractionsupport}
Every member of $\FH$ has a Hahn normal form with coefficients in $k$ and
exponents in the additive group $\Hh$.
\end{lemma}
\begin{proof}
This is immediate for hereditary forms. For $q\ne0$ factor its leading term,
$q=c\omega^a(1+u)$ with $a\in\Hh$, $c\in k^\times$ and $u$ a finite sum with
strictly negative exponents in $\Hh$. The geometric inverse
$q^{-1}=c^{-1}\omega^{-a}\sum_{n\ge0}(-u)^n$ is Hahn summable by the
positive-support lemma, and its exponents are finite sums of elements of $\Hh$
shifted by $-a$. Multiplication by $p$ keeps both properties.
\end{proof}

\begin{theorem}[The $\varepsilon_0$-shielded omnific calculus \src{05}]\label{onot:sh:thm:epsilon}
Let $\varepsilon_0$ be embedded as a surreal, so that $\varepsilon_0=\omega^{\varepsilon_0}$. Then
\begin{equation}\label{onot:sh:eq:calculus}
 \SE=\Bigl\{d+\sum_{j=1}^m\varepsilon_0^{\,j}f_j:\ d\in\Dd,\ f_j\in\FH\Bigr\}
\end{equation}
is a subring of $\Oz$ with decidable validity, equality and order on finite
polynomial--fraction names. It contains $\Dd$, $\varepsilon_0$ is transcendental
over $\FH$, and $\SE$ contains elements with infinite normal forms.
\end{theorem}
\begin{proof}
Take $G=\Hh$. \Cref{onot:sh:lem:epsilonbound} gives $|g|<\varepsilon_0$ for all
$g\in G$, and \cref{onot:sh:lem:fractionsupport} puts $\FH$ in the corresponding
Hahn field. Apply \cref{onot:thm:shield} with $\Lambda=\varepsilon_0$ and
$X=\omega^{\varepsilon_0}=\varepsilon_0$; thus $\SE=\mathcal
S_{\varepsilon_0}(\Dd,\FH)$. The denominator and omnific-base tests are
decidable, as are coefficient arithmetic and equality, so polynomial
normalization and the sign rule give the algorithms. Injectivity of evaluation
is transcendence of $\varepsilon_0$ over $\FH$. For an infinite support,
$\varepsilon_0/(1-\omega^{-1})=\sum_{n\ge0}\omega^{\varepsilon_0-n}$ has positive
exponents and all coefficients one.
\end{proof}

The polynomial representation is canonical only after coefficient values have
been identified; fractions need not be reduced to a unique quotient, since cross
multiplication already decides equality. An inefficient least-code
canonicalizer is available (\cref{onot:prop:canonical}), but no multivariate gcd
algorithm is hidden in the procedure (05).

\begin{proposition}[An algebraic limit of the fraction core \src{05}]\label{onot:sh:prop:notrealclosed}
$1+\omega$ has no square root in $\FH$. In particular $\FH$ is not real closed,
even for $k=\Ralg$.
\end{proposition}
\begin{proof}
Suppose $p/q$ were a square root with $p,q\in\Hh$, $q\ne0$. Take the
finite-dimensional $\Q$-vector space spanned by $1$ and the outer exponents of
$p$ and $q$, with a basis $\beta_1=1,\beta_2,\ldots,\beta_r$; clearing the
rational coordinates gives $N>0$ such that $p,q$ are Laurent polynomials over $k$
in $x_i=\omega^{\beta_i/N}$. The $x_i$ are algebraically independent over $k$:
distinct integer exponent vectors give distinct exponents by $\Q$-independence,
and normal-form uniqueness prevents a relation. So $1+x_1^N$ would be a square in
$k(x_1,\ldots,x_r)$. But $1+x_1^N$ is nonconstant and squarefree in
characteristic zero (it is coprime to $Nx_1^{N-1}$), and in the factorial ring
$k[x_1,\ldots,x_r]$ a square in the fraction field has even order at every
irreducible factor.
\end{proof}

\begin{corollary}[One-state supports of type $\omega^r$ \src{05}]\label{onot:sh:cor:omegar}
For every positive integer $r$,
\begin{equation}\label{onot:sh:eq:omegar}
 \zeta_r=\varepsilon_0\prod_{j=0}^{r-1}\bigl(1-\omega^{-\omega^j}\bigr)^{-1}
\end{equation}
lies in $\SE$. Its normal form has all coefficients one and support
$\{\varepsilon_0-\sum_{j<r}n_j\omega^j:(n_0,\ldots,n_{r-1})\in\N^r\}$, of
decreasing order type $\omega^r$.
\end{corollary}
\begin{proof}
Each denominator is a nonzero hereditary form, so the product of reciprocals lies
in $\FH$, and the expansion of the geometric factors is jointly Hahn summable. The
monomials $1,\omega,\ldots,\omega^{r-1}$ are $\Q$-independent, so distinct tuples
give distinct exponents; their nonnegative integer combinations in increasing
order are the ordinals below $\omega^r$, and subtraction from $\varepsilon_0$
reverses the order.
\end{proof}

``One-state'' refers to the matrix language of \cref{onot:sh:sec:matrices}: every
matrix is the $1\times1$ matrix $(1)$. Already for $r=2$ an interval of exponents
below a fixed positive Hahn bound contains infinitely many terms, so a proof by
scanning all terms below a numerical cutoff would not work. This agrees with the
bound $<\omega^{r+1}$ for omnific supports in a rank-$(r+1)$ lattice, since
$\zeta_r$ uses $r+1$ scales including $\varepsilon_0$
(\cref{onot:hol:thm:support,onot:grid:thm:supportbound}). The ring $\SE$ is
deliberately not all of $\Oz$: it is not asserted to be closed under the omega
map on arbitrary inputs, under root extraction, or under ordinal-indexed sums,
and these restrictions keep its interface tractable (05).

\section{Finite-state series and finite leading-term certificates}\label{onot:sh:sec:matrices}

Let $k$ be an effective ordered field and $\delta_1,\ldots,\delta_r$ strictly
positive elements of an effective ordered additive group $G\subseteq\No$. For
dimension $m\ge1$ choose $u\in k^{1\times m}$, $v\in k^{m\times1}$ and pairwise
commuting $M_1,\ldots,M_r\in k^{m\times m}$, and put
\begin{equation}\label{onot:sh:eq:matrixseries}
 \Phi=\sum_{\mathbf n\in\N^r}a_{\mathbf n}t^{w(\mathbf n)},\qquad
 a_{\mathbf n}=uM_1^{n_1}\cdots M_r^{n_r}v,\qquad w(\mathbf n)=\sum_in_i\delta_i,
\end{equation}
adding contributions of coinciding weights. The positive-monoid support lemma
makes this a legitimate Hahn series even when $G$ is non-Archimedean, and a
shielding parameter turns $X\Phi$ into an omnific integer (\cref{onot:thm:shield}).
In commuting formal variables $z_i$ the associated series is
\begin{equation}\label{onot:sh:eq:resolvent}
 \widehat\Phi(\mathbf z)=u\prod_{i=1}^r(I-z_iM_i)^{-1}v.
\end{equation}
Finite linear representations and their zero tests are standard in rational
series and weighted-automata theory \cite{Fijalkow,Arvind}; the underlying
theory goes back to Sch\"utzenberger's rational series, a reference 05 does not
give. The question here is what the finite data certify after evaluation at
positive Hahn scales, particularly when different tuples collide.

\subsection{Independent scales: a small simplex controls the front}

If the scales are $\Q$-linearly independent, $w:\N^r\to G$ is injective and the
coefficient at $w(\mathbf n)$ is $a_{\mathbf n}$.

\begin{theorem}[Independent-scale finite front \src{05}]\label{onot:sh:thm:front}
If $\Phi\ne0$ in \eqref{onot:sh:eq:matrixseries} and the scales are
$\Q$-independent, its least Hahn exponent occurs at a multi-index with
\begin{equation}\label{onot:sh:eq:simplex}
 |\mathbf n|=n_1+\cdots+n_r<m.
\end{equation}
Consequently zero, the leading exponent, the leading coefficient and the sign are
determined by the finitely many coefficients with $|\mathbf n|<m$. For two
representations of dimensions $m$ and $m'$ on the same scale frame, use $m+m'$
after forming the direct-sum difference.
\end{theorem}
\begin{proof}
If $v=0$ the series is zero. Otherwise take $\mathbf n$ of least weight with
nonzero coefficient, which exists by well ordering; choose a word of length
$N=|\mathbf n|$ with $n_i$ letters $i$ and put $v_0=v$, $v_j=M_{i_j}v_{j-1}$.
Commutativity identifies $uv_N$ with $a_{\mathbf n}$. If $N\ge m$, some first
$j\le m$ has $v_j=\sum_{l<j}c_lv_l$, since the earlier vectors are independent in
dimension $m$. Apply the remaining suffix matrices and $u$; the left side is
nonzero, so some summand on the right is, and by commutativity it corresponds to
deleting the nonempty block of letters between positions $l+1$ and $j$. The new
multi-index is coordinatewise at most $\mathbf n$ and different, so of smaller
weight, a contradiction; thus $N<m$. The rest follows because independence rules
out cancellation at the least weight; a direct sum with block-diagonal commuting
matrices represents the difference.
\end{proof}

There are $\binom{m+r-1}{r}$ candidate multi-indices: a finite bound, not a
polynomial one when $r$ is part of the input. Reachable-state linear algebra can
reduce the zero-test work; the theorem additionally locates the leading Hahn term
without enumerating the possibly infinite set of weights below it. For $r=1$ this
is the Cayley--Hamilton bound, and it is sharp: a nilpotent shift of dimension $m$
can give $\widehat\Phi(z)=z^{m-1}$ (05; the same bound and example are 09's
\cref{onot:grid:prop:recurrence}).

\subsection{Dependent scales: collect collisions in a finite numerator}

Checking individual small coefficients is unsound when different multi-indices
have the same weight. Define
\begin{align}
 \mathsf Q(\mathbf z)&=\prod_{i=1}^r\det(I-z_iM_i),\label{onot:sh:eq:Q}\\
 \mathsf P(\mathbf z)&=u\prod_{i=1}^r\adj(I-z_iM_i)v,\label{onot:sh:eq:P}
\end{align}
so $\widehat\Phi=\mathsf P/\mathsf Q$, $\mathsf Q(0)=1$, and each variable has
degree at most $m-1$ in $\mathsf P$. Let $\ev_\delta$ map
$\mathbf z^{\mathbf n}$ to $t^{w(\mathbf n)}$ and collect identical weights.

\begin{theorem}[Collision-aware finite front \src{05}]\label{onot:sh:thm:collision}
For arbitrary strictly positive scales, including dependent and non-Archimedean
ones, let $\overline{\mathsf P}=\ev_\delta(\mathsf P)$. Then
\begin{equation}\label{onot:sh:eq:collisionzero}
 \Phi=0\iff\overline{\mathsf P}=0,
\end{equation}
and when nonzero, $\Phi$ and $\overline{\mathsf P}$ have the same leading Hahn
exponent and coefficient. In particular the valuation $v(\Phi)$ satisfies
\begin{equation}\label{onot:sh:eq:frontbound}
 0\le v(\Phi)\le(m-1)(\delta_1+\cdots+\delta_r).
\end{equation}
The polynomial $\mathsf P$, its at most $m^r$ monomial positions and the exact
collection of equal weights provide an equality and sign certificate.
\end{theorem}
\begin{proof}
The adjugate identity gives $\widehat\Phi=\mathsf P/\mathsf Q$ in formal power
series. Positive-scale evaluation is defined on these series and commutes with
products, since supports lie in the positive finitely generated monoid and fibres
are finite. $\overline{\mathsf Q}=\ev_\delta(\mathsf Q)$ has constant coefficient
one and all other exponents strictly positive, with no collision at zero because
every $\delta_i>0$; so it is invertible with leading term $1$, and
$\Phi=\overline{\mathsf P}\,\overline{\mathsf Q}^{-1}$ is zero iff
$\overline{\mathsf P}$ is, with the same leading term otherwise. Every monomial of
$\mathsf P$ has $0\le n_i\le m-1$, so its weight lies in the stated range;
collecting collisions removes or combines terms but creates no new weights.
\end{proof}

The proof uses only the ordered product in the adjugate formula; commutativity is
part of the chosen series language and of the shift-state interpretation, not a
hidden positivity assumption, and entries and coefficients may be negative (05).

\begin{example}[The independent-scale bound cannot be reused \src{05}]\label{onot:sh:ex:collision}
Let $N_0=\begin{pmatrix}0&0\\1&0\end{pmatrix}$, $M_1=N_0$, $M_2=I-N_0$, $u=(0,1)$,
$v=(1,0)^{\mathsf T}$; the matrices commute and $m=2$. Then
\[
 \widehat\Phi(z_1,z_2)=\frac{z_1-z_2-z_1z_2}{(1-z_2)^2},
\]
with coefficients $1$ and $-1$ at $(1,0)$ and $(0,1)$. If
$\delta_1=\delta_2=\delta>0$ these cancel and
$\Phi=-t^{2\delta}/(1-t^\delta)^2$: the leading exponent is $2\delta$, not the
weight of a surviving coefficient of total degree below $2$, and the
collision-aware bound $(m-1)(\delta_1+\delta_2)=2\delta$ is attained.
\end{example}

\subsection{Closure and transfinite support}

Direct sums give addition of descriptions; tensor products give coefficientwise
(Hadamard) products of the multivariate arrays before evaluation, which with
colliding scales need not be the Hadamard product of the collected Hahn
coefficients. Hahn multiplication is effective within a fixed frame.

\begin{proposition}[Cauchy-product closure \src{05}]\label{onot:sh:prop:matrixproduct}
Two commuting matrix-series descriptions of dimensions $m$ and $m'$ on $r$ fixed
scales have a Cauchy-product description by commuting matrices of dimension at
most $(m+m')^r$, and positive-scale evaluation of the multivariate product is the
Hahn product, even when scales collide.
\end{proposition}
\begin{proof}
Before evaluation write the series as $\mathsf P/\prod_iq_i(z_i)$ and
$\mathsf R/\prod_is_i(z_i)$ by the adjugate formulas. The product numerator has
degree at most $m+m'-2$ in each variable, and $q_is_i$ has degree at most
$b=m+m'$ and constant coefficient one. Multiplying by the denominator shows that
the coefficients satisfy a homogeneous constant-coefficient recurrence in each
coordinate once that coordinate is at least $b$, the others held fixed (missing
high recurrence coefficients are zero). Store the initial array on
$\{0,\ldots,b-1\}^r$ as a vector of dimension $b^r$; in coordinate $i$ the
companion shift advances that coordinate using its recurrence at the last
position. The shifts act on separate tensor factors and commute, and selecting
the all-zero coordinate after the shifts recovers every coefficient. Evaluation
respects the Cauchy product by the support lemma, including the finite
regrouping of colliding weights.
\end{proof}

For $\zeta_r$ of \eqref{onot:sh:eq:omegar} take $m=1$, $u=v=1$, $M_i=(1)$ and
$\delta_i=\omega^{i-1}$: the state dimension controlling the front is independent
of the ordinal type of the support, by finite algebraic dependence of the
coefficient states, not because transfinite supports became finite. Not every
rational multivariate function has this separated-resolvent form:
$1/(1-z_1-z_2)$ is rational, but as the second coordinate varies the sequences in
the first coordinate include polynomials of unbounded degree, which cannot all
satisfy one fixed-order recurrence. Such functions can still be treated in the
fraction-field calculus; the finite-state sublanguage is a certificate format,
not an identification of all rational multivariate series (05).

\part{Rational--omega expressions}\label{onot:grid:part}

This part is source 09's. It studies one particular language, not a universal
decision claim. Write $\Om(x)=\omega^x$ for Conway's omega map and consider
finite terms
\begin{equation}\label{onot:grid:eq:grammar}
 s::=q\mid -s\mid s+s\mid s\,s\mid s/s\mid\Om(s),\qquad q\in\Q,
\end{equation}
where division requires a nonzero denominator. Let $\QO$ be the set of their
values and $\IO=\QO\cap\Oz$. The integer terms are the valid terms whose
\emph{values} lie in $\IO$; intermediate values need not be omnific, and
infinitesimal intermediate expressions are essential for the compressed infinite
series below. The symbol $\Om$ is not a shorthand for the full surreal
exponential: no unrestricted $\exp$, $\log$, trigonometric function, arbitrary
computable real constant or program-controlled infinite sum is included, and the
restriction is substantive (09). Rational constants are given as exact integer
numerator--denominator pairs.

The three main theorems of the part are: validity, equality, order and
membership in $\Oz$ are decidable, and from names of $x,a\in\QO$ one computes
names of $\fl x$ and $\tr a x$ and the rational coefficient of $\omega^a$ in $x$,
all exactly (\cref{onot:grid:thm:equality,onot:grid:thm:effectivefloor}); the
ordinals of $\QO$ are exactly those below $\varepsilon_0$, while every support has
order type below $\omega^\omega$, sharply at depth two
(\cref{onot:grid:thm:ordinaltrace,onot:grid:thm:supportbound,onot:grid:thm:sharp});
and nevertheless equality of an always-valid stream language of omnific integers
is $\Pi^0_1$-complete on a subfamily whose values lie in the depth-two field with
at most one term (\cref{onot:neg:thm:notranslator}). The $\varepsilon_0$ bound is
about which ordinal values occur, the $\omega^\omega$ bound about the order type
of the support of one value, and the complexity result about a different input
language; none is a birthday bound (09).

\section{Finite rational--omega terms and decidable equality}\label{onot:grid:sec:equality}

\subsection{The depth hierarchy}

Set
\begin{equation}\label{onot:grid:eq:hierarchy}
 \QO_0=\Q,\qquad \QO_{d+1}=\Frac\bigl(\Q[\omega^{\QO_d}]\bigr),\qquad \QO=\bigcup_{d<\omega}\QO_d,
\end{equation}
where $\Q[\omega^{\QO_d}]$ consists of \emph{finite} sums $\sum q_i\omega^{g_i}$
with $g_i\in\QO_d$, negative exponents allowed: the group algebra embedded by
Conway monomials, not a Hahn completion, with fraction field taken in $\No$. The
hierarchy is increasing ($\Q\subseteq\QO_1$, and an expression of an element of
$\QO_d$ as a fraction over exponents in $\QO_{d-1}$ is one over $\QO_d$). A term
has depth zero if it has no $\Om$ node; $+$, $\times$, $/$ take the maximum of the
input depths and $\Om$ adds one. Induction shows that $\QO_d$ is exactly the set
of values of terms of depth at most $d$.

\begin{lemma}[Sparse-fraction representation \src{09}]\label{onot:grid:lem:sparse}
Every term of depth at most $d+1$ can be computed as a pair $(P,Q)$, $Q\ne0$, of
finite rational-coefficient sums of monomials with exponents of depth at most $d$.
Addition and multiplication of these sums need only rational arithmetic, addition
of exponent terms and equality of exponent values.
\end{lemma}
\begin{proof}
A rational uses the exponent zero, and $\Om(u)$ is the singleton monomial with
exponent $u$. Fractions are added and multiplied by the usual identities, and
division swaps a nonzero numerator and denominator. Multiplying finite sums uses
$\omega^g\omega^h=\omega^{g+h}$; equal exponents are collected and zero
coefficients deleted. All loops are over finite lists.
\end{proof}

\begin{theorem}[Recursive ground decision procedure \src{09}]\label{onot:grid:thm:equality}
Validity of \eqref{onot:grid:eq:grammar}, equality of valid terms and their
order are decidable. For a nonzero term of depth at most $d+1$ its leading
exponent is computable as a term in $\QO_d$ and its leading coefficient as a
nonzero rational.
\end{theorem}
\begin{proof}
Induct on depth; depth zero is rational arithmetic, including rejection of zero
denominators. Assume comparison and equality are decidable at depth $d$. Build
the representation of \cref{onot:grid:lem:sparse}, sort each exponent list with
the depth-$d$ comparator and collect equal values. By uniqueness of normal forms
a finite sum is zero exactly when its list is empty, and otherwise its greatest
exponent and coefficient are its leading data; so a division at the next depth is
valid exactly when the numerator of its denominator is nonzero. For valid
fractions
\begin{equation}\label{onot:grid:eq:cross}
 \frac PQ=\frac RS\iff PS-RQ=0,
\end{equation}
decided by the same collection; if the difference is nonzero, its sign is the sign
of $\lc(PS-RQ)$ times the signs of $Q$ and $S$. Finally
\begin{equation}\label{onot:grid:eq:fractionleading}
 \ddeg(P/Q)=\ddeg P-\ddeg Q,\qquad \lc(P/Q)=\lc(P)/\lc(Q),
\end{equation}
the exponent being a subtraction at depth $d$. Inputs of different depths are
viewed at the larger depth. Every exponent comparison descends strictly in depth,
and the rest is finite polynomial arithmetic, which proves termination.
\end{proof}

Semantic exponent equality matters: $u+v$ and $v+u$ must be merged, and treating
the written $\Om$ subterms as independent indeterminates would miss
$\Om(u+v)=\Om(u)\Om(v)$ and give an incorrect equality algorithm (09).

\begin{corollary}[Canonical representatives \src{09}]\label{onot:grid:cor:canonical}
There is a computable canonical finite-string representative for every value of
$\QO$, and likewise for every value of $\IO$ once membership has been decided.
\end{corollary}
\begin{proof}
Apply \cref{onot:prop:canonical} to the decidable set of valid codes. For $\IO$
the same least representative works, since equality to a member of $\IO$ implies
membership.
\end{proof}
This is an existence proof for a canonicalizer, not a practical implementation;
the prototype uses semantic comparison instead. No polynomial-time or
elementary-time claim is made for unrestricted nesting: expression expansion,
coefficient growth and repeated comparison can be expensive (09).

\subsection{Rational coefficients throughout the infinite normal form}

\begin{lemma}[Rational supports \src{09}]\label{onot:grid:lem:rationalsupport}
Every $x\in\QO$ has only rational normal-form coefficients. More precisely, if
$x=P/Q$ is a finite monomial fraction and $\Gamma$ is the additive group
generated by the supports of $P$ and $Q$, then $\supp x\subseteq\Gamma$.
\end{lemma}
\begin{proof}
Write $Q=c\omega^g(1-h)$ with $c\in\Q^\times$ and $h$ of finite support strictly
below exponent zero. The Hahn inverse
$Q^{-1}=c^{-1}\omega^{-g}\sum_{n\ge0}h^n$ is strongly summable by the
Hahn--Neumann support lemma \cite[Proposition~3.7]{Camacho}; every exponent is an
integer combination of the original ones and every coefficient a finite sum of
rational products. Multiplying by $P$ preserves both assertions. The argument
uses a semantic inverse to prove a support statement; the equality algorithm
does not enumerate this expansion.
\end{proof}

\begin{corollary}[Real trace \src{09}]\label{onot:grid:cor:realtrace}
$\QO\cap\R=\Q$. In particular $\QO$ is not real closed.
\end{corollary}
\begin{proof}
A nonzero real $r$ has the one-term normal form $r\omega^0$, so $r\in\Q$ by the
lemma; and $X^2-2$ has no root in $\QO$.
\end{proof}

\section{Effective separated grids}\label{onot:grid:sec:grid}

Equality alone does not decide omnific membership: a rational function may hide
infinitely many negative exponents. The next reduction makes the sign of every
possible exponent visible in finitely many coordinates.

\subsection{Elimination by leading scale}

For positive surreals write $u\gg v$ when $u>nv$ for every positive integer $n$.
For positive elements with rational leading coefficients, $\ddeg u>\ddeg v$
implies $u\gg v$.

\begin{lemma}[Effective separated basis \src{09}]\label{onot:grid:lem:basis}
Given $h_1,\ldots,h_m\in\QO$ one can compute a basis $\gamma_1,\ldots,\gamma_r$
of their rational span, with all coordinates, such that
\begin{equation}\label{onot:grid:eq:separation}
 \gamma_1\gg\gamma_2\gg\cdots\gg\gamma_r>0;
\end{equation}
the basis is empty if all inputs are zero, and it can be chosen in the same
$\QO_d$ as the inputs.
\end{lemma}
\begin{proof}
Discard zero inputs; take a vector of greatest leading exponent, make it positive,
and call it $\gamma_1$. Replace every other vector $v$ with the same leading
exponent by $v-(\lc(v)/\lc(\gamma_1))\gamma_1$; the multiplier is rational by
\cref{onot:grid:lem:rationalsupport} and the leading term cancels exactly, while
vectors of smaller leading exponent are unchanged. Discard zero residuals and
repeat. Each stage removes one pivot, so there are at most $m$ stages, and each
pivot has strictly smaller leading exponent than the previous one; after sign
normalization \eqref{onot:grid:eq:separation} holds. Distinct leading exponents
make the pivots rationally independent, the reversible row operations preserve
the span, and tracking them gives the coordinates. All leading data and
comparisons are computable by \cref{onot:grid:thm:equality}, and all operations
stay in the original level.
\end{proof}

The restriction to rational coefficients is essential: the rational span of $1$
and $\sqrt2$ is two-dimensional but has only one Archimedean scale, so it has no
separated rational basis. Substituting an arbitrary computable ordered
coefficient field into this lemma would be false (09).

\subsection{From a finite fraction to a lexicographic lattice}

\begin{theorem}[Finite separated-grid compilation \src{09}]\label{onot:grid:thm:grid}
Given $x\in\QO$ one can compute positive separated $\gamma_1,\ldots,\gamma_r\in\QO$
and $R\in\Q(X_1,\ldots,X_r)$ with $x=R(\omega^{\gamma_1},\ldots,\omega^{\gamma_r})$.
The map
\begin{equation}\label{onot:grid:eq:lattice}
 \Z^r_{\mathrm{lex}}\to\No,\qquad(n_1,\ldots,n_r)\mapsto\sum_in_i\gamma_i
\end{equation}
is an injective ordered-group map, the monomials $\omega^{\gamma_i}$ are
algebraically independent over $\Q$, and $\supp x$ lies in their exponent
lattice.
\end{theorem}
\begin{proof}
Apply \cref{onot:grid:lem:basis} to the union of the exponent lists of a
representation $P/Q$, and clear denominators in the coordinate matrix by
replacing $\gamma_i$ with $\gamma_i/N$; every original exponent then has integer
coordinates, and replacing $\omega^{\sum n_i\gamma_i}$ by $\prod X_i^{n_i}$ gives a
Laurent rational function, an element of $\Q(X_1,\ldots,X_r)$. For a nonzero
integer vector with first nonzero coordinate $n_i$, separation makes $n_i\gamma_i$
dominate the rest, so the sign of \eqref{onot:grid:eq:lattice} is that of $n_i$;
this gives order preservation and injectivity. A nonzero Laurent polynomial then
evaluates to a finite normal form with distinct exponents, which cannot vanish;
this is algebraic independence. Support containment is
\cref{onot:grid:lem:rationalsupport}.
\end{proof}

The basis depends on the input and its members may have infinite normal forms:
leading-scale elimination may need to distinguish $\omega+1/(\omega-1)$ from
$\omega$, for which the leading data of $1/(\omega-1)$ suffice. For a fixed grid
the relevant Hahn field is the iterated Laurent field
\begin{equation}\label{onot:grid:eq:iterated}
 \Q((X_r^{-1}))((X_{r-1}^{-1}))\cdots((X_1^{-1})),
\end{equation}
with $X_1^{-1}$ the \emph{outermost} variable: a positive power of $X_1$
dominates every monomial in the later variables, and a negative power is smaller
than every positive real multiple of such a monomial. This is the lexicographic
order of \eqref{onot:grid:eq:lattice}, not a numerical substitution. After
compilation the grid fields are the effective rational monomial core
\replabel{cas:thm-core} of the computer-algebra report, with $\Z$-independent
exponents; what 09 adds is the compilation of arbitrary nested terms into such a
frame, and the omnific recognition, floor and truncation below (\mergetag).

\begin{corollary}[Finite-grid transplantation \src{09}]\label{onot:grid:cor:transplant}
Replacing one positive separated $r$-tuple of scales by another induces an
ordered-field isomorphism of the corresponding rational grid fields which
preserves zero-exponent coefficients, the omnific condition and truncation at
zero.
\end{corollary}
\begin{proof}
Both evaluations identify their fields with the same rational function field
inside \eqref{onot:grid:eq:iterated}, and their lattices have the same
lexicographic order, so signs and formal coefficients agree, including the test
that all exponents are nonnegative.
\end{proof}
This concerns rational functions in the selected monomials; the isomorphism is
not asserted to commute with $\Om$ on arbitrary elements (09).

\section{Computable truncation and omnific recognition}\label{onot:grid:sec:truncation}

\subsection{A finite polynomial-part recursion}

Fix a positive separated grid and put $\KG_r=\Q(X_1,\ldots,X_r)$, with Hahn
interpretation \eqref{onot:grid:eq:iterated}; $\KG_0=\Q$. For $R\in\KG_r$ let
$\tr0 R$ be its nonnegative-exponent part and $\ct(R)$ its coefficient at
exponent zero in this interpretation.

\begin{theorem}[Recursive polynomial-part formula \src{09}]\label{onot:grid:thm:polypart}
For $R\in\KG_r$, $r\ge1$, divide in $X_1$ over $\KG_{r-1}=\Q(X_2,\ldots,X_r)$:
\begin{equation}\label{onot:grid:eq:polydiv}
 R=\sum_{j=0}^mc_jX_1^j+\frac UV,\qquad \deg_{X_1}U<\deg_{X_1}V,
\end{equation}
a zero polynomial part being allowed, with then $c_0=0$. Then
\begin{align}
 \tr0R&=\sum_{j=1}^mc_jX_1^j+\tr0(c_0),\label{onot:grid:eq:T}\\
 \ct(R)&=\ct(c_0),\label{onot:grid:eq:C}
\end{align}
where $c_0$ is read in the tail grid, and at rank zero $\tr0q=\ct(q)=q$. Both
outputs are computable, $\tr0R\in\KG_r$ and $\ct(R)\in\Q$.
\end{theorem}
\begin{proof}
Expand at $X_1=\infty$. The proper fraction $U/V$ has only strictly negative
powers of $X_1$: factoring the leading powers of $X_1$ from $U$ and $V$ gives a
negative initial power times a power series in $X_1^{-1}$ over $\KG_{r-1}$,
including the zero remainder. After expanding coefficients in the later variables,
every exponent of such a block has negative first coordinate and so is negative,
whatever the other coordinates. Every term of $c_jX_1^j$, $j>0$, has positive
first coordinate and positive exponent, and these whole blocks are kept even when
$c_j$ has an infinite normal form. Only the block $j=0$ needs a lower-rank
decision. Polynomial division over a rational function field is finite, and the
recursion follows only the constant block, so there are at most $r$ univariate
polynomial-part computations; their algebraic cost can grow, but the recursion
cannot become an infinite expansion.
\end{proof}

\begin{lstlisting}[caption={Truncation at zero on a separated grid (09).},label={onot:grid:lst:truncate}]
TRUNCATE_ZERO(R, [X1, ..., Xr]):
    if r = 0:
        return (R, R)
    P := polynomial part of R in X1 over Q(X2, ..., Xr)
    c := constant coefficient of P as a polynomial in X1
    (T, C) := TRUNCATE_ZERO(c, [X2, ..., Xr])
    return (P - c + T, C)
\end{lstlisting}
The first return value is the nonnegative part and the second the rational
zero-exponent coefficient.

\subsection{The omnific part of a grid}

Let $\IO_r=\KG_r\cap\Oz$ in the separated grid, and
$\IO_{r-1}^{\mathrm{tail}}=\Q(X_2,\ldots,X_r)\cap\Oz$; $\IO_0=\Z$.

\begin{theorem}[Recursive grid integer ring \src{09}]\label{onot:grid:thm:gridring}
For every $r\ge1$,
\begin{equation}\label{onot:grid:eq:gridring}
 \IO_r=\IO_{r-1}^{\mathrm{tail}}+X_1\Q(X_2,\ldots,X_r)[X_1],
\end{equation}
with a unique decomposition. Equivalently, a rational function is omnific exactly
when it is a polynomial in the largest grid variable whose positive-degree
coefficients are arbitrary elements of the tail field and whose constant
coefficient is omnific in the tail grid.
\end{theorem}
\begin{proof}
If $R$ is omnific its negative part vanishes, so $R=\tr0R$, which by
\cref{onot:grid:thm:polypart} is a polynomial in $X_1$ with positive-degree
coefficients $c_j$ and constant coefficient $\tr0(c_0)$, whose zero coefficient
must be an integer; so the constant lies in $\IO_{r-1}^{\mathrm{tail}}$.
Conversely, on the right side every positive-degree block has strictly positive
exponents and the constant block has nonnegative exponents and integer zero
coefficient. Uniqueness is uniqueness of polynomial expressions in the
algebraically independent $X_1$.
\end{proof}

This is a pullback-shaped description, not a claim to invent the ring
construction: in the notation of \cref{onot:sec:shield} it says that $\IO_r$ is
exactly the shielded ring $\mathcal S_{\gamma_1}(\IO^{\mathrm{tail}}_{r-1},\KG_{r-1})$
(\cref{onot:rem:characterizations}). Its use is algorithmic: it identifies which
rational denominators are harmless at a smaller scale. In rank one
$\IO_1=\Z+X_1\Q[X_1]$; in rank two
$\IO_2=(\Z+X_2\Q[X_2])+X_1\Q(X_2)[X_1]$, so a denominator in $X_2$ is allowed in a
positive $X_1$ block though it may make a constant block non-omnific (09).

\subsection{The effective integer-part theorem}

\begin{theorem}[Omnific recognition, floor, and internal truncations \src{09}]\label{onot:grid:thm:effectivefloor}
There are terminating exact algorithms on finite rational--omega terms that, for
$x,a\in\QO$,
\begin{enumerate}[label=\textup{(\roman*)}]
\item decide $x\in\Oz$;
\item return a term for $\fl x\in\IO$;
\item return a term for $\tr a x\in\QO$ and compute $[\omega^a]x\in\Q$.
\end{enumerate}
Consequently the valid omnific terms form a decidable finite-string notation
system with decidable equality and order.
\end{theorem}
\begin{proof}
Compile $x$ into a separated grid (\cref{onot:grid:thm:grid}) and apply
\cref{onot:grid:thm:polypart} to compute a rational function for $\tau=\tr0x$ and
the rational $c=\ct(x)$; evaluating back in the grid gives a finite term of $\QO$,
since every grid monomial is the omega map of a represented exponent. Then
\begin{equation}\label{onot:grid:eq:membership}
 x\in\Oz\iff x=\tau\ \text{and}\ c\in\Z,
\end{equation}
the first test by \cref{onot:grid:thm:equality} and the second by reducing a
rational. For the floor, compute $x_{<0}=x-\tau$ and compare it with zero, and
substitute $x_{>0}=\tau-c$, $c$ and $x_{<0}$ in \eqref{onot:eq:floor}; the rational
floor and the test $c\in\Z$ are computable, also for negative rationals. For an
arbitrary represented cutoff $a$ use
\begin{align}
 \tr a x&=\Om(a)\,\tr0\bigl(x/\Om(a)\bigr),\label{onot:grid:eq:shifttrunc}\\
 [\omega^a]x&=\ct\bigl(x/\Om(a)\bigr),\label{onot:grid:eq:shiftcoeff}
\end{align}
which hold because multiplication by $\omega^{-a}$ subtracts $a$ from every
exponent and keeps coefficients. The shifted input is again a finite term; $a$
need not lie in the original lattice, and $x=0$ is covered.
\end{proof}

The theorem supplies \emph{effective} truncation at a named cutoff. Abstract
truncation closure of a field generated by a truncation-closed set is already
known (Camacho's Proposition~4.1(i) \cite{Camacho}); the novelty assessment is
not a claim to abstract closure. What the proof adds is a computable route from
this syntax to a finite output syntax, an omnific recognizer, and a floor with the
infinitesimal boundary handled explicitly (09).

\begin{corollary}[Effective division with remainder \src{09}]\label{onot:grid:cor:division}
$\IO$ is a discretely ordered subring and an integer part of $\QO$. Given
$a,b\in\IO$ with $b>0$, the unique $q,r\in\Oz$ with $a=qb+r$ and $0\le r<b$ lie in
$\IO$ and are computable; $\ct:\IO\to\Z$ is a computable ring retraction.
\end{corollary}
\begin{proof}
The ring and discrete-order properties come from \eqref{onot:eq:oz}, and
$\fl x\in\IO$ for $x\in\QO$ by \cref{onot:grid:thm:effectivefloor}. Existence and
uniqueness of $q=\fl{a/b}$ and $r=a-qb$ in $\Oz$ are \replabel{odg:prop:division}
of the omnific Diophantine report; since $a/b\in\QO$, the theorem computes $q$ in
$\IO$, and then $r\in\IO$. The constant term is computable by
\eqref{onot:grid:eq:shiftcoeff} with $a=0$.
\end{proof}

This is not a claim that a Euclidean gcd algorithm terminates. The positive
omnific integers are not well ordered ($\omega,\omega-1,\omega-2,\ldots$ strictly
decreases), and successively smaller remainders need a separate termination
argument, which is not supplied (09); the omnific Diophantine report exhibits a
Euclidean algorithm in $\Oz$ that never stops (\replabel{odg:ex:euclid}).

\subsection{Worked examples}\label{onot:grid:sub:examples}

Put $X=\omega^\omega$ and $Y=\omega$, so the scales $\omega$ and $1$ are
separated. The rational expression
\begin{equation}\label{onot:grid:eq:Bexample}
 \sB=\frac X{1-Y^{-1}}=\sum_{n\ge0}\omega^{\omega-n}
\end{equation}
has infinitely many terms, all with positive exponents; it is omnific by
\cref{onot:grid:thm:gridring}, its coefficient $1/(1-Y^{-1})$ being allowed in a
positive $X$ block. Geometric examples of this kind have classical precedents
\cite{LM}; here they are tests for the decision procedure. On the other hand
$1/(Y-1)=\sum_{n\ge1}\omega^{-n}$ has zero nonnegative part, so $\sB+Y^{-1}$ is
not omnific, and neither is $\sB+\tfrac12$: one fails by a negative exponent, the
other by a nonintegral zero coefficient. Polynomial division gives
\begin{equation}\label{onot:grid:eq:divisionexample}
 \frac{X^2+Y}{X-1}=X+1+\frac{Y+1}{X-1},\qquad \tr0\Bigl(\frac{X^2+Y}{X-1}\Bigr)=X+1:
\end{equation}
every term of the remainder has negative $X$ coordinate, so the remainder is
infinitesimal although its numerator contains $Y$; similarly $\tr0(Y/(X-Y))=0$.
The floor correction appears without series expansion:
$\fl{\sB-Y^{-1}}=\sB-1$, $\fl{\sB+Y^{-1}}=\sB$, $\fl{\sB-\tfrac12}=\sB-1$. For a
cutoff off the lattice, $\tau_{\ge\omega-3/2}\sB=\omega^\omega+\omega^{\omega-1}$
and $[\omega^{\omega-3/2}]\sB=0$; the shifted-grid algorithm proves these by
finite algebra, without searching for a last term in an infinite stream.

\section{Ordinal traces}\label{onot:grid:sec:ordinals}

\subsection{Depth gives a tower bound, not a birthday bound}

Define $\beta_0=\omega$ and $\beta_{d+1}=\omega^{\beta_d}$; their supremum is
$\varepsilon_0$, the least positive fixed point of $\alpha\mapsto\omega^\alpha$.

\begin{lemma}[Strong ordinal majorants \src{09}]\label{onot:grid:lem:majorants}
For every $d$ and $x\in\QO_d$ there is an ordinal $\alpha<\beta_d$ with
$|x|<\alpha$.
\end{lemma}
\begin{proof}
At $d=0$ a rational is bounded by a natural number. For nonzero
$x\in\QO_{d+1}$, \cref{onot:grid:thm:equality} gives $e=\ddeg x\in\QO_d$ and a
nonzero rational leading coefficient; by induction choose $\delta<\beta_d$ with
$|e|<\delta$, and $\delta+1<\beta_d$ since $\beta_d$ is an infinite limit. Then
$|x|<\omega^{e+1}\le\omega^{\delta+1}$, because $|x|/\omega^e$ is a positive
finite surreal with rational standard part, below $\omega$; and
$\omega^{\delta+1}<\beta_{d+1}$. Zero is bounded by $1$.
\end{proof}
This supplies an ordinal strictly below the boundary that dominates the value,
which matters because a surreal below a limit ordinal need not be below an
ordinal strictly below that limit (09).

\begin{theorem}[Exact ordinal trace at every depth \src{09}]\label{onot:grid:thm:ordinaltrace}
For every $d<\omega$,
\begin{equation}\label{onot:grid:eq:ordinaltrace}
 \QO_d\cap\On=[0,\beta_d),\qquad \QO\cap\On=[0,\varepsilon_0).
\end{equation}
In particular $\varepsilon_0$ has no rational--omega notation.
\end{theorem}
\begin{proof}
The inclusion $\subseteq$ is \cref{onot:grid:lem:majorants}. Conversely induct on
$d$: the ordinals below $\omega$ are rational. If $\alpha<\omega^{\beta_d}$, its
Cantor normal form $\omega^{\alpha_1}n_1+\cdots+\omega^{\alpha_k}n_k$ has
$\beta_d>\alpha_1>\cdots>\alpha_k$, each $\alpha_i\in\QO_d$ by induction, so it is
a term of depth at most $d+1$. The decreasing Cantor form agrees with the Conway
normal form, so although surreal operations on ordinals are the natural ones, this
decreasing sum has the required value. The union over $d$ gives the second
identity.
\end{proof}

Division and negative or fractional intermediate values do not enlarge the
ordinal trace beyond the $\varepsilon_0$ range of hereditary finite Cantor
descriptions, but they enlarge the nonordinal omnific values very substantially.
The smaller hereditary hierarchy $H_0=\Q$, $H_{d+1}=\Q[\omega^{H_d}]$, with finite
outer support and no fractions, has the same ordinal trace by the same argument
but is strictly smaller than $\QO$: $1/(\omega-1)$ has infinite outer support. The
difference is invisible in the representable ordinals (09). The union of the
$H_d$ is 05's ring $\Hh_\Q$ (\cref{onot:sh:thm:ordinalfragment}), since both
consist of the finite $\Q$-linear combinations of $\omega$-powers of earlier
elements, starting from $\Q$ (\mergetag).

\subsection{Three systems, one trace}

The three systems of this report have the same ordinal trace,
$\On_{<\varepsilon_0}$: 08's labelled union by \cref{onot:hol:thm:epsilon}, 05's
hereditary ring by \cref{onot:sh:thm:ordinalfragment}, and $\QO$ by
\cref{onot:grid:thm:ordinaltrace}, which alone refines the trace depth by depth
(08 refines it by rank, \cref{onot:hol:thm:ordinalslice}). The trace does not
determine the system.

\begin{proposition}[The omnific rings are incomparable \mergetag]\label{onot:grid:prop:incomparable}
Let $\mathcal R_\varepsilon$ be the omnific part of the labelled holonomic union
over a coefficient field $k$, let $\SE$ be the $\varepsilon_0$-shielded ring over
$\Q$ or $\Ralg$, and let $\IO=\QO\cap\Oz$.
\begin{enumerate}[label=\textup{(\roman*)}]
\item $\SE$ and $\IO$ are incomparable, and so are $\SE$ and $\mathcal
R_\varepsilon$, for every choice of coefficient fields.
\item $\IO\not\subseteq\mathcal R_\varepsilon$ for every $k$, and
$\mathcal R_\varepsilon\not\subseteq\IO$ whenever $k\ne\Q$.
\end{enumerate}
All three contain every ordinal below $\varepsilon_0$ and no other ordinal, except
that $\SE$ also contains $\varepsilon_0$ and other ordinals $\ge\varepsilon_0$.
\end{proposition}
\begin{proof}
First, an element $x\in\SE$ which is bounded in absolute value by an ordinal
$\beta<\varepsilon_0$ has finite support. Indeed write
$x=d+\sum_{j=1}^m\varepsilon_0^{\,j}f_j$ with $f_m\ne0$, $m\ge1$. By
\cref{onot:thm:shield} the leading exponent of $x$ is
$m\varepsilon_0+\ddeg f_m$, and $\ddeg f_m$ lies in the group $\Hh$
(\cref{onot:sh:lem:fractionsupport}), so $|\ddeg f_m|<\alpha$ for some ordinal
$\alpha<\varepsilon_0$ (\cref{onot:sh:lem:epsilonbound}). Hence
$\ddeg x>\varepsilon_0-\alpha$. For $\beta<\varepsilon_0$ write
$\beta<\omega^{\gamma}$ with $\gamma<\varepsilon_0$; the natural sum
$\alpha\oplus\gamma$ is below $\varepsilon_0$, so $\varepsilon_0-\alpha>\gamma$ and
$|x|>\omega^\gamma>\beta$. So a bounded $x$ has all $f_j=0$ and $x=d\in\Dd$, a
finite hereditary form.

(i) $\varepsilon_0=\varepsilon_0^{\,1}\cdot1\in\SE$, while $\varepsilon_0\notin\IO$
by \cref{onot:grid:thm:ordinaltrace} and $\varepsilon_0\notin\mathcal R_\varepsilon$
because $\mathcal F_\varepsilon\subset(-\varepsilon_0,\varepsilon_0)$
(\cref{onot:hol:cor:ambientbound}). Conversely $\sB=\sum_{n\ge0}\omega^{\omega-n}$
lies in $\IO$ (\cref{onot:grid:sub:examples}) and in $\mathcal R_\varepsilon$ (it is
$U\Geo(t)$ with labels $\{0,1\}$), is bounded by the ordinal $\omega^{\omega+1}$,
and has infinite support, so $\sB\notin\SE$.

(ii) $\omega^{1/2}\in\IO$. Every exponent of an element of $\mathcal F_A$ lies in
$G_A=\sum_{\alpha\in A}\Z\omega^\alpha$ (\cref{onot:hol:lem:embedding}), and the
Conway normal form of an element of $G_A$ has integer coefficients, whereas
$\tfrac12=\tfrac12\omega^0$; so $\omega^{1/2}\notin\mathcal R_\varepsilon$. If
$c\in k\setminus\Q$, then $c\,\omega=cU_0\in\mathcal R_{\{0\}}$ by
\cref{onot:hol:thm:triangular}, while every element of $\QO$ has rational
coefficients (\cref{onot:grid:lem:rationalsupport}).

The trace statements are \cref{onot:hol:thm:epsilon,onot:sh:thm:ordinalfragment,onot:grid:thm:ordinaltrace};
$\varepsilon_0\cdot2$ and $\varepsilon_0^2$ lie in $\SE$.
\end{proof}

For $k=\Q$ it is not established here whether $\mathcal R_\varepsilon\subseteq\IO$;
for instance, whether $U\Fac(t)=\sum_nn!\,\omega^{\omega-n}$ or $U\Ex(t)$ has a
rational--omega term is not decided by any source, and neither is whether they
lie in the field $\FH$.

\subsection{What can and cannot be inferred from the Church--Kleene bound}

\begin{proposition}[An effective ordinal-order barrier \src{09; 05}]\label{onot:grid:prop:CK}
Let $N\subseteq\N$ be decidable and $\nu:N\to\On$, with equality and order of
the values decidable. Then the ordered range of $\nu$ has order type strictly
below $\CK$; in particular it cannot contain all computable ordinals. If in
addition the interpretation uniformly supplies computable presentations of the
well-orders of its values, the range is not cofinal in the computable ordinals.
\end{proposition}
\begin{proof}
Put $N^*=\{e\in N:\text{no }f<e\text{ in }N\text{ has }\nu(f)=\nu(e)\}$, a
decidable set since only finitely many smaller codes are tested. The semantic
order on $N^*$ is a computable well-order, so its type is a computable ordinal,
below $\CK$. For the cofinality statement, take the supplied presentation for
$e\in N$ and the empty order otherwise; the ordinal sum over $e\in\N$ (05 uses
$(L_0+1)+(L_1+1)+\cdots$) is a computable well-order on a decidable set of pairs,
comparing first the outer index and then the component order, and its type is a
computable ordinal at least as large as every value; its successor is a computable
strict upper bound.
\end{proof}

05 states the second half for a computably enumerable family of support
descriptions with uniformly computable presentations of their decreasing
well-orders: their order types are bounded by a computable ordinal and so are not
cofinal in $\CK$ (the finite case is a finite ordinal sum plus one). 09 adds the
first half and the following caveat.

\begin{remark}[The uniform-presentation hypothesis is essential \src{09}]\label{onot:grid:rem:CK}
The order-type conclusion alone does \emph{not} bound the supremum of an
arbitrarily interpreted range: a noncomputable increasing sequence of computable
ordinals can be cofinal in $\CK$ and still have the trivially decidable index
order of type $\omega$. For completely arbitrary semantic maps only the first
conclusion and the prohibition on containing every computable ordinal apply;
equality alone, or equality with an abstract order, is often overinterpreted.
\end{remark}

This is a classical computability observation rather than a new theorem, and it
does not apply directly to $\IO$ with its whole order, which is not well ordered;
\cref{onot:grid:thm:ordinaltrace} is a stronger language-specific statement that
does not use it (09). A bare finite symbol denoting an ordinal is not a computable
copy of its well-order, and the proposition supplies a bound for each sound
system, not a single algorithm extracting a least bound from every description.
A sound effective support calculus can be very expressive but cannot effectively
cover all recursive well-order types; adding larger named ordinals or rank
constructors may extend a calculus without creating a universal validator. The
one-state supports of type $\omega^r$ (\cref{onot:sh:cor:omegar}) against the hard
stream family of support type $\omega$ (\cref{onot:neg:sec:raw}) show that support
order type alone does not measure equality difficulty; coefficient representation
and the available certificates are decisive (05).

\section{Support complexity and values outside the language}\label{onot:grid:sec:support}

The outer normal form is a reverse well-ordered sequence of monomials whose length
$\ord(\supp x)$ need not resemble the ordinal value of $x$, even when $x$ is an
ordinal.

\begin{lemma}[Lexicographic support bound \src{09}]\label{onot:grid:lem:lexbound}
If $S\subseteq\Z^r$ is reverse well ordered lexicographically, its decreasing
order type is at most $\omega^r$; if $r\ge1$ and $S\subseteq(\Z^r)_{\ge0}$, it is
strictly below $\omega^r$.
\end{lemma}
\begin{proof}
For $r=0$ there is at most one element. For $r\ge1$ the first coordinates of a
nonempty reverse well-ordered set are bounded above, and listed decreasingly they
form a finite or $\omega$-sequence; each fibre has type at most $\omega^{r-1}$ by
induction, so the sum is at most $\omega^r$. If $S$ is nonnegative there are
finitely many fibres: for $r=1$ this is finiteness, and for $r>1$ the positive
fibres have type at most $\omega^{r-1}$ and the zero fibre, nonnegative in the
tail group, less than $\omega^{r-1}$, so the finite sum is below $\omega^r$.
\end{proof}

\begin{theorem}[The support barrier \src{09}]\label{onot:grid:thm:supportbound}
If a term compiles into a rank-$r$ grid, its support has type at most $\omega^r$,
and strictly less for an omnific value in a positive-rank grid. Consequently
\begin{equation}\label{onot:grid:eq:supportbarrier}
 x\in\QO\Longrightarrow\ord(\supp x)<\omega^\omega,
\end{equation}
independently of the omega-nesting depth.
\end{theorem}
\begin{proof}
Apply \cref{onot:grid:lem:lexbound} to the lattice support of
\cref{onot:grid:thm:grid}; omnific support lies in the nonnegative cone, the rank
is finite, and rank zero gives at most one term.
\end{proof}
The theorem does not assert that every support below the bound occurs: it is
necessary, not a characterization (09). The same bounds hold in the towers
(\cref{onot:hol:thm:support}), where every $\eta<\omega^r$ is realized.

\begin{theorem}[Sharpness at depth two \src{09}]\label{onot:grid:thm:sharp}
For every $d\ge1$ there is $\sB^\sharp_d\in\QO_2\cap\Oz$ with
$\ord(\supp\sB^\sharp_d)=\omega^d$; so the supremum of the support types of
$\QO_2\cap\Oz$ is already $\omega^\omega$.
\end{theorem}
\begin{proof}
Put $\gamma_0=\omega^d$ and $\gamma_i=\omega^{d-i}$ for $1\le i\le d$; they lie in
$\QO_1$ and $\gamma_0\gg\cdots\gg\gamma_d>0$. Let
\begin{equation}\label{onot:grid:eq:sharp}
 \sB^\sharp_d=\omega^{\gamma_0}\prod_{i=1}^d\bigl(1-\omega^{-\gamma_i}\bigr)^{-1}
 =\sum_{\mathbf n\in\N^d}\omega^{\gamma_0-n_1\gamma_1-\cdots-n_d\gamma_d}.
\end{equation}
Every exponent is positive since $\gamma_0$ dominates the later scales, distinct
multi-indices give distinct exponents, and the decreasing order corresponds to
lexicographically increasing $(n_1,\ldots,n_d)$, a $d$-fold ordinal sum of type
$\omega^d$. The sum is strong, all coefficients are one and no exponent is zero.
\end{proof}
The grid here has $d+1$ scales, consistent with the nonnegative-cone bound
$<\omega^{d+1}$ (09; compare \cref{onot:sh:cor:omegar}).

\begin{theorem}[Explicit summation escape \src{09}]\label{onot:grid:thm:escape}
There is an explicitly computable sequence of terms $\mathsf E_d\in\QO_2\cap\Oz$,
$d\ge1$, whose strong sum is an omnific integer $\mathsf E\notin\QO$ with support
type exactly $\omega^\omega$.
\end{theorem}
\begin{proof}
For $d\ge1$ put $b_d=\omega^{-d}$, $\delta_{d,i}=\omega^{-(d+i)}$ for $1\le i\le d$,
and
\begin{equation}\label{onot:grid:eq:escape}
 \mathsf E_d=\omega^{b_d}\prod_{i=1}^d\bigl(1-\omega^{-\delta_{d,i}}\bigr)^{-1},
\end{equation}
so $\mathsf E_d\in\QO_2$ with support
$S_d=\{b_d-\sum_in_i\delta_{d,i}:\mathbf n\in\N^d\}$ of type $\omega^d$, all
coefficients one and all exponents positive. For every tuple,
$b_d-\sum_in_i\delta_{d,i}>b_{d+1}$, since the difference has leading term
$\omega^{-d}$ with coefficient one; so $S_d$ lies above $S_{d+1}$, whose greatest
member is $b_{d+1}$. The union in decreasing order is the concatenation
$S_1,S_2,\ldots$, of type $\omega+\omega^2+\cdots=\omega^\omega$, reverse well
ordered with each exponent in one block; $\mathsf E=\sum_d\mathsf E_d$ is a strong
sum with coefficient one everywhere and positive support, so $\mathsf E\in\Oz$, and
$\mathsf E\notin\QO$ by \cref{onot:grid:thm:supportbound}. The map from $d$ to the
syntax \eqref{onot:grid:eq:escape} is computable.
\end{proof}
This exhibits a specific limit of the language, not a cardinality obstruction:
finite rational--omega arithmetic is not closed under explicitly generated strong
sums of its own omnific values (09).

\begin{proposition}[The finite-lattice obstruction \src{09}]\label{onot:grid:prop:denominators}
The omnific integer $\mathsf H=\sum_{n\ge1}\omega^{1/n}$ has support type
$\omega$ but is not in $\QO$.
\end{proposition}
\begin{proof}
The exponents $1,\tfrac12,\tfrac13,\ldots$ decrease and are positive, so this is a
valid omnific normal form of type $\omega$. If it were in $\QO$,
\cref{onot:grid:lem:rationalsupport} would put its support in a finitely generated
group $\Gamma$; then $\Gamma\cap\Q$ is finitely generated, hence cyclic, inside
$(1/N)\Z$ for some $N$, which does not contain all $1/n$.
\end{proof}
Three restrictions should therefore be remembered: rational coefficients, a
finitely generated exponent lattice for each input, and a support-length bound;
the last alone is far from sufficient. The example is a diagnostic consequence of
finite-grid support, not a new series (09).

\section{Certified infinite descriptions and Diophantine scope}\label{onot:grid:sec:certificates}

The contrast between finite expressions and arbitrary programs is not a
finite-versus-infinite dichotomy: a finite certificate can control an infinite
coefficient sequence. The two intermediate languages below use elementary linear
and polynomial algebra and are not new general theorems about automata or
recurrences (09).

\subsection{Linear-recurrence stream certificates}

Let $M\in\Q^{m\times m}$ and $u,v\in\Q^m$, and put $a_n=u^{\mathsf T}M^nv$. The
finite triple $(u,M,v)$ has the omnific value
\begin{equation}\label{onot:grid:eq:matrixstream}
 \sum_{n\ge0}a_n\omega^{\omega-n}=\omega^\omega u^{\mathsf T}(I-\omega^{-1}M)^{-1}v,
\end{equation}
the inverse existing because its determinant is a polynomial in $\omega^{-1}$
with constant coefficient one; the right side lies in $\QO_2\cap\Oz$, and the
formal geometric identity proves \eqref{onot:grid:eq:matrixstream}
coefficientwise.

\begin{proposition}[Finite recurrence witness bound \src{09}]\label{onot:grid:prop:recurrence}
For an $m$-dimensional triple the value is zero iff $a_0=\cdots=a_{m-1}=0$. Two
triples of dimensions $m$ and $m'$ are equal iff their first $m+m'$ coefficients
agree, and their order is decidable. The zero-test bound $m$ is sharp.
\end{proposition}
\begin{proof}
By Cayley--Hamilton, $M^m$ is a rational combination of $I,\ldots,M^{m-1}$, so the
$a_n$ satisfy a recurrence of order at most $m$ and vanish if the first $m$ do. The
block-diagonal matrix with output row $(u^{\mathsf T},-u'^{\mathsf T})$ represents
the difference in dimension $m+m'$; its first nonzero coefficient occurs before the
bound and gives the sign. For sharpness take the nilpotent shift $Me_i=e_{i+1}$,
$Me_m=0$, with $v=e_1$ and $u=e_m$: $a_n=0$ for $n<m-1$ and $a_{m-1}=1$.
\end{proof}
This is the case $r=1$ of 05's finite-front theorem
(\cref{onot:sh:thm:front}), shifted by $\omega^\omega$. It asks whether \emph{all}
coefficients vanish; it does not solve the Skolem problem of whether \emph{some}
coefficient of a general recurrence vanishes (09).

\subsection{Rational jets in several separated directions}

Let $P,Q,R,S\in\Q[t_1,\ldots,t_d]$ with $Q(0)S(0)\ne0$, so $P/Q$ and $R/S$ have
formal power-series expansions over $\N^d$. Choose $b_i\ge0$ bounding the degree
of $PS-RQ$ in each variable, for example
$b_i=\max\{0,\deg_iP+\deg_iS,\deg_iR+\deg_iQ\}$, omitting the degree terms of zero
polynomials.

\begin{theorem}[Finite rectangular jet certificate \src{09}]\label{onot:grid:thm:jets}
$P/Q=R/S$ iff their formal coefficients agree on the box
\begin{equation}\label{onot:grid:eq:jetbox}
 0\le n_i\le b_i\qquad(1\le i\le d).
\end{equation}
After substituting $t_i=\omega^{-\gamma_i}$ for positive separated
$\gamma_1\gg\cdots\gg\gamma_d$ the same criterion decides equality, and
multiplication by $\omega^{\gamma_0}$ with $\gamma_0\gg\gamma_1$ makes both series
omnific without changing the criterion.
\end{theorem}
\begin{proof}
Let $F=P/Q-R/S$ and $N_F=PS-RQ=QSF$. Equality implies agreement. Conversely, if
the coefficients of $F$ vanish on the box, the coefficient of $QSF$ at a box
index is a finite sum of products with coefficients of $F$ at indices
$\mathbf n-\mathbf j$ that stay in the box, which is closed downward; so $N_F$
vanishes on the box, and outside it by the degree bounds, hence $N_F=0$ and
$F=0$. Separated scales make $\mathbf n\mapsto-\sum n_i\gamma_i$ injective and
order-compatible, so no coefficients merge; after the shift every exponent is
positive, and multiplication by a nonzero monomial preserves equality.
\end{proof}
The certificate depends on the finite algebraic input; there is no universal
prefix length for all rational series, since the first nonzero coefficient can be
postponed by a high monomial, and the certificate cannot be transferred to
arbitrary primitive-recursive streams (\cref{onot:neg:thm:notranslator}) (09). The
theorem is a special case of the fraction-field setting of \cref{onot:sh:part}.

\subsection{Ground equality is not Diophantine decidability}

\begin{proposition}[Classical transfer consequence \src{09}]\label{onot:grid:prop:diophantine}
For every $P\in\Z[Y_1,\ldots,Y_m]$, $P$ has a zero in $\IO^m$ iff it has one in
$\Z^m$. Thus solvability of integer polynomial equations over $\IO$ is
undecidable, although equality of named elements is decidable.
\end{proposition}
\begin{proof}
Apply the ring retraction $\ct:\IO\to\Z$ to a solution, and use $\Z\subseteq\IO$
for the converse; the undecidability is the negative solution of Hilbert's tenth
problem \cite{Matiyasevich}. This is the argument of \replabel{odg:eq:existencetransfer}
and \replabel{odg:cor:H10} of the omnific Diophantine report, which prove the same
equivalence over all of $\Oz$, for positive existential formulas, and the
c.e.-completeness of the yes-instances.
\end{proof}
This is included to prevent a scope error, not as a new omnific-Diophantine
result: ground equality tests given tuples, while an existential question ranges
over all tuples and an equality algorithm only semidecides it by search. 09 wrote
that ``similar constant-term arguments already belong to the repository's omnific
arithmetic program'', citing the catalogue; the precise reference is
\replabel{odg:cor:H10} (\cref{onot:app:stale}).

\part{The complexity of equality, positivity, and validity}\label{onot:neg:part}

The positive systems carry certificates telling when an entire coefficient
sequence vanishes. This part removes that certificate without introducing any
uncertainty about totality or support validity (\cref{onot:neg:sec:raw,onot:neg:sec:rational,onot:neg:sec:ershov}),
and then releases the support geometry as well, exposing an independent
well-foundedness obstruction (\cref{onot:neg:sec:validity}). The base is 08's
\S\S11--13, with 05's \S\S7--9 and 09's \S8. Throughout $U=\omega^\omega$ and
$t=\omega^{-1}$ (05 wrote $X$ for $\omega^\omega$).

\section{Primitive-recursive stream names}\label{onot:neg:sec:raw}

\subsection{An everywhere-valid syntax}

Fix a standard decidable grammar of primitive-recursive functions $\N\to\N$. A
rational-valued program is a triple $(p,q,b)$ of such programs with value
$(p(n)-q(n))/(1+b(n))$, so its denominator is positive by syntax, not by an
undecidable promise; integer-valued programs are the case $b=0$. Every code is
total, and universal evaluation is computable although not itself primitive
recursive. Primitive-recursive syntax is used so that totality is guaranteed
syntactically; running arbitrary purported programs for total functions would
introduce a separate validity problem (05, 08). Boolean output can be enforced
syntactically by a final reduction modulo two (09).

For fixed $d\ge1$, a raw $d$-block name is a tuple of rational programs
$a_1,\ldots,a_d$ with denotation
\begin{equation}\label{onot:neg:eq:raw}
 Z(a_1,\ldots,a_d)=\sum_{j=1}^d\sum_{n\ge0}a_j(n)\omega^{j\omega-n}=\sum_{j=1}^dU^j\sum_{n\ge0}a_j(n)t^n.
\end{equation}
Here $j\omega$ is the surreal integer multiple of $\omega$, the ordinal
$\omega\cdot j$, not the ordinal product $j\cdot\omega$; this avoids a collapse of
the blocks (08). The case $d=1$, $Z(a)=\sum_na(n)\omega^{\omega-n}$, is 05's stream
notation with rational $a$ and 09's Boolean stream notation with $a:\N\to\{0,1\}$.

\begin{proposition}[Total validity and effective ring operations \src{08}]\label{onot:neg:prop:rawvalid}
Every raw $d$-block name denotes an omnific integer, and distinct pairs $(j,n)$
have distinct exponents. Allowing finitely many blocks and an integer constant
gives a notation ring with computable addition, negation and multiplication of
codes.
\end{proposition}
\begin{proof}
All exponents $j\omega-n$ are positive and decrease with $n$ within a block, and
every exponent of block $j+1$ exceeds every exponent of block $j$; a finite union
of such blocks is reverse well ordered, also after zero coefficients are
omitted. Distinctness is the independence of $1$ and $\omega$. Addition and
negation are coefficientwise; in a product the coefficient in block $j$ at index
$n$ is a finite sum of products with block indices adding to $j$ and ordinary
indices adding to $n$, and bounded sums and products of primitive-recursive
rational functions are primitive recursive. The number of blocks stays finite and
all exponents except the constant stay positive.
\end{proof}

For the positivity classification of \cref{onot:neg:sec:ershov} the zero-constant
syntax \eqref{onot:neg:eq:raw} is used; adding a positive ordinary constant changes
the initial-default bookkeeping of the mind-change bound and is not silently
included (08). Zero coefficients disappear from the semantic support; every name
is valid, its support lies in a fixed reverse well-ordered set, and for $d=1$ its
support type is at most $\omega$ (09).

\subsection{The exact equality complexity}

\begin{theorem}[Raw equality and zero testing \src{08; 05, 09}]\label{onot:neg:thm:rawequality}
For every fixed $d\ge1$, equality of raw $d$-block names is $\Pi^0_1$-complete,
and so is the zero test. Inequality is c.e., and equality is not c.e. Hardness
holds on a uniformly generated family whose values have at most one nonzero
monomial, with coefficients in $\{0,1\}$. For $d=1$, each strict comparison
relation is $\Sigma^0_1$-complete.
\end{theorem}
\begin{proof}
Normal-form uniqueness and the fixed grid of distinct exponents give
\[
 Z(a_1,\ldots,a_d)=Z(b_1,\ldots,b_d)\iff\forall n\,\forall j\le d\ [a_j(n)=b_j(n)],
\]
a computable matrix under one universal quantifier; a mismatch is found by
evaluating the programs at $n=0,1,\ldots$, so inequality is c.e. For hardness let
$h_e(n)=1$ if machine $e$ on empty input halts for the first time at step $n$ and
$0$ otherwise; bounded simulation makes $h_e$ primitive recursive uniformly in
$e$. With $a_1=h_e$ and the other blocks zero,
\begin{equation}\label{onot:neg:eq:spike}
 Z_e=\sum_{n\ge0}h_e(n)\omega^{\omega-n}=
 \begin{cases}0,&e\text{ never halts},\\ \omega^{\omega-N},&e\text{ first halts at step }N,\end{cases}
\end{equation}
so $Z_e=0$ is exactly nonhalting, a $\Pi^0_1$-complete property, and every value
has at most one monomial; comparison with a fixed zero name gives the pair
version. If equality were c.e., halting and nonhalting would both be c.e.

For $d=1$ and strict comparison, search for the first index where the
coefficients differ; its sign decides, so $Z(a)<Z(b)$ has the c.e.\ witness
$\exists n\,(a(n)<b(n)\wedge\forall i<n\ a(i)=b(i))$, and the search need not stop
on equal inputs. Finally $0<Z(h_e)$ exactly when $e$ halts, which gives
$\Sigma^0_1$-hardness; reversing the pair handles the other relation.
\end{proof}

\paragraph{Second reduction.} 05 reduces with the cumulative predicate
$\hb_e(n)=1$ iff $e$ has halted within $n$ steps: $Z(\hb_e)=0$ iff $e$ never
halts, and $Z(\hb_e)=U t^s/(1-t)$ if $e$ first halts at step $s$, so the hard
values have full tails rather than one term (\cref{onot:neg:sec:rational}).

\paragraph{Credit.} Both reductions are those of \replabel{cas:thm-zerotest} of
the computer-algebra report, which proves that total coefficient access does not
decide zero for Hahn series at $t=\omega^{-1}$: its Reduction~A is $\hb_e$, with
the moral that values in a decidable rational field become undecidable when
presented by unrestricted programs (\replabel{cas:rem-moralA}); its Reduction~B is
the first-halt spike $h_e$, with the moral that the obstruction survives a promise
of at most one nonzero coefficient (\replabel{cas:rem-moralB}). The report was
present at the sources' pin; 08 read only its README and counted the
one-monomial localization as its own, and 05 and 09 did not cite it
(\cref{onot:app:stale}). What is new here is the placement inside always-valid
omnific names on positive supports, the $d$-block form, the translator
corollaries and the exact realization dimension of \cref{onot:neg:thm:statebound}.
The method is a standard halting reduction, not a new undecidability principle
(08, 05, 09).

The coefficient field here is $\Q$, the support is prescribed, and every
candidate exponent is an elementary hereditary expression; none of the hardness
comes from exotic real coefficients, uncertain well-foundedness or large ordinal
notations, but from the universal demand that no later nonzero coefficient
exists (05). It cannot be evaded by requiring a support certificate, restricting
coefficients to $\{0,1\}$ or assuming that every program is total (09).

\begin{corollary}[Known leading data do not decide equality \src{08}]\label{onot:neg:cor:knownleading}
There is no equality algorithm for a uniformly named family of nonzero omnific
integers with the same leading exponent, leading coefficient and constant term,
even if each support has at most two terms.
\end{corollary}
\begin{proof}
Compare $Y_e=U^2+Z_e$ with $U^2$: the leading exponent $2\omega$, the leading
coefficient $1$ and the constant $0$ do not depend on $e$, and equality holds
exactly when $e$ never halts.
\end{proof}

\subsection{The same values in incompatible representations}

\begin{theorem}[Same values, no translator \src{08, 09}]\label{onot:neg:thm:notranslator}
Every value $Z_e$ of \eqref{onot:neg:eq:spike} has support of size at most one,
lies in the tower ring $R_2$ (labels $0,1$) as the rational monomial $Ut^N$ or $0$,
and lies in $\QO_2\cap\Oz$ as
\[
 \omega^{\omega-N}=\frac{\Om(\Om(1))}{\Om(1)^N}.
\]
Nevertheless no total computable map translates all these raw names into value-
equivalent names of any notation system with decidable equality and an available
name for zero. In particular there is no uniform compiler from these stream names
into the holonomic towers or into rational--omega terms, although a suitable finite
term exists for each value.
\end{theorem}
\begin{proof}
A machine has at most one first halting step, which gives the values. If a
translator existed, translate $Z_e$ and compare with the translated zero name;
this would decide nonhalting, contrary to \cref{onot:neg:thm:rawequality}.
\end{proof}

The obstruction is not that the values lack exact finite descriptions but the
absence of a uniform algorithm extracting them from this input representation
(08). Here the \emph{semantic normal form is finite}, with at most one term; what
is missing is an effective certificate of whether a term occurs and where. A
promise of finite support is not an explicit support list or a computable bound
on its positions (09). The existence of canonical Conway normal forms does not
give a computable canonicalizer: on the decidable term language one exists
(\cref{onot:grid:cor:canonical}), but on these stream indices a canonicalizer into
that language would be the forbidden compiler; the input representations differ
(09).

\subsection{What finite-support information suffices}

\begin{proposition}[An exact support count is enough \src{08}]\label{onot:neg:prop:count}
For a raw finite-block name promised to have exactly $s<\omega$ nonzero
coefficients, with $s$ supplied, its finite normal form is computable. Replacing
the exact count by the upper bound ``at most one'' does not give an equality
algorithm.
\end{proposition}
\begin{proof}
If $s=0$ return zero. Otherwise evaluate all blocks on successive complete
prefixes until $s$ nonzero coefficients are found; the promise makes the search
terminate and guarantees no other terms. Sort and return. The failure for an upper
bound is \eqref{onot:neg:eq:spike}.
\end{proof}
This is a promise algorithm and does not certify the count; an erroneous count
can make it diverge or return an incorrect truncation, so the count must come from
a trusted certificate or be an explicit precondition (08).

\begin{corollary}[No uniform finite observation bound \src{08}]\label{onot:neg:cor:nojet}
No total computable function of a raw name gives a prefix length whose vanishing in
every block guarantees that the value is zero. There is no sound, c.e., complete
proof system for equality of these names.
\end{corollary}
\begin{proof}
A prefix bound would reduce the zero test to finitely many coefficient tests. A
sound and complete enumerable proof system would make equality c.e.; inequality is
c.e., so dovetailing would decide equality. Both contradict
\cref{onot:neg:thm:rawequality}.
\end{proof}

\subsection{A single computable real coefficient}

\begin{proposition}[One-term obstruction for Cauchy names \src{05, 09}]\label{onot:neg:prop:realcoeff}
A finite-list language allowing arbitrary computable real coefficients, given by
fast rational Cauchy names, need not have decidable equality; the failure occurs
already for one-term names $r\omega$.
\end{proposition}
\begin{proof}
Let $r_e=0$ if machine $e$ never halts and $r_e=2^{-s}$ if it first halts at step
$s$ (09; 05 uses $2^{-(s+1)}$). At precision $n$, simulate $n$ steps and output the
corresponding value if $e$ has halted and $0$ otherwise; a later halt changes the
value by at most $2^{-(n+1)}$, so this is a uniform fast Cauchy name. Every
$r_e\omega$ is omnific and is zero exactly when $e$ never halts.
\end{proof}
Every $r_e$ here is a dyadic rational; the obstruction lies in extracting an exact
finite answer from an intensional name (09). This is why exact algebraic codes,
not arbitrary Cauchy names, were specified for the positive systems (05). The
constants are those of \replabel{thm:inversehard}(iv) of the computable-surreals
report, which proves $\Pi^0_1$-complete equality on raw numerical names with the
same family (\cref{onot:sec:neighbours}).

\section{The rational-realization gap}\label{onot:neg:sec:rational}

This section is 05's. The obstruction can be strengthened to a family whose values
all have simple rational generating functions and identical supports. Let
\begin{equation}\label{onot:neg:eq:switch}
 a_e(n)=1+\hb_e(n),\qquad A^{\mathrm{sw}}_e=U\sum_{n\ge0}a_e(n)t^n,\qquad \sB=\frac U{1-t}.
\end{equation}
All coefficients of $A^{\mathrm{sw}}_e$ are $1$ or $2$, and every support is
exactly $\{\omega-n:n\in\N\}$.

\begin{theorem}[Rational full-support hardness \src{05}]\label{onot:neg:thm:rationalhard}
The uniformly primitive-recursive family \eqref{onot:neg:eq:switch} satisfies:
\begin{enumerate}[label=\textup{(\roman*)}]
\item every value is rationally generated: it is $\sB$ if $e$ never halts, and
$A^{\mathrm{sw}}_e=U(1+t^s)/(1-t)$ if $e$ first halts at step $s$;
\item the equality $A^{\mathrm{sw}}_e=\sB$ is $\Pi^0_1$-complete, and equality among
the $A^{\mathrm{sw}}_e$ is $\Pi^0_1$-complete after choosing a known nonhalting
baseline code;
\item no total computable translator takes these stream codes to explicit
rational-function codes denoting the same values, nor to any notation system with
decidable equality and an effective notation for $\sB$.
\end{enumerate}
\end{theorem}
\begin{proof}
If halting first occurs at $s$, the added coefficients are $0$ before $s$ and $1$
from $s$ on, a geometric tail $t^s/(1-t)$; without halting the added stream is
zero. This gives (i), and equality to $\sB$ is exactly nonhalting; the
coefficientwise upper bound gives (ii). A translator as in (iii), followed by
cross multiplication of rational functions or by the target system's equality test
against its name for $\sB$, would decide nonhalting.
\end{proof}

The difficulty is not that arbitrary streams denote nonrational functions: every
member of this syntactically valid family is rational, but the rational
description cannot be extracted uniformly from the stream program (05).

\begin{corollary}[The same decidable ring admits hard names \src{05}]\label{onot:neg:cor:samering}
There is a uniformly primitive-recursive family of stream names whose values all
lie in $\SE$ (\cref{onot:sh:thm:epsilon}), but whose equality to one fixed member is
$\Pi^0_1$-complete; the same values with explicit polynomial--fraction names in
$\SE$ have decidable equality.
\end{corollary}
\begin{proof}
Replace $U$ by $\varepsilon_0$ in \eqref{onot:neg:eq:switch}. Each stream then
denotes $\varepsilon_0/(1-t)$ or $\varepsilon_0(1+t^s)/(1-t)$, whose rational
coefficient lies in $\FH$, so the value lies in $\SE$; equality to
$\varepsilon_0/(1-t)$ still means nonhalting. The explicit-name algorithm is
\cref{onot:sh:thm:epsilon}.
\end{proof}
Restricting attention to a fixed ring with a decidable presentation does not remove
representation dependence: it is the uniform availability of the finite
certificate, not its existence somewhere, that supplies the algorithm (05).

\paragraph{Credit.} Parts (i)--(iii) are anticipated in substance by
\replabel{cas:rem-moralA}, whose values $t^N/(1-t)$ lie in a decidable rational
field. New in 05 are the omnific placement with full support and coefficients $1$
and $2$, the same-ring corollary, and the exact dimension count that follows.

\subsection{Exact minimal dimension and absence of an effective bound}

For a rational power series $f(t)$ over $k$, a univariate linear realization of
dimension $m$ is $f(t)=u(I-tM)^{-1}v$ with $M\in k^{m\times m}$; no minimality
promise is imposed on the input.

\begin{theorem}[Unbounded realization complexity \src{05}]\label{onot:neg:thm:statebound}
Let $f_e(t)=\sum_na_e(n)t^n$. If $e$ never halts its minimal nonzero realization
dimension is $1$; if $e$ first halts at step $s\ge0$ it is exactly $s+1$.
Consequently no total computable function of the stream program code bounds the
minimal realization dimension on this family.
\end{theorem}
\begin{proof}
The nonhalting series $(1-t)^{-1}$ and the series $2/(1-t)$ for $s=0$ have
one-dimensional realizations. For $s\ge1$, $f_e=(1+t^s)/(1-t)$, and numerator and
denominator are coprime since the numerator is $2$ at $t=1$. In dimension $m$ the
adjugate formula writes $f_e$ with numerator degree at most $m-1$, and reduction
cannot raise that degree; the reduced numerator has degree $s$, so $m\ge s+1$.
Conversely the sequence is $s$ ones followed by twos and satisfies
$a(n+s+1)=a(n+s)$, so a companion shift on windows of $s+1$ coefficients realizes
it in dimension $s+1$. If a computable bound $b(e)$ existed, simulating $e$ for
$b(e)$ steps (one more if the convention requires) would decide halting, since a
later first halt at $s$ would give $s+1>b(e)$.
\end{proof}

This is more precise than saying normalization is difficult: a finite-state
certificate exists for every input, but even its size cannot be bounded computably
from this kind of name. The finite-front theorem (\cref{onot:sh:thm:front}) is not
contradicted: it needs the realization itself, not a promise that one exists (05).

\begin{corollary}[No complete effective equality certificates \src{05}]\label{onot:neg:cor:proofobstruction}
No computably enumerable sound proof system proves every true equation
$A^{\mathrm{sw}}_e=\sB$, and there is no total computable stopping bound for the
number of coefficients to check.
\end{corollary}
\begin{proof}
Enumerating proofs would enumerate the nonhalting codes; a stopping bound would
decide the same set by bounded checking.
\end{proof}
An incomplete proof-producing simplifier is compatible with this: a sound
implementation distinguishes ``equal'', ``different'' and ``not established'' and
never turns a timeout into a false equality (05).

\section{The finite mind-change spectrum of positivity}\label{onot:neg:sec:ershov}

This section is 08's. Equality of raw names is controlled by one universal
quantifier; order has a different obstruction, since a later discovery in a higher
block overrides every conclusion based on a lower one. The result is an exact
classification in the finite Ershov hierarchy rather than another undecidability
proof.

\begin{definition}[$d$-c.e.\ set]\label{onot:neg:def:dce}
For $d\ge1$, $A\subseteq\N$ is $d$-c.e.\ if there is a total computable
approximation $a(e,s)\in\{0,1\}$ with $a(e,0)=0$, $\lim_sa(e,s)=\mathbf 1_A(e)$, and
at most $d$ changes of value for each $e$. The $1$-c.e.\ sets are the c.e.\ sets,
the $2$-c.e.\ sets the d.c.e.\ sets; the finite levels and their strictness are
standard \cite{Ershov}.
\end{definition}

\begin{theorem}[Dominance-layer completeness \src{08}]\label{onot:neg:thm:ershov}
For $d\ge1$ let $\Pos_d$ be the set of raw $d$-block names \eqref{onot:neg:eq:raw}
denoting a strictly positive omnific integer. $\Pos_d$ is many-one complete for
the $d$-c.e.\ sets. The hardness reduction can be chosen with at most one nonzero
coefficient per block, each equal to $\pm1$, and total support of at most $d$
terms. Equality on the unrestricted raw $d$-block syntax remains
$\Pi^0_1$-complete.
\end{theorem}
\begin{proof}[Upper bound]
At stage $s\ge1$ compute the coefficients with $0\le n<s$ in every block; a block
is discovered if this prefix contains a nonzero coefficient, and then its first
nonzero coefficient and its sign are known permanently. If no block is discovered
output $0$; otherwise let $j_s$ be the largest discovered block and output $1$
exactly when its first nonzero coefficient is positive; stage zero outputs $0$. The
output can change only when a higher block replaces the current highest, the
record indices increase strictly and there are at most $d$ of them, so there are at
most $d$ changes, including a possible initial change from $0$ to $1$. Every
nonzero block is eventually discovered; if all blocks vanish the output stays $0$,
and otherwise the final highest block is the highest nonzero block, whose first
nonzero coefficient is the leading coefficient of the number.
\end{proof}
\begin{proof}[Lower bound]
Let $A$ be $d$-c.e.\ with a fixed approximation starting at $0$. For input $e$ run
a machine computing the successive approximation values and emitting a change event
whenever the value flips, at most $d$ times; computing a stage may take unbounded
finite time, which is harmless because the predicate below uses machine steps, not
stages. Let $\chi_{i,e}(n)$ say that machine step $n$ is exactly the $i$th change
event; bounded simulation makes it primitive recursive uniformly in $(i,e,n)$, and
it is $1$ for at most one $n$. Put
\begin{equation}\label{onot:neg:eq:mindchangecode}
 a_{i,e}(n)=(-1)^{i+1}\chi_{i,e}(n)\qquad(1\le i\le d),
\end{equation}
with codes computable from $e$. With no changes the value is $0$ and $e\notin A$.
With $j\ge1$ changes the highest occupied block is $j$, with sign $(-1)^{j+1}$, so
the value is positive exactly when $j$ is odd, which is exactly when the final
approximation value is $1$. Hence $e\in A\iff Z(a_{1,e},\ldots,a_{d,e})>0$. The
equality statement is \cref{onot:neg:thm:rawequality}.
\end{proof}

\begin{corollary}[Strict order hierarchy with fixed equality complexity \src{08}]\label{onot:neg:cor:hierarchy}
One raw block has c.e.-complete positivity, two blocks d.c.e.-complete
positivity, and $d$ blocks the $d$th finite mind-change level; these levels are
strict, although all the equality problems are $\Pi^0_1$-complete.
\end{corollary}
\begin{proof}
Apply \cref{onot:neg:thm:ershov} and the strictness of the hierarchy
\cite{Ershov}: a complete set for level $d+1$ at level $d$ would collapse the two
levels, since the $d$-c.e.\ sets are closed under computable preimages.
\end{proof}
The case $d=1$ is 05's $\Sigma^0_1$-complete strict comparison
(\cref{onot:neg:thm:rawequality}). The theorem is insensitive to the times of the
change events: they only fix the offsets $-n$ in the exponents, while any term in a
higher block dominates every term in a lower one. This is the mechanism converting
logical mind changes into surreal dominance changes (08).

\begin{example}[A sign obstruction under a nonzero promise \src{08}]\label{onot:neg:ex:We}
For the spike $Z_e$ put
\[
 W_e=U-UZ_e=U-\sum_{n\ge0}h_e(n)\omega^{2\omega-n}.
\]
Every $W_e$ is nonzero with at most two monomials. If $e$ does not halt, $W_e=U>0$;
if it halts, the negative term in block two dominates and $W_e<0$. So a nonzero
promise does not make the two-block sign problem decidable. For one block under a
nonzero promise, scanning to the first nonzero coefficient terminates and decides
the sign.
\end{example}

This example is \replabel{cas:prop-leading} of the computer-algebra report (the
series $t^{(1,0)}-\sum_na_nt^{(0,n)}$ with first-halt coefficients, always nonzero
with one or two terms), and the one-block contrast is its
\replabel{cas:rem-rankone}. \Cref{onot:neg:thm:ershov} strengthens that
proposition from undecidability of the sign to an exact classification of
positivity at every finite number of blocks.

\begin{proposition}[Known highest occupied block \src{08}]\label{onot:neg:prop:highestblock}
For a raw finite-block name supplied with the index of its highest nonzero block,
and the promise that all higher blocks vanish, sign and leading monomial are
computable. Neither this information nor the full leading monomial decides
equality in general.
\end{proposition}
\begin{proof}
Scan the specified block to its first nonzero coefficient, which gives the leading
monomial and the sign; the promise makes this correct and terminating. The
equality counterexample is \cref{onot:neg:cor:knownleading}.
\end{proof}

An ordinary halting oracle $0'$ decides whether each of the finitely many blocks is
identically zero, finds the highest nonzero block and then applies
\cref{onot:neg:prop:highestblock}; this gives a uniform $0'$ upper bound for order
over the union of all finite-block syntaxes. The one-block equality lower bound
shows that this oracle strength is necessary for a general equality decision, and
no stronger oracle is concealed in the finite-level results (08).

\section{Support validity is a separate coanalytic problem}\label{onot:neg:sec:validity}

The raw block construction made validity automatic. To expose the separate cost of
recognizing a valid support it suffices to use dyadic exponents in a fixed bounded
interval and coefficient one everywhere, with no large ordinal labels, exotic real
constants or varying exponent algebra (08, 05).

\subsection{Trees and the Kleene--Brouwer order}

A tree $T\subseteq\N^{<\N}$ is closed under prefixes (and contains the empty
sequence when nonempty). Its Kleene--Brouwer order puts $s\kb u$ when $u$ is a
proper prefix of $s$, or, at the first differing coordinate, $s$ has the smaller
entry; it is a computable linear order on $\N^{<\N}$.

\begin{lemma}[Tree well-foundedness and its linearization \src{08; 05, 09}]\label{onot:neg:lem:kbtree}
A tree $T\subseteq\N^{<\N}$ has no infinite branch iff $(T,\kb)$ is a well-order.
\end{lemma}
\begin{proof}
An infinite branch gives the strictly $\kb$-descending chain of its longer and
longer prefixes. Conversely, if $T$ has no infinite branch, proper extension is
well founded on $T$; induct on it. The subtree at $s$ is, in the
Kleene--Brouwer order, the ordinal sum in increasing $n$ of the subtrees at the
children $s^\frown n$, followed by the singleton $s$; a countable ordered sum of
well-orders followed by one point is a well-order, which applies at the root. No
finite-branching assumption is used. (09 argues the converse directly: a
descending chain yields a path by successively stabilizing its natural
coordinates, since there is no infinite strict decrease of a natural coordinate.)
\end{proof}

Well-foundedness of recursive, and of primitive-recursive, trees is
$\Pi^1_1$-complete: this standard fact is recorded in
\cite[Theorem~2.6(g)]{CMR} (08), in \cite[Theorems~5.2.10 and 7.1.4]{Pohlers}
(05) and in the tree normal forms discussed in \cite[Section~2]{DowneyMontalban}
(09). Tree predicates can be made primitive recursive by including bounded
computation witnesses in the nodes; this is the usual effective tree normal form
(09). Treehood is enforced syntactically by taking as $T$ the set of nodes all of
whose prefixes pass a primitive-recursive test (08, 09). $\Pi^1_1$ is not
$\Pi^0_1$.

\subsection{An explicit dyadic coding with decidable image}

Associate an open interval $(a_s,b_s)$ to each $s\in\N^{<\N}$: start with
$(a_\varnothing,b_\varnothing)=(0,1)$ and, with $L=b_s-a_s$, put
\begin{equation}\label{onot:neg:eq:intervals}
 a_{s^\frown n}=a_s+L(1-2^{-n}),\qquad b_{s^\frown n}=a_s+L(1-2^{-(n+1)}).
\end{equation}
Let $\iota(s)=b_s$ and
\begin{equation}\label{onot:neg:eq:dyadic}
 \kappa(s)=2-\iota(s).
\end{equation}
For example $\iota(\varnothing)=1$, $\iota(0)=\tfrac12$, $\iota(1)=\tfrac34$,
$\iota(0,0)=\tfrac14$, $\iota(0,1)=\tfrac38$.

\begin{lemma}[Dyadic coding \src{05}]\label{onot:neg:lem:dyadic}
$\iota$ is a computable bijection from $\N^{<\N}$ onto the dyadic rationals in
$(0,1]$ which preserves the Kleene--Brouwer order and has a computable inverse.
Hence $\kappa$ is an order-reversing computable bijection onto the dyadic rationals
in $[1,2)$ with computable inverse.
\end{lemma}
\begin{proof}
Every child interval lies strictly left of its parent's right endpoint;
consecutive children occur from left to right with disjoint interiors, and all
descendants stay in their child interval. So extensions precede prefixes and the
first differing child index decides the order, which gives order preservation and
injectivity. For the inverse, write a dyadic $q\in(0,1)$ in lowest terms
$q=c/2^\ell$ with $c$ odd, write $c-1$ as a binary word of exactly $\ell$ bits
with leading zeros (it ends in zero), and parse it uniquely into blocks
$1^{n_0}0\,1^{n_1}0\cdots1^{n_j}0$; the sequence is $(n_0,\ldots,n_j)$, because each
child choice appends exactly the indicated block to the binary address of the
interval. The endpoint $q=1$ is the empty sequence. This gives surjectivity and a
computable inverse, and subtraction from $2$ reverses the order.
\end{proof}

For a decidable tree $T$ put
\begin{equation}\label{onot:neg:eq:treesupport}
 \kappa(T)=\{\kappa(s):s\in T\},\qquad \Sigma(T)=\sum_{s\in T}\omega^{\kappa(s)}.
\end{equation}
This is a \emph{candidate} notation: it is assigned a surreal value only when the
displayed support is admissible. An invalid candidate is not given a default value
such as zero, which would silently change the representation and its decision
problem (05).

\begin{theorem}[Bounded dyadic support validity \src{05}]\label{onot:neg:thm:pionesone}
There is a uniformly effective family of candidates \eqref{onot:neg:eq:treesupport}
in which every support is a decidable subset of the dyadic rationals in $[1,2)$,
every listed coefficient is one, and
\[
 \Sigma(T)\text{ is an admissible omnific normal form}\iff T\text{ is well founded}.
\]
Consequently admissibility of these candidates is $\Pi^1_1$-complete; every
admissible candidate denotes an omnific integer.
\end{theorem}
\begin{proof}
Given a rational, test whether it is dyadic and in $[1,2)$; if so, apply the
inverse of \cref{onot:neg:lem:dyadic} to $2$ minus it and test the resulting
sequence for membership in $T$. This decides the support, uniformly in a total
tree predicate. The decreasing order on $\kappa(T)$ is isomorphic to $(T,\kb)$, so
admissibility is equivalent to well-foundedness by \cref{onot:neg:lem:kbtree}.
Admissible sums have positive exponents and constant zero, so they are omnific.
This reduces well-founded recursive trees to admissibility; conversely the absence
of an infinite increasing sequence of support exponents is a $\Pi^1_1$ statement,
and the tree completeness theorem gives the classification.
\end{proof}

No numerical estimate repairs this validity problem: all exponents lie between one
and two and all coefficients are one. The compactness of a real interval has no
bearing on whether a subset is well ordered; the relevant order is that of the
support, not a metric estimate for a real series (05).

\begin{remark}[Second routes: midpoint codings \src{08, 09}]\label{onot:neg:rem:secondroute}
08 and 09 prove the same completeness with a different coding. Enumerate
$\N^{<\N}$ without repetition; at stage $m$ find the position of the new sequence
among the finitely many already placed in the \emph{reverse} Kleene--Brouwer order,
and assign the midpoint $\mu(s)=(l+u)/2$ of the values $l<u$ of its placed
neighbours, using $1$ and $2$ when a neighbour is absent. Induction on stages shows
that $\mu:\N^{<\N}\to\Q\cap(1,2)$ is a computable injection into the dyadics with
$s\kb u\iff\mu(s)>\mu(u)$ (08's lemma); to compute $\mu(s)$ run the construction to
the stage at which $s$ appears. 09 applies the same midpoint construction to any
computable countable linear order, on all finite sequences first and then
restricted to $T$, so that it is uniform and unaffected by the eventual
well-foundedness of $T$. With $S_T=\sum_{s\in T}\omega^{\mu(s)}$, 08 proves that
validity is $\Pi^1_1$-complete with coefficient one and exponents in the fixed
dyadic set $\mu(\N^{<\N})\subset(1,2)$, and 09 that the property ``this is a valid
Conway normal form'' of a program notation for sums with coefficient one and
distinct rational exponents in $(1,2)$ is $\Pi^1_1$-complete. Neither source shows
that the image of $\mu$, or the support of $S_T$, is decidable, and neither claims
it; decidable supports are obtained here only with 05's coding $\kappa$.
\end{remark}

\begin{proposition}[Promised versus checked equality \src{08}]\label{onot:neg:prop:treeequality}
On admissible tree-support names,
\[
 \Sigma(T)=\Sigma(T')\iff T=T',
\]
so under the validity promise inequality is c.e.\ and equality is co-c.e. The
checked equality relation, requiring validity of both inputs and equality, is
$\Pi^1_1$-complete. The same holds for 08's names $S_T$.
\end{proposition}
\begin{proof}
The coefficient at the distinct exponent $\kappa(s)$ (or $\mu(s)$) is the indicator
of $s\in T$, so normal-form uniqueness gives the equivalence, and a finite node in
exactly one of the trees witnesses inequality. Checked equality is the conjunction
of two $\Pi^1_1$ validity conditions and the arithmetical condition $T=T'$, hence
$\Pi^1_1$; on the diagonal $(T,T)$ it is exactly validity, which gives hardness.
\end{proof}

08's statement reads ``equality under the validity promise has a co-c.e.\ test for
inequality''; its proof, and the summary table of its introduction (``co-c.e.\ on
valid inputs''), show that the intended and correct statement is the one printed
here: inequality is c.e., so promised equality is co-c.e. (\cref{onot:app:corrections}).

\begin{corollary}[No complete enumerable certificate system \src{08}]\label{onot:neg:cor:nocertificates}
No sound computably enumerable family of finite certificates recognizes exactly
the valid tree-support names.
\end{corollary}
\begin{proof}
An accepted finite certificate would make validity c.e., contrary to its
$\Pi^1_1$-completeness.
\end{proof}

\subsection{What these theorems classify}

The theorems classify \emph{validity}, not equality under a validity promise: once a
candidate is promised valid, its equation with itself is trivial, and
$\Pi^1_1$-hardness of promised equality cannot be inferred from the difficulty of
checking the promise; the fixed-support equality theorem was proved by a separate
reduction to avoid exactly this mistake (05). The statement concerns the validity
of this flexible support language, not equality in the fixed-stream language, and
certainly not validity of finite rational--omega terms, which is decidable;
conflating these domains turns correct statements into false ones (09). Nor does it
undermine the certified syntaxes of the positive parts, which certify deliberately
restricted, total families of valid supports and do not purport to accept every
semantically valid program-generated support (08).

The computer-algebra report proves by an arithmetical reduction that Hahn
admissibility of supports output by total bounded-time tests is undecidable
(\replabel{cas:prop-validation}). \Cref{onot:neg:thm:pionesone,onot:neg:prop:treeequality}
strengthen this to an exact $\Pi^1_1$ classification with bounded dyadic
exponents, and for checked equality. The computable-surreals report proves
$\Pi^1_1$-completeness of \emph{structural} validity for its cut-program coding
(\replabel{prop:validity}) by the same tree reduction in a different
representation.

\part*{Implementations, relations, questions, and records}
\addcontentsline{toc}{part}{Implementations, relations, questions, and records}

\section{Implementation specifications and the shipped programs}\label{onot:sec:implementation}

\subsection{Interfaces}

A practical implementation should keep separate interfaces (05). A
\emph{hereditary normal form} carries finite exact coefficient and exponent data.
A \emph{certified compressed form} carries finite fraction, holonomic or matrix
data together with its denominator, scale, jet and shielding contracts. A
\emph{stream form} supplies coefficient access on a stated support but need not
admit an equality decision. These are not names for one unchecked data structure:
an operation converting a stream into a compressed form is partial and
certificate-producing, and it cannot be a total hidden coercion by
\cref{onot:neg:thm:rationalhard,onot:neg:thm:statebound,onot:neg:thm:notranslator}.

For the towers (08), an exact coefficient backend supplies $k$ with decidable
validity, equality, order and arithmetic; a scale context stores a sorted finite
list of canonical ordinal labels and represents each lattice exponent by its
integer coordinate vector; and a recursive field name is either a $k$-name or a
fraction of certified Euler-series names whose operator and jet coefficients are
field names in the preceding context. An omnific name may use the triangular form
\eqref{onot:hol:eq:triangular} directly: an integer constant, positive outer
powers, and arbitrary valid field names one level lower as coefficients, which
makes omnific membership automatic. An importer of an unrestricted field name
must instead run the recognition algorithm of \cref{onot:hol:thm:triangular}; a
nonzero field name provides an exact Laurent valuation, zero being handled first.

\begin{lstlisting}[caption={Field-to-integer conversion in the towers (08).},label={onot:lst:normalize}]
OmnificNormalize(x, level):
    if level == 0:
        m := ordinary_floor(x)
        return Integer(m) if x == m else Reject
    if x == 0:
        return Zero
    N := LaurentValuation(x, top_variable)
    prefix := sum(Coefficient(x,n) * top_variable^n
                  for n from N through 0)
    # The sum is empty when N > 0.
    if x != prefix:
        return Reject
    lower := OmnificNormalize(Coefficient(x,0), level-1)
    if lower is Reject:
        return Reject
    return Triangular(lower, all nonzero coefficients at N <= n < 0)
\end{lstlisting}
Every equality sign invokes the certified field equality test; replacing it by a
numerical tolerance or by comparison of a fixed number of coefficients would
invalidate the proof (08).

\begin{lstlisting}[caption={Hereditary finite normalization (05).},label={onot:lst:hereditary}]
normalize(expression):
    recursively normalize the operands
    for addition: merge equal exponents and add coefficients
    for multiplication:
        for each pair of terms:
            add their exponents recursively
            multiply their coefficients
        merge the resulting finite list
    discard zero coefficients
    absorb exponent zero into the constant
    sort exponents by recursive exact comparison
    return the canonical tree
\end{lstlisting}

\begin{lstlisting}[caption={Equality in the $\varepsilon_0$-shielded calculus (05).},label={onot:lst:epsilon}]
equal(A, B):
    validate the hereditary integer parts
    validate every fraction denominator as nonzero
    subtract A and B as finite polynomials in epsilon_0
    for each coefficient p/q - r/s:
        if normalize(p*s - r*q) is nonzero:
            return False
    return True
\end{lstlisting}

\begin{lstlisting}[caption={Collision-aware leading term of a matrix series (05).},label={onot:lst:front}]
front(u, M[1..r], v, delta[1..r]):
    require every delta[i] > 0
    P = u * product(adjugate(I - z[i]*M[i])) * v
    weights = empty exact map
    for each monomial c*z^n of P:
        weights[sum(n[i]*delta[i])] += c
    remove zero coefficients from weights
    if weights is empty:
        return Zero
    lambda = least key in weights
    return LeadingTerm(lambda, weights[lambda])
\end{lstlisting}
The last algorithm is finite even if the exponent interval below its answer
contains infinitely many potential support points. For the omnific value $X\Phi$
it returns the omega-leading exponent $\Lambda-v(\Phi)$ and the same coefficient;
for $X^j\Phi$, use $j\Lambda-v(\Phi)$ (05).

\subsection{Finite certificates}

At one level of a tower, a zero answer for a certified series consists of its
normalized operator, the verified exceptional-index computation and the compatible
zero jet; a nonzero answer identifies the first nonzero jet coefficient, whose
nonzeroness is certified recursively; fractions use a certified numerator
cross-difference; exact rational or algebraic arithmetic discharges the claims at
the bottom. The integer-root procedure exposes a certificate---the leading
polynomial $p$, an integer Cauchy bound $M$ and the exact tests of $P(n)$ for
$0\le n\le M$---which an independent checker can verify without trusting a
black-box root solver; a faster implementation may factor $p$ first. These are
mathematical specifications: the prototype uses exact arithmetic but does not emit
a proof-assistant object or a complete independent proof trace (08). For the
hereditary layer, equality certificates are normalized finite trees with recursive
checks of exponent order and coefficient identities; for $\FH$, the identity
$ps-rq=0$ and the two nonzero denominators; for a shielded polynomial, the finite
coefficient lists; for a matrix series, the adjugate identities, the collected
numerator and the unit leading term of the denominator. For a general stream,
inequality has a finite first-difference witness; equality may sometimes have a
finite special-purpose proof, such as a recurrence identity, but no sound c.e.\
collection of such proofs covers every true equation of the switch family
(\cref{onot:neg:cor:proofobstruction}), which limits completeness, not usefulness
(05).

\subsection{Complexity and resource limits}

No uniform polynomial-time bound is claimed for any of the three systems.
Differential orders grow under tensor products and coefficient expressions under
repeated field operations; a singularity $M$ written in binary forces a jet of
length $M+1$ in the elementary certificate format (\cref{onot:hol:ex:resonance}). A
production system could use sparse compatible jets with a proved completion
procedure, shared subexpressions, cached zero tests and specialized annihilator
algorithms, each with its own correctness argument (08). The hereditary normalizer
and the grid compiler allow substantial expression swell (05, 09). A
resource-limited implementation should return a distinct ``not completed'' status
when a limit is exceeded; that is not a proof of inequality, invalidity or
nonexistence, and failing to import a raw stream into the holonomic backend says
nothing about whether its value has a holonomic or even rational name
(\cref{onot:neg:thm:notranslator}) (08).

\subsection{The shipped programs}

Delivered code is kept byte-identical under its placed names; the README of this
directory explains how to rerun it on a copy.

\paragraph{Source 08.} \fn{code/08-holonomic-towers-holonomic.py} implements
certified Euler series over $\Q$, sum and product annihilators by finite
differential-system linear algebra, their fraction field, and the guarded ring
$\Z+UF_1[U]\subset R_2$; \fn{code/08-holonomic-towers-ordinals.py} implements
hereditary Cantor-normal-form trees below $\varepsilon_0$ with comparison, natural
and ordinary addition, natural multiplication and omega exponentiation;
\fn{code/08-holonomic-towers-verify.py} runs the tests and
\fn{code/08-holonomic-towers-demo.py} is a short example. Its scope is strictly
smaller than the tower theorem: arbitrary coefficient backends, general
higher-rank operators, the full $R_2$ recognition and floor algorithms, and all
$\varepsilon_0$-labelled embeddings are mathematical constructions, not claimed
implementations. The recorded run (\fn{data/08-holonomic-towers-verification.json},
Python 3.13.5, SymPy 1.14.0) passed 41 test groups, including 521 finite
change-event schedules, 780 arbitrary-sign block assignments, and all 14{,}520
ordered pairs among 121 dyadic-coded tree nodes; the tests exercise recurrence
indices, exceptional jets, series identities through certificates, Laurent
leading data and finite instances of the dominance encodings. Tests of
infinite-series identities use generated certificates and the finite-jet rule,
not matching a fixed number of terms.

\paragraph{Source 05.} \fn{code/05-hereditary-shielding-verify.py} (Python 3.9 or
later, standard library only, exact rational arithmetic) checks hereditary
arithmetic and the floor correction, the dyadic coding and its order
preservation, matrix numerator identities and collision examples, and Hankel ranks
of switch sequences. The recorded run
(\fn{data/05-hereditary-shielding-verification.json}) passed 18{,}109 assertions
in 16 categories, the largest being 14{,}641 Kleene--Brouwer comparisons; these
are not 18{,}109 independent theorems.

\paragraph{Source 09.} \fn{code/09-rational-omega-omega\_notation.py} is an exact
prototype: its sparse-fraction kernel uses Python integers and
\texttt{fractions.Fraction}; SymPy is used only after grid compilation, for
rational-function algebra and polynomial division; floating-point input is
rejected. \texttt{om(x)} is $\Om(x)$; \texttt{equivalent} and \texttt{compare}
implement \cref{onot:grid:thm:equality}, \texttt{leading} the leading data,
\texttt{separated\_basis} and \texttt{grid\_form}
\cref{onot:grid:lem:basis,onot:grid:thm:grid}, and \texttt{nonnegative\_part},
\texttt{is\_omnific}, \texttt{omnific\_floor}, \texttt{truncate\_at} and
\texttt{coefficient\_at} \cref{onot:grid:thm:effectivefloor}. Semantic equality is
\texttt{equivalent}, not Python equality. For example, with
\texttt{w = om(1)} and \texttt{B = om(w) / (1 - 1/w)}, the prototype confirms that
$\sB$ is omnific, that $\sB+\tfrac12$ is not, that $\fl{\sB-\omega^{-1}}=\sB-1$,
that $[\omega^{\omega-3}]\sB=1$ and that
$\tau_{\ge\omega-3/2}\sB=\omega^\omega+\omega^{\omega-1}$. The recorded run
(\fn{data/09-rational-omega-verification.json}, Python 3.13.5, SymPy 1.14.0, seed
23092026) passed 197 exact assertions, covering nested monomial identities, a
positive infinitesimal exponent, infinite geometric omnific values, both
floor-boundary cases, absent coefficients, off-lattice cutoffs, a tail-sensitive
separated basis, polynomial-part fixtures in two and three scales, and rejection of
zero division and floating-point input; several expected truncations are supplied
as independent fixtures.

\paragraph{Reruns for this report.} All three suites were rerun on copies with the
delivery file names restored, under Python 3.14.4 and SymPy 1.14.0, and passed
with the same counts (41 groups; 18{,}109 assertions; 197 assertions). No test
computes a halting set, verifies a $\Pi^1_1$ reduction, enumerates the support of
$\mathsf E$, or proves a completeness, support or termination theorem.

\subsection{Formalization targets}

No part of this report is formalized, and \fn{docs/FORMALIZATION.md} maps none of
its labels. The sources propose the same layering. A manageable first target is
the Euler recurrence theorem over an abstract effective field, then the
integer-root lemma and the triangular theorem over iterated Laurent fields, with
the bridge to actual surreal monomials kept as a separate semantic layer (08).
05 suggests finite hereditary trees with rational coefficients first, then the
shielding theorem over an explicitly embedded set-sized Hahn field, then the
adjugate certificate theorem, which needs only finite algebra and a positive-scale
evaluation interface, and finally the computability reductions in a fixed
machine-index framework. 09 suggests an inductive syntax with finite sparse
fractions and explicit depth measures, a denotation into the surreal carrier or a
fixed Hahn workspace, and proof obligations for the collection of equal exponents,
the fraction leading-data formula, leading-scale elimination, the ordered
embedding of each lattice, and the polynomial-part identities, with the
negative-infinitesimal floor correction as a distinct branch. A type-level
assertion in Lean that equality is decidable is not by itself an executable
verified normalizer; classical instances must be distinguished from the
algorithms, and a classical formalized isomorphism, especially one using choice,
does not implement a computable normalizer for finite strings (05, 09). The
normal-form bridge note \fn{docs/NORMAL_FORM_BRIDGE.md} describes the existing
finite-evaluation and field-bridge modules.

\section{Relations to neighbouring reports}\label{onot:sec:neighbours}

\begin{itemize}[leftmargin=*]
\item \fn{docs/foundations-and-computation/computer-algebra} (\texttt{cas:}). The
two stream reductions are its \replabel{cas:thm-zerotest} with
\replabel{cas:rem-moralA} and \replabel{cas:rem-moralB}; 08's two-block sign
example is \replabel{cas:prop-leading} and the one-block contrast
\replabel{cas:rem-rankone}; \cref{onot:neg:thm:ershov} strengthens the former, and
\cref{onot:neg:thm:pionesone,onot:neg:prop:treeequality} strengthen
\replabel{cas:prop-validation}. The floor is \replabel{cas:eq-floor}. After
compilation the rational--omega grids are its core \replabel{cas:thm-core}; its
positive-grid extraction \replabel{cas:thm-grid} concerns coefficient extraction
with collisions, to which 05's collision-aware front adds a finite zero and
leading-term certificate for matrix coefficient rules. Its section on positive
counterweights says that singular points require a certified recurrence start and
treatment of exceptional indices; \cref{onot:hol:thm:integerroots,onot:hol:thm:certificate}
supply exactly that, over non-Archimedean coefficient fields. Its remark on the
2026 zero test (\replabel{cas:rem-cfh}) applies equally to the towers.
\item \fn{docs/foundations-and-computation/computable-surreals} (bare labels). The
one-coefficient obstruction uses the constants of \replabel{thm:inversehard}(iv);
its structural validity theorem \replabel{prop:validity} is the tree reduction in
another representation; its no-go theorem \replabel{thm:nogo} lists a more
informative real representation among the ways out, the route taken by all three
systems here. Its research directions ``support certificates beyond
left-finiteness'' and ``uniform translation complexity'' are \emph{not} answered:
the transfinite supports here arise only inside finite-rank lattices or towers,
which is the excluded case, with exact rather than Cauchy coefficients
(\cref{onot:grid:thm:escape} and the questions below only point beyond it), and
the translator theorems concern stream-to-finite names, not numerical-to-structural
names.
\item \fn{docs/surreal/omnific-diophantine-geometry} (\texttt{odg:}). The floor is
\replabel{odg:thm:floor}, division with remainder \replabel{odg:prop:division}, the
Diophantine transfer \replabel{odg:cor:H10} with \replabel{odg:eq:existencetransfer};
all are cited, with only their effective forms proved here. Its question
\replabel{odg:q:roots} on root decidability for finitely represented omnific
coefficients is \emph{not} addressed: the systems here decide ground equality,
order and floor, not the existence of roots; nor is \replabel{odg:q:search}. The
report is being extended concurrently from the same batch; only labels present at
the time of writing are cited.
\item \fn{docs/surcomplex/holonomic-rigidity-for-entire-hahn-functions}
(\texttt{hol:}). \Cref{onot:hol:thm:integerroots} is the effective form of its
\replabel{hol:lem:integer} (\cref{onot:hol:rem:holinteger}). Its question
\replabel{hol:nl:q:effective}, on computing candidate degree sets and optimal
excluded valuation cuts for nonlinear rigidity certificates, is \emph{not}
answered; the effective leading-data fields of \cref{onot:hol:part} are relevant
input only.
\item \fn{docs/foundations-and-computation/foundations} (\texttt{found:}): the
definition $\Oz=\Pi\oplus\Z$ (\replabel{found:eq:omnific}) agrees with
\cref{onot:def:oz}.
\item \fn{docs/foundations-and-computation/definable-surreals-and-omnific-integers},
written concurrently from the same batch, restates the same floor; its subject is
set-theoretic definability, not computability, and it explicitly declines to
identify the two. It is cited by directory only.
\end{itemize}

\section{Questions}\label{onot:sec:questions}

The sources call these directions suggested by their results, not completed results
or previously published open problems (08, 05, 09).

\begin{question}[Compact certificates \src{08, 05}]\label{onot:q:compact}
Can the full jet through the last exceptional index be replaced by a compressed
certificate listing the free singular coefficients with a small independently
checkable proof that the coefficients between them follow from the recurrence
(08)? Can the adjugate numerator be computed or certified by reachability and
observability reduction, structured polynomial identity testing or sparse
group-algebra arithmetic with substantially smaller certificates, and what is the
bit complexity when the number of scales is part of the input (05)?
\end{question}
\emph{Status: open.} Univariate polynomial-time weighted-automata algorithms should
not be quoted as a bound for every multiscale representation (05).

\begin{question}[Larger certified coefficient classes \src{08, 05, 09}]\label{onot:q:coefficients}
Can the D-finite numerator and denominator class be replaced by richer uniquely
specified differential-algebraic classes, with a precise common field, leading-data
access, truncation closure and compatible parameter semantics (08)? Can algebraic
extensions of $\FH$ with explicit branch and order data serve as coefficients for
the shielding theorem (05)? For algebraic real constants in rational--omega
expressions, is there a block version of leading-scale elimination over effective
number fields that computes the finite subgroup order and truncation cuts uniformly
(09)?
\end{question}
\emph{Status: open.} The D-algebraic zero test \cite{CFH} gives tools but the
obligations must be verified separately rather than inferred from the phrase
``zero-test'' (08); the separated rational basis fails with irrational coefficients
(\cref{onot:grid:lem:basis}).

\begin{question}[Beyond finite rank \src{05, 09}]\label{onot:q:beyond}
Is there a transfinite iteration constructor, or a certified $\omega$-summation
operation with syntactic support certificates and algebraic restrictions on block
coefficients, that extends the one-state supports $\omega^r$ or the escape sum
$\mathsf E$ while keeping equality decidable?
\end{question}
\emph{Status: open.} \Cref{onot:grid:prop:CK,onot:neg:thm:pionesone} show that no
sound effective universal constructor covers all recursive support orders, and
\cref{onot:neg:thm:notranslator} excludes all primitive-recursive coefficient
programs; a specific stronger language needs its own positive theorem (05, 09).
This is also the direction ``support certificates beyond left-finiteness'' of the
computable-surreals report.

\begin{question}[Beyond finitely many dominance priorities \src{08}]\label{onot:q:priorities}
Can higher effective difference hierarchies be assigned to appropriately
well-founded systems of dominance priorities?
\end{question}
\emph{Status: open.} The finite proof uses a computable uniform bound on the number
of record priorities; an ordinal-indexed version must specify a validity
certificate and an effective descending control measure, and the finite
classification is not asserted to extend automatically (08).

\begin{question}[Intrinsic versus presentation-based invariants \src{08}]\label{onot:q:intrinsic}
Which representation morphisms preserve leading-term certificates or exact support
counts, and which natural subclasses of raw names admit effective translations?
\end{question}
\emph{Status: open.} \Cref{onot:neg:thm:notranslator} shows that an extensional
inclusion of denotation sets is much too weak to answer this (08).

\begin{question}[Richer exponent languages and proof-theoretic strength \src{05}]\label{onot:q:exponents}
Which additional decision assumptions are needed for exponent languages admitting
transcendental real constants, global exponential identities or cut-defined
exponents? What proof-theoretic strength is needed to establish termination of
normalizers for richer ordinal-indexed constructors?
\end{question}
\emph{Status: open.} Any extension must state its decision assumptions rather than
subsume them under a generic ``simplify'' operation; the extensional question of
which surreals are denoted stays distinct (05).

\begin{question}[Size, sharing, and mechanization \src{09}]\label{onot:q:complexity}
What complexity bounds hold for the recursive comparator and the grid compiler in
terms of nesting depth, the number of distinct exponent subterms, grid rank,
degrees and coefficient bit lengths? Can the ring $\IO$, whose ground word problem
is decidable and whose existential polynomial solvability is not, serve as a
testbed for mechanizing omnific division?
\end{question}
\emph{Status: open.} Finite test timings do not justify an asymptotic claim (09).

\begin{question}[Comparison of the three systems \mergetag]\label{onot:q:compare}
For $k=\Q$, is the tower ring $\mathcal R_\varepsilon$ contained in the
rational--omega ring $\IO$? In particular, does $\sum_{n\ge0}n!\,\omega^{\omega-n}$
have a rational--omega term, or lie in $\FH$?
\end{question}
\emph{Status: open;} no source addresses it (\cref{onot:grid:prop:incomparable}).

None of these is a question named by an earlier report, and no question of an
earlier report is answered here: not the research directions of the
computable-surreals report, not \replabel{odg:q:roots} or \replabel{odg:q:search},
and not \replabel{hol:nl:q:effective}.

\appendix
\section{Provenance}\label{onot:app:provenance}

\subsection{Section by section}

{\renewcommand{\arraystretch}{1.15}
\begin{longtable}{>{\raggedright\arraybackslash}p{0.27\textwidth}>{\raggedright\arraybackslash}p{0.63\textwidth}}
\caption{Provenance of each part of the report.}\label{onot:tab:provenance}\\
\toprule
Part of this report & Sources and choices\\
\midrule
\endfirsthead
\toprule
Part of this report & Sources and choices\\
\midrule
\endhead
\Cref{onot:sec:intro} & 08 \S1 with 05 \S1 and 09 \S1; the summary table combines 08's and 05's tables.\\
\Cref{onot:sec:normal} & 08 \S\S2--3, 05 \S2, 09 \S2, printed once. The floor formula is cited from \replabel{odg:thm:floor} and \replabel{cas:eq-floor}; 05's Theorem~3.2, 08's Theorem~7.4 and 09's equation (8) prove it again. Canonicalization: 08's Proposition~2.3, 05 \S2.2, 09's Corollary~3.3. Size obstruction: all three.\\
\Cref{onot:sec:shield} & 05 \S4 (Theorems 4.1, 4.2), printed as the general embedding; \cref{onot:rem:characterizations} is the merge's.\\
\Cref{onot:hol:part} & 08 \S\S4--10. \Cref{onot:hol:rem:holinteger} (credit to \replabel{hol:lem:integer}) is the merge's; the floor theorem keeps only its effectivity.\\
\Cref{onot:sh:part} & 05 \S\S3, 5, 6. The floor theorem keeps only its effectivity; the Sch\"utzenberger attribution is added.\\
\Cref{onot:grid:sec:equality,onot:grid:sec:grid,onot:grid:sec:truncation} & 09 \S\S3--5; the division corollary cites \replabel{odg:prop:division} for existence and uniqueness.\\
\Cref{onot:grid:sec:ordinals} & 09 \S6 as the main statement; \cref{onot:grid:prop:incomparable} is the merge's; the $\CK$ proposition is 09's Proposition~6.3 with 05's Proposition~10.1 as its uniform half.\\
\Cref{onot:grid:sec:support,onot:grid:sec:certificates} & 09 \S\S7, 9; the Diophantine transfer cites \replabel{odg:cor:H10}.\\
\Cref{onot:neg:sec:raw} & 08 \S11 (base; $d$ blocks) with 05 \S7 and 09 \S\S8.1--8.3 as the case $d=1$; strict comparison from 05's Theorem~7.1, non-semidecidability from 09's Theorem~8.2; the no-translator theorem is 08's Theorem~11.4 and 09's Theorem~8.3 printed once; credit to \replabel{cas:thm-zerotest}.\\
\Cref{onot:neg:sec:rational} & 05 \S8; credit to \replabel{cas:rem-moralA}.\\
\Cref{onot:neg:sec:ershov} & 08 \S12; credit to \replabel{cas:prop-leading}, \replabel{cas:rem-rankone}.\\
\Cref{onot:neg:sec:validity} & 05 \S9 (coding and theorem, main route); 08 \S13 (Kleene--Brouwer lemma, second route, promised and checked equality, corrected) and 09's Theorem~8.5 (second route).\\
\Cref{onot:sec:implementation} & 08 \S14, 05 \S11, 09 \S10.\\
\Cref{onot:sec:neighbours,onot:sec:questions} & The merge, with the sources' directions (08 \S16, 05 \S12, 09 \S11.3).\\
\bottomrule
\end{longtable}}

Section, theorem and equation numbers of the sources refer to their own
manuscripts, which are not shipped. Where the sources overlap, the statement with
the weakest hypotheses is printed. No source proof needed a hypothesis that the
printed statement lacks: the sum-of-squares step of the transfer theorem needs an
ordered coefficient field, which the Hahn setting supplies; 08's floor proof uses
only \cref{onot:hol:cor:oztest} and the rank-zero integer search; the stream
reductions use only $\Q$-valued primitive-recursive coefficients.

\subsection{Dependency map}

For the towers the dependencies run as follows (08): exact ordered real
coefficients; leading-polynomial integer-root bounds; finite Euler certificates;
effective D-finite ring operations; fraction fields and leading data; finite-rank
recursion; truncation and triangular integer parts; floor and ordinal recognition;
effective insertion of finitely many labels. The next rank is never used to justify
the preceding rank's zero test, and the root procedure needs only real leading
coefficients. For rational--omega expressions (09): finite recursive equality;
separated grid; polynomial-part recursion; omnific recognition and floor; the
support barrier follows from finite lattice rank, not from the floor theorem. For
05: hereditary normalization and ordinal bounds; shielding; fraction supports and
the $\varepsilon_0$ calculus; matrix adjugates and positive-scale evaluation. The
negative results are independent of all three constructions: they use normal-form
uniqueness, fixed positive support blocks, bounded machine simulation, and, only
for validity, the completeness of well-founded recursive trees. The translator
theorems join the tracks by observing that every spike value has a rational
monomial name.

\subsection{Renamed symbols}

{\renewcommand{\arraystretch}{1.15}
\begin{longtable}{>{\raggedright\arraybackslash}p{0.12\textwidth}>{\raggedright\arraybackslash}p{0.32\textwidth}>{\raggedright\arraybackslash}p{0.44\textwidth}}
\caption{Symbols renamed from the sources.}\label{onot:tab:renames}\\
\toprule
Source & Source symbol & Here, and why\\
\midrule
\endfirsthead
\toprule
Source & Source symbol & Here, and why\\
\midrule
\endhead
08 & $\ell(x)$ (leading exponent) & $\ddeg x$, the collection's symbol; 09's $\ell$ is support length.\\
08 & $C$ (base field) & $k$; 05's $\mathcal C$ and 09's $C$ are other objects.\\
08 & $D$ (code set) & $N$; $D$ is the constant ring of the shielding theorem.\\
08 & $N(a)$ (canonicalization) & $\can(a)$; $N$ is the code set.\\
08 & $E^{\mathrm{chk}}$ & $\Eq^{\mathrm{chk}}_\nu$.\\
08 & $L$ (Euler operator) & $\mathsf L$; $L$ is the shielded coefficient field.\\
08 & $\mathcal A(K,t)$ (modular class) & $\mathfrak C(K,t)$.\\
08 & $\operatorname{Tr}_{\ge g},\operatorname{Tr}_{>g}$ & $\tr g,\trs g$, shared with 09.\\
08 & $G_r=\prod(1-t_j)^{-1}$, $G(t)$ & $\Geo_r$, $\Geo(t)$; $G_r$ is the exponent group.\\
08 & $A(t)$, $E(t)$, $E_-(t)$ & $\Fac(t)$, $\Ex(t)$, $\Ex(-t)$.\\
08 & $B$ in $A\subseteq B$ (label sets) & $A'$; $B(P)$ is the jet bound.\\
08 & $S_e$, $q$ (tree sums, midpoint map) & $S_T$, $\mu$; $\Sigma(T)$ and $\kappa$ are 05's.\\
08 & $k$ (block index), $\ell$ (product index) & $j$, $l$; $k$ is the coefficient field.\\
05 & $K$ (coefficient field), $\mathbb A$ & $k$, $\Ralg$.\\
05 & $E$ ($=\varepsilon_0$) & $\varepsilon_0$; $\Ex$ is the exponential series.\\
05 & $\mathcal F=\Frac\Hh$, $\mathcal C$ & $\FH$, $\SE=\mathcal S_{\varepsilon_0}(\Dd,\FH)$; $\mathcal F_A$ and $\QO$ are the other fields.\\
05 & $F$, $D$ in the shielding theorem & $L$, $D$; $\Eq_F$ is $\Eq_L$.\\
05 & $C(p)$ (canonicalizer) & $\can(p)$.\\
05 & $P+c+\eta$ (floor decomposition) & $x_{>0}+\ct(x)+x_{<0}$.\\
05 & $F$, $\widehat F$, $P$, $Q$ (matrix series) & $\Phi$, $\widehat\Phi$, $\mathsf P$, $\mathsf Q$.\\
05 & $d$ (matrix dimension), $B=d+e$ & $m$, $b=m+m'$; $d$ is the number of blocks.\\
05 & $\lambda$ (leading Hahn exponent) & $v(\Phi)$, a valuation in the $t$-convention.\\
05 & $N$ (nilpotent matrix) & $N_0$; $N$ is the code set.\\
05 & $Z_r$ & $\zeta_r$; $Z$ denotes raw block names.\\
05 & $X=\omega^\omega$ (streams), $\mathsf S(a)$ & $U$, $Z(a)$, as in 08.\\
05 & $h_e$ (halted within $n$ steps) & $\hb_e$; $h_e$ is the first-halt spike of 08 and 09.\\
05 & $A_e$, $B$ (switch family) & $A^{\mathrm{sw}}_e$, $\sB$.\\
05 & $q$, $e(s)$, $(\ell_s,r_s)$, $\mathsf Z(T)$, $S_T$ & $\iota$, $\kappa(s)$, $(a_s,b_s)$, $\Sigma(T)$, $\kappa(T)$; $e$ is a machine index.\\
09 & $\mathcal F$, $\mathcal F_d$, $\mathcal A$ & $\QO$, $\QO_d$, $\IO$; $\mathcal F_A$ is 08's field.\\
09 & $\lambda(x)$, $\ell(x)$ & $\ddeg x$, $\ord(\supp x)$.\\
09 & $K_r$, $A_r$, $A^{\mathrm{tail}}_{r-1}$ & $\KG_r$, $\IO_r$, $\IO^{\mathrm{tail}}_{r-1}$; $K$ is 08's leading-data field.\\
09 & $T_r(R)$, $C_r(R)$ & $\tr0R$, $\ct(R)$.\\
09 & $D$ (code set; $\sum\omega^{1/n}$) & $N$; $\mathsf H$.\\
09 & $B$, $B_d$, $C_d$, $C$ & $\sB$, $\sB^\sharp_d$, $\mathsf E_d$, $\mathsf E$.\\
09 & $D_i$ (jet degree bounds), $G$ (support group) & $b_i$, $\Gamma$.\\
09 & $t$ (terms), $S(a)$, $Z_i$ & $s$, $Z(a)$, $Y_i$.\\
\bottomrule
\end{longtable}}

\subsection{Stale repository statements}\label{onot:app:stale}

All three sources pin \texttt{a45d722}, 30 commits before the placement. At the pin
the computer-algebra, computable-surreals, omnific Diophantine and
holonomic-rigidity reports already contained every label cited here. The sources'
reviews were targeted, and the following statements are corrected; the pins are
kept as provenance.
\begin{enumerate}[leftmargin=*]
\item 08 read the computer-algebra README (lines 1--180) and describes the report
correctly as containing effective rational monomial cores and zero-test and
high-rank leading-term obstructions. Its novelty ledger then presents the
localization of raw equality to at most one monomial as part of its contribution,
and prints the two-block sign example; both were at the pin
\replabel{cas:rem-moralB} and \replabel{cas:prop-leading}, and are credited here
(\cref{onot:neg:sec:raw,onot:neg:sec:ershov}).
\item 05 read the omnific Diophantine report only through its line 150 and says
that it ``uses normal-form constant terms to organize the ring''. That is correct,
but at the pin the report already contained the floor, division and
Hilbert's-tenth-problem results (\replabel{odg:thm:floor},
\replabel{odg:prop:division}, \replabel{odg:cor:H10}), which 05's and 09's
corresponding results duplicate without citing.
\item 09 writes that ``similar constant-term arguments already belong to the
repository's omnific arithmetic program'', citing the catalogue; the precise
reference is \replabel{odg:cor:H10} with \replabel{odg:eq:existencetransfer}.
09's description of the computable-surreals report as ``the full 51-page report''
still matches that report's README.
\item Neither 05 nor 09 cites the computer-algebra report, whose Reductions~A and~B
anticipate their stream reductions, and no source cites
\replabel{hol:lem:integer}; the credits are added.
\item At the pin the computer-algebra report said that an omnific-integer
convention ``is a different target again'' from its $\Pi\oplus\Z$ floor. This was
corrected in commit \texttt{1e3f14b} before the placement: the report now
identifies $\Pi\oplus\Z$ with $\Oz$ (\replabel{found:eq:omnific}).
\item All three sources say that the repository has no Lean formalization of
omnific integers; this remains true, and \fn{docs/FORMALIZATION.md} maps no label of
this report.
\item 08 cites the Chen--Fang--van der Hoeven zero test as a 2026 arXiv preprint;
the computer-algebra report records its publication in the ISSAC 2026 proceedings
\cite{CFH}.
\end{enumerate}

\subsection{Corrections}\label{onot:app:corrections}

\begin{enumerate}[leftmargin=*]
\item 08's Proposition~13.4 stated that ``equality under the validity promise has a
co-c.e.\ test for inequality''. Its proof (a finite node in exactly one tree
witnesses inequality) and its own summary table (``co-c.e.\ on valid inputs'') show
that inequality is c.e.\ and promised equality co-c.e.; the statement is corrected
(\cref{onot:neg:prop:treeequality}). Wording only.
\item 05 attributes its matrix zero test to weighted-automata expositions only; the
underlying rational-series theory is Sch\"utzenberger's, which is now named
(\cref{onot:sh:sec:matrices}).
\end{enumerate}
No theorem of the three sources was found to be false. The placement dossier
re-read every proof and independently rechecked 30 finite instances (among them
05's collision example, the Hankel ranks $s+1$ for $s\le12$, the dyadic coding on
all sequences of length at most three over $\{0,\ldots,3\}$, 08's certificates and
integer-root example, and 09's division example).

\section{What this report does not claim}\label{onot:app:nonclaims}

Every limitation, disclaimer and priority caveat stated by a source is kept here,
numbered per source, followed by those of the merge.

\subsection{Source 08 (27 items)}
\begin{enumerate}[leftmargin=*,widest=00,label=08.\arabic*]
\item The proofs are mathematical arguments, not Lean formalizations; no Lean file
is supplied, no theorem is asserted to have passed the repository's axiom audit,
and classical normal-form semantics is used directly, not a claim that every
semantic bridge in the repository is formalized.
\item Several combined results are proposed as original formulations; the
literature search is not a proof of priority and supports only ``to our
knowledge''; it cannot certify absence from published literature, unpublished work
or uninspected repository companions.
\item Conway normal forms, the omnific support criterion and Hahn arithmetic are
classical; sparse exact monomial algebra and representation-sensitive obstructions
are classical in substance and already in the repository.
\item D-finite ring closure and linear-equation zero tests are classical; the
construction is not presented as the first zero test beyond rational functions,
and more general transseries zero tests are acknowledged.
\item The integer-root procedure is an explicit elementary construction with
standard ingredients; no broad priority claim is made for a root-bound
observation.
\item Raw equality hardness is a standard halting method, localized; the Ershov
hierarchy and the reduction techniques are established; the $\Pi^1_1$ result is a
localized formulation, not a new discovery of the completeness of
well-foundedness.
\item The differential operator at a level treats the lower field as constant; no
compatibility with a global surreal derivation is assumed.
\item No algorithm is asserted for recognizing that a coefficient program is
D-finite or for finding its equation, and a general routine deciding whether a raw
program admits a certificate is neither supplied nor implied.
\item No analytic convergence is claimed, and no real closedness of the tower
fields.
\item The modular generalization does not justify inserting arbitrary
differential-equation solvers; a zero test, a solution convention and terminating
coefficient algorithms remain hypotheses.
\item The triangular form is not a claim that every coefficient has a finite
Conway normal form; a truncation is returned as an expression, not a finite list.
\item Ordinal recognition does not decide nonnegativity of arbitrary recurrences.
\item $\varepsilon_0$ is a design choice, not a ceiling for effective omnific
notation; ``least ordinal upper bound'' is not ``surreal supremum''.
\item The construction is not a notation system for every omnific integer: it
needs an exact coefficient field, finitely many scales per name and a specified
certificate class, and it is not closed under all analytic functions or cuts.
\item Support bounds concern rank, not birthdays; the $\varepsilon_0$ trace
concerns the chosen label syntax.
\item The $d$-c.e.\ classification concerns zero-constant raw $d$-block names with
primitive-recursive rational coefficients and finitely many blocks; the $\Pi^1_1$
result concerns a different raw syntax; moving any conclusion to another
representation needs a new proof.
\item The count-promise algorithm does not certify the count and can diverge or be
wrong on a false count.
\item No uniform complexity bound; jets can have length $M+1$.
\item A ``not completed'' status is not a proof, and failing to import a raw stream
says nothing about its value.
\item The prototype lacks arbitrary coefficient backends, higher-rank operators, the
full $R_2$ recognition and floor, and all $\varepsilon_0$-labelled embeddings; these
are mathematical algorithms, not software capabilities.
\item The code is not formally verified, is not a hardened parser for untrusted
input, and is not a security boundary (callers can mutate Python objects); it emits
no proof-assistant object or complete proof trace.
\item No finite test establishes a halting reduction, a completeness theorem, a
support theorem for arbitrary trees, or correctness on every input.
\item During development a test bug (reversed neighbour selection in the dyadic
audit) was found and fixed; the record reports only the final run, not that no
intermediate test failed; the lemma itself did not rely on the erroneous code.
\item The repository inspection was targeted (the computable-surreals article
lines 1--230 and the computer-algebra README lines 1--180, with the catalogue); the
repository's formalization status is not transferred; nothing was modified.
\item The ordinal-notation overview is motivation, not a technical source.
\item The further questions are not asserted results; the finite mind-change
classification is not asserted to extend to ordinal priorities.
\item The PDF checks were document-quality checks, not proof of correctness.
\end{enumerate}

\subsection{Source 05 (26 items)}
\begin{enumerate}[leftmargin=*,widest=00,label=05.\arabic*]
\item An AI-assisted research manuscript, not refereed; no theorem has been
independently refereed or newly formalized in Lean, and no repository-wide Lean
build was run.
\item The formulations appear new after a limited search; they are contribution
candidates, not a literature-wide priority claim; some narrowly quoted queries
returned nothing, which is weak evidence.
\item The audit was targeted (catalogue, the computable-surreals README and article
lines 1--295, the omnific Diophantine article lines 1--150); nothing is claimed
absent from uninspected parts, and no repository theorem is assumed formally
verified because a related module exists; no repository files were changed.
\item Hereditary normalization is a natural effective use of classical
normal-form uniqueness; the one-variable matrix zero test is classical and close to
weighted-automata equivalence; undecidability of zero testing for unrestricted data
is classical; these are infrastructure, not new discoveries.
\item Fixed-support hardness is a standard-style reduction, not a new halting
theorem; the dyadic validity result is a restricted realization of a classical
obstruction; effective boundedness is an elementary specialization.
\item No decidability of factorization, of elementary theories or of general
Diophantine solvability is inferred.
\item The hereditary normalizer has no polynomial bit-complexity bound.
\item $\Lambda$ exists, but no uniform algorithm produces it; it is presentation
data.
\item The transfer theorem is for ordered coefficient fields; over arbitrary fields
there is only a truth-table reduction; the abstract existence of $L$ does not make
its raw-name equality decidable.
\item Shielding does not embed $L$ as a subfield of $\Oz$.
\item $\SE$ is not closed under $\Om$ on arbitrary inputs, roots or ordinal sums,
and is not all of $\Oz$; fractions need not be reduced, and no multivariate gcd is
hidden.
\item $\FH$ is not real closed.
\item Not every rational multivariate series has a separated-resolvent form
($1/(1-z_1-z_2)$).
\item The front bound is finite, not polynomial, when $r$ is part of the input;
univariate polynomial-time automata bounds are not quoted for the multiscale case;
the variable-rank bit complexity is not settled.
\item With colliding scales the tensor product need not give the Hadamard product
of collected coefficients.
\item The $\Pi^1_1$ theorem classifies validity, not promised equality.
\item The boundedness proposition is not a uniform algorithm, and a bare symbol
for an ordinal is not a computable copy of its well-order.
\item $\varepsilon_0-1$ is also an upper bound: no surreal supremum is claimed.
\item Stream-to-compressed conversion is partial and certificate-producing; a
sound implementation reports ``not established'' rather than a false equality.
\item The finite tests are regression evidence, not proofs; no test computes the
halting set or decides well-foundedness; the package does not implement arbitrary
omnific arithmetic or validate arbitrary well-orders.
\item A rerun rewrites the verification record, whose checksum then changes.
\item A type-level decidability assertion in Lean is not an executable verified
normalizer, and a Python test is not a substitute for a proof of semantic
preservation.
\item Lurie's paper was consulted through its abstract and record, not reanalysed;
L'Innocente--Mantova is context, not an algorithm for notation equality; Arvind et
al.\ are cited for context, not as the same algorithm.
\item The ordinal-notation overview is motivational only.
\item The extensions of its final section are not completed results.
\item Extensions to richer exponent languages must state their extra decision
assumptions.
\end{enumerate}

\subsection{Source 09 (24 items)}
\begin{enumerate}[leftmargin=*,widest=00,label=09.\arabic*]
\item An AI-assisted research draft, not independently refereed or formalized in
Lean; no repository build was run and no repository files were changed.
\item The specific effective theorems and representation comparisons are proposed
contributions after a targeted check, not certified priority claims.
\item The article concerns one language, not a universal decision claim;
``computable'' refers to a stated representation, and the language does not capture
every reasonable computable surreal (Hamkins--Turetsky, Lurie).
\item $\Om$ is Conway's map, not the surreal exponential; no decision result for
unrestricted $\exp$ and $\log$ is claimed.
\item Constants are exact rationals, not arbitrary computable reals; the real trace
is $\Q$, and $\QO$ is not real closed.
\item $\QO$ is not a completion of $\Oz$ and not all countably or effectively
described surreals.
\item The support bound is necessary, not sufficient.
\item No polynomial-time or other complexity bound; expression swell is allowed.
\item Ground equality is decidable, existential Diophantine solvability is not; the
transfer is included to prevent a scope error, not as a new result.
\item The $\CK$ cofinality statement needs uniform presentations; the proposition is
a classical observation and does not apply directly to $\IO$.
\item Abstract truncation closure is Camacho's; geometric omnific examples are
classical; $\sum\omega^{1/n}$ is not a new series.
\item Classical ingredients are not relabelled: normal forms, $\Oz$ and its
constant-term retraction, Hahn geometric inversion, Cantor normal forms, computable
well-orders, the halting and tree completeness inputs, Cayley--Hamilton tests,
clearing denominators.
\item The separated rational basis fails with irrational coefficients.
\item Grid transplantation does not commute with $\Om$.
\item No Euclidean termination is claimed.
\item The recurrence test is not the Skolem problem.
\item The certified intermediate languages use elementary algebra and are not new
general theorems about automata or recurrences.
\item The prototype is not a CAS, a parser or production software, has no halting
oracle, rejects floating-point input, and does not support directly constructed
malformed objects.
\item The finite tests do not prove arbitrary-depth termination, support
well-ordering, ordinal bounds or undecidability, do not verify the $\Pi^1_1$
reduction, and do not enumerate the support of $\mathsf E$.
\item A classical formalized isomorphism, especially one using choice, does not
implement a computable normalizer.
\item The repository review read the catalogue, the computable-surreals README and
the normal-form bridge note, not the full computability article or every Lean
declaration.
\item The $\Pi^1_1$ statement concerns the flexible support language, not
fixed-stream equality or validity of rational--omega terms.
\item The ordinal-notation overview is orientation, not a source of proofs.
\item The research directions are not claims.
\end{enumerate}

\subsection{The merge (6 items)}
\begin{enumerate}[leftmargin=*,widest=00,label=M.\arabic*]
\item The incomparability of the three omnific rings
(\cref{onot:grid:prop:incomparable}) is not established for the pair
$(\mathcal R_\varepsilon,\IO)$ when $k=\Q$; whether $\sum n!\,\omega^{\omega-n}$ has a
rational--omega term or lies in $\FH$ is open (\cref{onot:q:compare}).
\item The characterizations of \cref{onot:rem:characterizations} are the triangular
and grid-ring theorems; the shielding theorem alone gives only the embedding.
\item Only 05's coding gives decidable supports; the midpoint codings are not
claimed to.
\item Results printed by citation (the floor, division, the Diophantine transfer,
the stream reductions, the two-scale sign obstruction, the rational monomial core,
integer evaluation) are proved in the cited reports; this report adds only the
stated effective or omnific refinements.
\item No question named by an earlier report is answered.
\item The reruns for this report used Python 3.14.4 and checked finite instances
only; the dossier's 30 independent checks are likewise finite.
\end{enumerate}

\begin{thebibliography}{99}
\raggedright
\bibitem{Arvind}
V.~Arvind, A.~Chatterjee, R.~Datta and P.~Mukhopadhyay, \emph{Equivalence testing
of weighted automata over partially commutative monoids}, arXiv:2002.08633, v2.
\url{https://arxiv.org/abs/2002.08633}. Cited by 05 for broader context, not as
the same algorithm.

\bibitem{Bancerek}
G.~Bancerek, \emph{Epsilon numbers and Cantor normal form}, Formalized Mathematics
(2009). \href{https://doi.org/10.2478/v10037-009-0032-8}{doi:10.2478/v10037-009-0032-8}.

\bibitem{Camacho}
S.~Camacho, \emph{Truncation in unions of Hahn fields with a derivation},
arXiv:1610.10058, 2016; Propositions 3.7 and 4.1.
\url{https://arxiv.org/abs/1610.10058}.

\bibitem{CMR}
D.~Cenzer, V.~W.~Marek and J.~B.~Remmel, \emph{Index sets for finite normal predicate
logic programs}, arXiv:1303.6555, 2013; Section~2, Theorem 2.6(g).
\url{https://arxiv.org/abs/1303.6555}.

\bibitem{CFH}
S.~Chen, H.~Fang and J.~van der Hoeven, \emph{A zero-test for D-algebraic
transseries}, Proceedings of ISSAC 2026, ACM, 2026, pp.~105--113;
\href{https://arxiv.org/abs/2602.01188}{arXiv:2602.01188}; author-hosted version
dated 6 February 2026, \url{https://www.texmacs.org/joris/zttr/zttr.html}. 08 cites
the arXiv version; the publication data are those recorded in the computer-algebra
report.

\bibitem{Conway}
J.~H.~Conway, \emph{On Numbers and Games}, 2nd ed., A K Peters, 2001 (first edition
Academic Press, 1976).

\bibitem{DowneyMontalban}
R.~Downey and A.~Montalb\'an, \emph{The isomorphism problem for torsion-free abelian
groups is analytic complete}, author manuscript, 7 June 2007; Section~2.
\url{https://math.berkeley.edu/~antonio/papers/torsionFree.pdf}.

\bibitem{Ershov}
Yu.~L.~Ershov, \emph{A hierarchy of sets. I}, Algebra and Logic \textbf{7} (1968),
25--43. \href{https://doi.org/10.1007/BF02218750}{doi:10.1007/BF02218750}.

\bibitem{Fijalkow}
N.~Fijalkow, \emph{A polynomial time algorithm for the equivalence problem of
weighted automata over a field}, expository note, 2019, clarified 2020.
\url{https://games-automata-play.github.io/blog/equivalence_weighted_automata/}.

\bibitem{Gonshor}
H.~Gonshor, \emph{An Introduction to the Theory of Surreal Numbers}, London
Mathematical Society Lecture Note Series 110, Cambridge University Press, 1986.

\bibitem{Hamkins}
J.~D.~Hamkins, \emph{The computable surreal numbers}, Notre Dame Logic Seminar,
3 December 2024, announcement and notes (joint work with D.~Turetsky, discussing
J.~Lurie's earlier work).
\url{https://jdh.hamkins.org/the-computable-surreal-numbers-notre-dame-logic-seminar-december-2024/}.

\bibitem{Kleene}
S.~C.~Kleene, \emph{On notation for ordinal numbers}, Journal of Symbolic Logic
\textbf{3} (1938), 150--155. \href{https://doi.org/10.2307/2267778}{doi:10.2307/2267778}.

\bibitem{LM}
S.~L'Innocente and V.~Mantova, \emph{A factorisation theory for generalised power
series and omnific integers}, Advances in Mathematics (2024), article 109513,
\href{https://doi.org/10.1016/j.aim.2024.109513}{doi:10.1016/j.aim.2024.109513};
arXiv:1710.07304v5. \url{https://arxiv.org/abs/1710.07304}.

\bibitem{Lurie}
J.~Lurie, \emph{The effective content of surreal algebra}, Journal of Symbolic Logic
\textbf{63} (1998), 337--371. \href{https://doi.org/10.2307/2586835}{doi:10.2307/2586835}.

\bibitem{Matiyasevich}
Yu.~V.~Matiyasevich, \emph{Hilbert's Tenth Problem}, MIT Press, 1993; see also
J.~Bayer, M.~David, M.~Hassler, Yu.~Matiyasevich and D.~Schleicher, \emph{Diophantine
equations over $\mathbb Z$: universal bounds and parallel formalization},
arXiv:2506.20909, 2025.

\bibitem{ordinalwiki}
Wikipedia contributors, \emph{Ordinal notation}, accessed 23 September 2026.
\url{https://en.wikipedia.org/wiki/Ordinal_notation}. A user-supplied orientation,
not a source of proofs.

\bibitem{Pohlers}
W.~Pohlers, \emph{Computability Theory of Hyperarithmetical Sets}, lecture notes,
University of M\"unster, 1996; Theorems 5.2.10, 7.1.4, 7.6.6.
\url{https://www.uni-muenster.de/imperia/md/content/logik/Skripte/pohlers._computability_theory_of_hyperarithmetical_sets.pdf}.

\bibitem{Stanley}
R.~P.~Stanley, \emph{Differentiably finite power series}, European Journal of
Combinatorics \textbf{1} (1980), 175--188.
\url{https://math.mit.edu/~rstan/pubs/pubfiles/45.pdf}.
\end{thebibliography}
\end{document}
